% Dissertation : produire le plan détaillé d’une dissertation — Français 1re Tech — Cours signé OnBuch+
% Programme officiel MINESEC. Compiler avec : tectonic N01-dissertation-produire-plan-detaille-dissertation.tex
\newcommand{\DOCMATIERE}{Français}
\newcommand{\DOCNIVEAU}{1re Tech}
\newcommand{\DOCMODULE}{Français – classe de Première technique}
\newcommand{\DOCLECON}{N01}
\newcommand{\DOCTITRE}{Dissertation : produire le plan détaillé d’une dissertation}
% OnBuch+ — préambule « cours signés » ENRICHI (compilé avec tectonic / XeLaTeX)
% Système visuel « Pop » d'OnBuch+ : fond crème (jamais blanc), bordures encre
% épaisses, ombres « dures » décalées, titres Archivo Black en capitales.
% Version MATHÉMATIQUES : graphes (pgfplots/tikz), boîtes pédagogiques
% (définition, propriété, méthode, attention, le savais-tu, exercice résolu,
% exercice, TP, bilan), page de garde et signature OnBuch+ ancrée en bas.
%
% Macros attendues dans main.tex AVANT \input{preamble} :
%   \DOCMATIERE  \DOCNIVEAU  \DOCMODULE  \DOCLECON (numéro)  \DOCTITRE
\documentclass[11pt]{article}

\usepackage{fontspec}
\usepackage[french]{babel}
\usepackage[a4paper,margin=2cm,top=3.4cm,bottom=2.8cm,headheight=1.4cm,footskip=1.3cm]{geometry}
\usepackage{amsmath,amssymb,amsfonts}
\usepackage{mathtools} % \xmapsto, \coloneqq... souvent utilisés par les modèles
\usepackage{cancel} % simplifications barrées (bilans de forces)
% \abs{z} (module, valeur absolue) et \norm{u} (norme), utilisés par les modèles
\providecommand{\abs}{}\let\abs\relax\DeclarePairedDelimiter{\abs}{\lvert}{\rvert}
\providecommand{\norm}{}\let\norm\relax\DeclarePairedDelimiter{\norm}{\lVert}{\rVert}
\usepackage[table]{xcolor} % [table] : fournit aussi \rowcolors (alternance de lignes)
\usepackage{pagecolor}
\usepackage{tikz}
\usetikzlibrary{arrows.meta,positioning,calc,shapes.geometric,shapes.misc,shapes.arrows,decorations.pathmorphing,decorations.pathreplacing,decorations.markings,patterns,shadows,mindmap,trees,fit,backgrounds,matrix,intersections,angles,quotes}
\usepackage{pgfplots}
\usepackage[version=4]{mhchem} % équations nucléaires \ce{^{235}_{92}U}
\usepackage[european,siunitx]{circuitikz} % schémas électriques
\pgfplotsset{compat=1.18}
\usepgfplotslibrary{fillbetween,groupplots} % aires entre courbes (calcul intégral)
\usepackage{enumitem}
\usepackage{fancyhdr}
% [most] charge toutes les bibliothèques tcolorbox (listings, minted, raster,
% documentation...) et gonfle inutilement la mémoire TeX sur un document long
% (~300 boîtes) : on ne charge que ce qui est réellement utilisé.
\usepackage[skins,breakable]{tcolorbox}
\usepackage{graphicx}
\usepackage{colortbl}
\usepackage{booktabs}
\usepackage{tabularx}
\usepackage{multirow}
\usepackage{diagbox} % case d'en-tête diagonale des tableaux à double entrée
% Colonnes extensibles centrée / gauche / droite (C, L, R), souvent utilisées par les modèles
\newcolumntype{C}{>{\centering\arraybackslash}X}
\newcolumntype{L}{>{\raggedright\arraybackslash}X}
\newcolumntype{R}{>{\raggedleft\arraybackslash}X}
\usepackage{array}
\usepackage{multicol}
\usepackage{siunitx}
% parse-numbers=false : accepte aussi \SI{2,5 \times 10^{-3}}{...}
\sisetup{locale=FR,output-decimal-marker={,},group-minimum-digits=4,parse-numbers=false}
\DeclareSIUnit\ohm{\ensuremath{\symup{\Omega}}} % U+2126 absent de la police
\DeclareSIUnit\division{div} % réglages d'oscilloscope (ms/div, V/div)
\DeclareSIUnit\tr{tr} % tours (vitesse de rotation, ex. \SI{1500}{\tr\per\minute})
\DeclareSIUnit\molar{M} % concentration molaire (ex. \SI{3}{\molar}, \SI{1}{\micro\molar})
\DeclareSIUnit\an{an} % année (le modèle confond parfois avec \year, un compteur TeX) — ex. \SI{5}{\kilo\meter\per\an}
\DeclareSIUnit\FCFA{FCFA} % franc CFA (ex. \SI{500}{\FCFA\per\kilo\gram})
\DeclareSIUnit\degreeBrix{\textdegree Brix} % teneur en sucre d'un sirop/jus (réfractométrie)
\DeclareSIUnit\month{mois} % le modèle utilise parfois \month, un compteur TeX, comme unité de durée
\DeclareSIUnit\microliter{\micro\liter} % le modèle écrit parfois \microliter en un seul token
\DeclareSIUnit\million{millions} % dénombrement (ex. \SI{15}{\million\per\mL} spermatozoïdes)
\usepackage{qrcode}
% \degree : commande fréquemment utilisée par les modèles pour un angle ou une
% température en dehors de \SI{}{} (non fournie par les paquets chargés).
\providecommand{\degree}{^{\circ}}
% \qed (fin de démonstration), utilisé par les modèles sans amsthm
\providecommand{\qed}{\hfill$\square$}
% \barre : double barre des valeurs interdites dans un tableau de variations (tkz-tab)
\providecommand{\barre}{\ensuremath{\|}}
% environnement proof (amsthm non chargé) utilisé par les modèles
\ifdefined\proof\else\newenvironment{proof}[1][Démonstration]{\par\noindent\textit{#1.}\ }{\qed\par}\fi
\usepackage{lastpage}
\usepackage{hyperref}
\hypersetup{hidelinks}

% ============================================================
% Palette « Pop » OnBuch+ — mêmes hex que la classe P (plus_ui.dart)
% ============================================================
\definecolor{poporange}{HTML}{F59321}
\definecolor{popdark}{HTML}{E07A0C}
\definecolor{popcream}{HTML}{FFF6E8}
\definecolor{popink}{HTML}{1C1714}
\definecolor{poppurple}{HTML}{7A5AE0}
\definecolor{popgreen}{HTML}{2E9E6A}
\definecolor{poppink}{HTML}{E0457B}
\definecolor{popblue}{HTML}{2F7DE1}
\definecolor{popgold}{HTML}{C98A12}
\definecolor{popmuted}{HTML}{8A7F76}
\definecolor{popline}{HTML}{EADFD2}
\definecolor{popgray}{HTML}{8A7F76}
\colorlet{popgrey}{popgray}
% Alias tolérants (le modèle nomme parfois les couleurs différemment)
\colorlet{popwhite}{white}
\colorlet{popblack}{popink}
\colorlet{popred}{poppink}
\colorlet{popyellow}{popgold}
% Teintes claires (fonds de tableaux/encadrés) — le modèle nomme parfois
% les couleurs avec un suffixe L (ex. popblueL) sans les avoir définies.
\colorlet{poporangeL}{poporange!20}
\colorlet{popdarkL}{popdark!20}
\colorlet{poppurpleL}{poppurple!20}
\colorlet{popgreenL}{popgreen!20}
\colorlet{poppinkL}{poppink!20}
\colorlet{popblueL}{popblue!20}
\colorlet{popgoldL}{popgold!20}
\colorlet{popredL}{popred!20}
\colorlet{popgrayL}{popgray!20}
% Teintes claires pour fonds de tableaux / zones
\colorlet{popblueL}{popblue!10!white}
\colorlet{poporangeL}{poporange!14!white}
\colorlet{popgreenL}{popgreen!12!white}
\colorlet{poppurpleL}{poppurple!12!white}
\colorlet{poppinkL}{poppink!12!white}
\colorlet{popgoldL}{popgold!14!white}

\pagecolor{popcream}

% ============================================================
% Typographie
% ============================================================
\defaultfontfeatures{Ligatures=TeX}
\setmainfont{PlusJakartaSans-Medium.ttf}[Path=../../../fonts/,BoldFont=PlusJakartaSans-Bold.ttf]
% Maths en sans-serif (Fira Math) avec chiffres et lettres droites de Plus
% Jakarta, pour que les formules se fondent dans le texte.
\usepackage[math-style=ISO,bold-style=ISO]{unicode-math}
\setmathfont{FiraMath-Regular.otf}[Path=../../../fonts/]
\setmathfont{PlusJakartaSans-Medium.ttf}[Path=../../../fonts/,range={up/num,up/{Latin,latin},\mathup}]
\usepackage{icomma} % virgule décimale sans espace en mode math
% Caractères mathématiques Unicode absents des polices de texte (Plus Jakarta,
% Archivo Black) : rendus par leur équivalent mathématique, en texte comme en
% formule — sinon ils s'affichent en carré vide (ex. « Arithmétique dans ℤ »).
\usepackage{newunicodechar}
\newunicodechar{ℝ}{\ensuremath{\mathbb{R}}}
\newunicodechar{ℤ}{\ensuremath{\mathbb{Z}}}
\newunicodechar{ℂ}{\ensuremath{\mathbb{C}}}
\newunicodechar{ℕ}{\ensuremath{\mathbb{N}}}
\newunicodechar{ℚ}{\ensuremath{\mathbb{Q}}}
\newunicodechar{𝔻}{\ensuremath{\mathbb{D}}}
\newunicodechar{α}{\ensuremath{\alpha}}
\newunicodechar{β}{\ensuremath{\beta}}
\newunicodechar{γ}{\ensuremath{\gamma}}
\newunicodechar{δ}{\ensuremath{\delta}}
\newunicodechar{ε}{\ensuremath{\varepsilon}}
\newunicodechar{θ}{\ensuremath{\theta}}
\newunicodechar{λ}{\ensuremath{\lambda}}
\newunicodechar{μ}{\ensuremath{\mu}}
\newunicodechar{σ}{\ensuremath{\sigma}}
\newunicodechar{τ}{\ensuremath{\tau}}
\newunicodechar{φ}{\ensuremath{\varphi}}
\newunicodechar{ψ}{\ensuremath{\psi}}
\newunicodechar{ω}{\ensuremath{\omega}}
\newunicodechar{Σ}{\ensuremath{\Sigma}}
% majuscules produites par \MakeUppercase dans les titres (α-aminés -> Α)
\newunicodechar{Α}{\ensuremath{\alpha}}
\newunicodechar{Β}{\ensuremath{\beta}}
\newunicodechar{Γ}{\ensuremath{\Gamma}}
\newunicodechar{Δ}{\ensuremath{\Delta}}
\newunicodechar{Ε}{\ensuremath{\varepsilon}}
\newunicodechar{Θ}{\ensuremath{\Theta}}
\newunicodechar{Λ}{\ensuremath{\Lambda}}
\newunicodechar{Μ}{\ensuremath{\mu}}
\newunicodechar{Τ}{\ensuremath{\tau}}
\newunicodechar{Φ}{\ensuremath{\Phi}}
\newunicodechar{Ψ}{\ensuremath{\Psi}}
\newunicodechar{Ω}{\ensuremath{\Omega}}
\newunicodechar{↦}{\ensuremath{\mapsto}}
\newunicodechar{∈}{\ensuremath{\in}}
\newunicodechar{∉}{\ensuremath{\notin}}
\newunicodechar{∧}{\ensuremath{\wedge}}
\newunicodechar{⇔}{\ensuremath{\Leftrightarrow}}
\newunicodechar{⟺}{\ensuremath{\Longleftrightarrow}}
\newunicodechar{⇒}{\ensuremath{\Rightarrow}}
\newunicodechar{⟹}{\ensuremath{\Longrightarrow}}
\newunicodechar{∘}{\ensuremath{\circ}}
\newunicodechar{∪}{\ensuremath{\cup}}
\newunicodechar{∩}{\ensuremath{\cap}}
\newunicodechar{≡}{\ensuremath{\equiv}}
\newunicodechar{⊂}{\ensuremath{\subset}}
\newunicodechar{⊥}{\ensuremath{\perp}}
\newunicodechar{∃}{\ensuremath{\exists}}
\newunicodechar{∀}{\ensuremath{\forall}}
\newunicodechar{∛}{\ensuremath{\sqrt[3]{\,}}}
\newunicodechar{∜}{\ensuremath{\sqrt[4]{\,}}}
\newunicodechar{⃗}{\ensuremath{{}^{\to}}}
\newunicodechar{ⁿ}{\ifmmode^{n}\else\textsuperscript{\ensuremath{n}}\fi}
\newunicodechar{ˣ}{\ifmmode^{x}\else\textsuperscript{\ensuremath{x}}\fi}
\newunicodechar{ᵇ}{\ifmmode^{b}\else\textsuperscript{\ensuremath{b}}\fi}
\newunicodechar{ʳ}{\ifmmode^{r}\else\textsuperscript{\ensuremath{r}}\fi}
\newunicodechar{ᵘ}{\ifmmode^{u}\else\textsuperscript{\ensuremath{u}}\fi}
\newunicodechar{ʸ}{\ifmmode^{y}\else\textsuperscript{\ensuremath{y}}\fi}
\newunicodechar{ᵃ}{\ifmmode^{a}\else\textsuperscript{\ensuremath{a}}\fi}
\newunicodechar{ᵉ}{\ifmmode^{e}\else\textsuperscript{\ensuremath{e}}\fi}
\newunicodechar{ᵏ}{\ifmmode^{k}\else\textsuperscript{\ensuremath{k}}\fi}
\newunicodechar{ᵖ}{\ifmmode^{p}\else\textsuperscript{\ensuremath{p}}\fi}
\newunicodechar{ᴺ}{\ifmmode^{N}\else\textsuperscript{\ensuremath{N}}\fi}
\newunicodechar{ᶜ}{\ifmmode^{c}\else\textsuperscript{\ensuremath{c}}\fi}
\newunicodechar{ʰ}{\ifmmode^{h}\else\textsuperscript{\ensuremath{h}}\fi}
\newunicodechar{ᵠ}{\ifmmode^{q}\else\textsuperscript{\ensuremath{q}}\fi}
\newunicodechar{ₙ}{\ifmmode_{n}\else\textsubscript{\ensuremath{n}}\fi}
\newunicodechar{ₐ}{\ifmmode_{a}\else\textsubscript{\ensuremath{a}}\fi}
\newunicodechar{ᵢ}{\ifmmode_{i}\else\textsubscript{\ensuremath{i}}\fi}
\newunicodechar{ᵣ}{\ifmmode_{r}\else\textsubscript{\ensuremath{r}}\fi}
\newunicodechar{ₖ}{\ifmmode_{k}\else\textsubscript{\ensuremath{k}}\fi}
\newunicodechar{ᵦ}{\ifmmode_{\beta}\else\textsubscript{\ensuremath{\beta}}\fi}
\newunicodechar{ₓ}{\ifmmode_{x}\else\textsubscript{\ensuremath{x}}\fi}
\newfontfamily\popdisplay{ArchivoBlack-Regular.ttf}[Path=../../../fonts/]
\newfontfamily\poplogo{SpaceGrotesk-Bold.ttf}[Path=../../../fonts/]
\newfontfamily\popsemi{PlusJakartaSans-SemiBold.ttf}[Path=../../../fonts/]
\newfontfamily\popxbold{PlusJakartaSans-ExtraBold.ttf}[Path=../../../fonts/]
\newfontfamily\popxboldls{PlusJakartaSans-ExtraBold.ttf}[Path=../../../fonts/,LetterSpace=3.0]

\setlist[enumerate]{leftmargin=1.7em,itemsep=5pt,topsep=4pt}
\setlist[itemize]{leftmargin=1.7em,itemsep=4pt,topsep=4pt,label={\color{poporange}\textbullet}}
\setlength{\parindent}{0pt}
\setlength{\parskip}{5pt}
\renewcommand{\arraystretch}{1.25}

% pgfplots : style Pop par défaut
\pgfplotsset{
  every axis/.append style={axis line style={popink,line width=1pt},tick style={popink},
    grid style={popline,line width=0.5pt},label style={font=\small},tick label style={font=\footnotesize}},
  popcurve/.style={poporange,line width=1.8pt,no markers,smooth},
}
% Même style disponible dans un \draw TikZ simple (hors pgfplots)
\tikzset{popcurve/.style={poporange,line width=1.8pt}}
\tikzset{popgrid/.style={popline,very thin}}
\colorlet{popcurve}{poporange} % le nom du style est parfois utilisé comme couleur

% ============================================================
% Pill / squiggle
% ============================================================
\newcommand{\poppill}[3]{%
  \tikz[baseline=-0.55ex]\node[draw=popink,line width=1.2pt,rounded corners=2.6pt,%
    fill=#2,inner xsep=6.5pt,inner ysep=3pt]{%
    \popxboldls\fontsize{7}{7}\selectfont\color{#3}\MakeUppercase{#1}};%
}
\newcommand{\popsquiggle}[1]{%
  \tikz[baseline=-0.4ex]{
    \draw[#1,line width=2.6pt,line cap=round]
      plot [smooth] coordinates {(0,0.09) (0.34,0.01) (0.68,0.21) (1.02,0.01) (1.36,0.10)};
  }%
}
% Mot-clé surligné (marqueur orange sous le texte)
\newcommand{\cle}[1]{{\popxbold\color{popdark}#1}}

% ============================================================
% En-tête / pied
% ============================================================
\pagestyle{fancy}
\fancyhf{}
\renewcommand{\headrulewidth}{0pt}
\renewcommand{\footrulewidth}{0pt}
\lhead{\raisebox{-0.15ex}{\poplogo\fontsize{13}{13}\selectfont\color{popink}OnBuch\color{poporange}+}}
\rhead{\poppill{\DOCMATIERE\ \textbullet\ \DOCNIVEAU\ \textbullet\ Leçon \DOCLECON}{white}{popink}}
\lfoot{\popxbold\fontsize{7.6}{7.6}\selectfont\color{popmuted}\MakeUppercase{Cours signé OnBuch+}\ \textbullet\ {\popsemi\DOCTITRE}}
\rfoot{\tikz[baseline=-0.6ex]\node[rounded corners=8.5pt,fill=popink,minimum height=17pt,minimum width=17pt,inner xsep=5pt,inner sep=0]{\popxbold\fontsize{8}{8}\selectfont\color{popcream}\thepage\,/\,\pageref*{LastPage}};}

% ============================================================
% Titres
% ============================================================
\newcommand{\chaptitre}[2]{%
  \begin{tcolorbox}[enhanced,colback=poporange,colframe=popink,boxrule=2.6pt,arc=11pt,
      drop shadow=popink,left=15pt,right=15pt,top=12pt,bottom=11pt,before skip=3pt,after skip=18pt]
    \poppill{#2}{popink}{popcream}\par\vspace{9pt}
    {\raggedright\hyphenpenalty=10000\exhyphenpenalty=10000\popdisplay\fontsize{22}{24}\selectfont\color{popcream}\MakeUppercase{#1}\par}
  \end{tcolorbox}
}
% Section numérotée : pastille orange avec le numéro + titre Archivo
\newcounter{popsec}
\newcommand{\coursec}[1]{%
  \stepcounter{popsec}\setcounter{popsub}{0}%
  \par\vspace{16pt}\noindent
  \begin{minipage}{\linewidth}
  \tikz[baseline=-0.7ex]\node[draw=popink,line width=1.6pt,fill=poporange,rounded corners=4pt,minimum size=22pt,inner sep=0]{\popdisplay\fontsize{12}{12}\selectfont\color{popcream}\thepopsec};%
  \hspace{8pt}{\popdisplay\fontsize{15}{17}\selectfont\color{popink}\MakeUppercase{#1}}\par
  \vspace{4pt}\noindent{\color{poporange}\rule{36pt}{3.6pt}}
  \end{minipage}\par\nopagebreak\vspace{8pt}
}
\newcounter{popsub}
\newcommand{\courssub}[1]{%
  \stepcounter{popsub}%
  \par\vspace{9pt}\noindent{\popxbold\fontsize{11.5}{13}\selectfont\color{popink}{\color{poporange}\thepopsec.\thepopsub}\hspace{6pt}#1}\par\nopagebreak\vspace{4pt}
}

% ============================================================
% Boîtes pédagogiques — même recette : fond blanc, bordure épaisse colorée,
% ombre dure encre, badge plein. Titre optionnel [#1].
% ============================================================
\tcbset{popbase/.style={breakable,enhanced,colback=white,boxrule=2.3pt,arc=9pt,
  shadow={3pt}{-3pt}{0pt}{popink},left=14pt,right=14pt,top=11pt,bottom=11pt,before skip=11pt,after skip=15pt,
  fontupper=\popsemi\fontsize{10.6}{14.5}\selectfont\color{popink},
  fontlower=\popsemi\fontsize{10.6}{14.5}\selectfont\color{popink}}}
% Coche dessinée (le glyphe ✓ n'existe pas dans les polices du document)
\newcommand{\popcheck}[1]{\tikz[baseline=-0.5ex]\draw[#1,line width=1.6pt,line cap=round,line join=round](-0.13,0.0)--(-0.04,-0.09)--(0.13,0.1);}
\providecommand{\coche}{\popcheck{popgreen}}
\providecommand{\resultat}[1]{\par\smallskip\noindent\fcolorbox{popgreen}{popgreenL}{\parbox{\dimexpr\linewidth-2\fboxsep-2\fboxrule\relax}{\textbf{Résultat.} #1}}\par\smallskip}
\newcommand{\popboxhead}[3]{\poppill{#1}{#2}{white}\ifx\relax#3\relax\else\hspace{6pt}{\popxbold\fontsize{10.6}{12}\selectfont\color{popink}#3}\fi\par\vspace{7pt}}

\newtcolorbox{aretenir}[1][]{popbase,colframe=popgreen,before upper={\popboxhead{À retenir}{popgreen}{#1}}}
\newtcolorbox{exemplebox}[1][]{popbase,colframe=poporange,before upper={\popboxhead{Exemple}{poporange}{#1}}}
\newtcolorbox{definition}[1][]{popbase,colframe=popblue,before upper={\popboxhead{Définition}{popblue}{#1}}}
\newtcolorbox{propriete}[1][]{popbase,colframe=poppurple,before upper={\popboxhead{Propriété}{poppurple}{#1}}}
\newtcolorbox{methode}[1][]{popbase,colframe=popgold,before upper={\popboxhead{Méthode}{popgold}{#1}}}
\newtcolorbox{attention}[1][]{popbase,colframe=poppink,before upper={\popboxhead{Attention}{poppink}{#1}}}
\newtcolorbox{experience}[1][]{popbase,colframe=popblue,colback=popblueL,before upper={\popboxhead{Expérience}{popblue}{#1}}}
% « Le savais-tu ? » : carte violette pleine (contexte, culture, Cameroun)
\newtcolorbox{savaistu}[1][]{popbase,colback=poppurple,colframe=popink,
  fontupper=\popsemi\fontsize{10.4}{14.2}\selectfont\color{white},
  before upper={\poppill{Le savais-tu\,?}{white}{poppurple}\ifx\relax#1\relax\else\hspace{6pt}{\popxbold\color{white}#1}\fi\par\vspace{7pt}}}
% Exercice résolu : énoncé en haut, \tcblower puis la solution sur fond crème
\newcounter{popexo}\newcounter{popexr}
\newtcolorbox{exoresolu}[1][]{popbase,colframe=poporange,colbacklower=popcream,
  segmentation style={popink,line width=1.2pt,dashed},
  before upper={\stepcounter{popexr}\popboxhead{Exercice résolu \thepopexr}{poporange}{#1}},
  before lower={\poppill{Solution}{popink}{popcream}\par\vspace{6pt}}}
% Exercice (non résolu en place) — la correction est dans la partie Corrigés
\newtcolorbox{exercice}[1][]{popbase,colframe=popink,
  before upper={\stepcounter{popexo}\popboxhead{Exercice \thepopexo}{popink}{#1}}}
% Corrigé
\newtcolorbox{corrige}[1][]{popbase,colframe=popgreen,colback=popgreenL,
  before upper={\popboxhead{Corrigé}{popgreen}{#1}}}
% Objectifs / prérequis / bilan
\newtcolorbox{objectifs}{popbase,colframe=popink,colback=poporangeL,
  before upper={\popboxhead{Objectifs de la leçon}{popink}{}}}
\newtcolorbox{prerequis}{popbase,colframe=popink,colback=popblueL,
  before upper={\popboxhead{Prérequis}{popblue}{}}}
\newtcolorbox{bilan}{popbase,colframe=popink,colback=white,boxrule=2.8pt,
  before upper={\popboxhead{Fiche bilan}{poporange}{L'essentiel à connaître pour l'examen}}}
% Figure encadrée avec légende
\newtcolorbox{popfigure}[1][]{enhanced,breakable=false,colback=white,colframe=popline,boxrule=1.4pt,arc=8pt,
  left=10pt,right=10pt,top=10pt,bottom=8pt,before skip=10pt,after skip=12pt,halign=center,
  lower separated=false,halign lower=center,
  fontlower=\popsemi\fontsize{9}{11.5}\selectfont\color{popmuted},
  #1}
\newcommand{\legende}[1]{\par\vspace{6pt}{\popsemi\fontsize{9}{11.5}\selectfont\color{popmuted}#1}}

% ============================================================
% Signature OnBuch+
% ============================================================
\newcommand{\poplogoinline}[1]{{\poplogo\fontsize{#1}{#1}\selectfont\color{popink}OnBuch\color{poporange}+}}
\newcommand{\signatureonbuch}{%
  \begin{tcolorbox}[enhanced,colback=popink,colframe=popink,boxrule=2.2pt,arc=10pt,
      drop shadow=poporange,left=14pt,right=14pt,top=10pt,bottom=10pt,before skip=0pt,after skip=0pt]
    \tikz[baseline=-0.7ex]\node[circle,fill=popgreen,draw=popcream,line width=1.3pt,minimum size=22pt,inner sep=0]{\popcheck{white}};%
    \hspace{9pt}\begin{minipage}[c]{0.62\linewidth}
      {\popxbold\fontsize{12}{13}\selectfont\color{popcream}Signé\ \textbullet\ {\poplogo OnBuch\color{poporange}+}}\par\vspace{2pt}
      {\raggedright\popsemi\fontsize{8}{10.5}\selectfont\color{popline}\MakeUppercase{Équipe pédagogique OnBuch+}\\\MakeUppercase{Conforme au programme officiel MINESEC}\par}
    \end{minipage}\hfill
    {\poplogo\fontsize{22}{22}\selectfont\color{popcream}OnBuch\color{poporange}+}
  \end{tcolorbox}%
}

% ============================================================
% Page de garde
% #1 titre · #2 matière · #3 niveau · #4 module · #5 n° leçon · #6 durée ·
% #7 description / objectifs (texte ou itemize)
% ============================================================
\newcommand{\pagegarde}[7]{%
  \thispagestyle{empty}
  \newgeometry{a4paper,margin=1.7cm,top=1.6cm,bottom=1.4cm}%
  \begin{tikzpicture}[remember picture,overlay]
    \fill[poporange!18] ([xshift=-3.2cm,yshift=-3.6cm]current page.north east) circle (4.6cm);
    \fill[poppurple!14] ([xshift=2.4cm,yshift=7.6cm]current page.south west) circle (3.8cm);
  \end{tikzpicture}%
  {\poplogo\fontsize{30}{30}\selectfont\color{popink}OnBuch\color{poporange}+}\hfill
  \raisebox{0.6ex}{\poppill{Cours signé \textbullet\ Premium}{popink}{popcream}}\par
  \vspace{1pt}\popsquiggle{poppurple}\par
  \vspace{8pt}
  \begin{tcolorbox}[enhanced,colback=poporange,colframe=popink,boxrule=2.8pt,arc=13pt,
      drop shadow=popink,left=18pt,right=18pt,top=12pt,bottom=13pt,after skip=10pt]
    \poppill{#2 \textbullet\ #3}{popink}{popcream}\hspace{5pt}\poppill{#4}{white}{popink}\par\vspace{8pt}
    {\popdisplay\fontsize{14}{14}\selectfont\color{popink}LEÇON #5}\par\vspace{4pt}
    {\raggedright\hyphenpenalty=10000\exhyphenpenalty=10000\popdisplay\fontsize{25}{27}\selectfont\color{popcream}\MakeUppercase{#1}\par}
    \vspace{8pt}
    {\popxbold\fontsize{9.5}{9.5}\selectfont\color{popink}\MakeUppercase{Durée conseillée : #6}}
  \end{tcolorbox}
  \begin{tcolorbox}[enhanced,colback=white,colframe=popink,boxrule=2.2pt,arc=10pt,
      drop shadow=popink,left=16pt,right=16pt,top=10pt,bottom=10pt,after skip=8pt]
    \poppill{Dans cette leçon}{poppurple}{white}\par\vspace{5pt}
    \popsemi\fontsize{9.3}{12.6}\selectfont\color{popink}#7
  \end{tcolorbox}
  \noindent\begin{minipage}[t]{0.487\linewidth}
    \begin{tcolorbox}[enhanced,colback=popgreen,colframe=popink,boxrule=2.2pt,arc=10pt,
        drop shadow=popink,left=11pt,right=11pt,top=10pt,bottom=10pt,height=3.1cm,valign=center]
      \tcbox[colback=white,colframe=popink,boxrule=1.3pt,arc=5pt,boxsep=0pt,nobeforeafter,
        left=3pt,right=3pt,top=3pt,bottom=3pt,valign=center]{\qrcode[height=1.9cm]{https://play.google.com/store/apps/details?id=com.luvvix.onbuch&hl=fr}}%
      \hspace{8pt}\begin{minipage}[c]{0.5\linewidth}
        {\popdisplay\fontsize{11}{12.5}\selectfont\color{white}\MakeUppercase{Télécharge OnBuch}}\par\vspace{4pt}
        {\raggedright\popsemi\fontsize{8.4}{11}\selectfont\color{white}Gratuit \textbullet\ Google Play\\Tuteur IA, annales, concours, résultats\par}
      \end{minipage}
    \end{tcolorbox}
  \end{minipage}\hfill
  \begin{minipage}[t]{0.487\linewidth}
    \begin{tcolorbox}[enhanced,colback=poppurple,colframe=popink,boxrule=2.2pt,arc=10pt,
        drop shadow=popink,left=11pt,right=11pt,top=10pt,bottom=10pt,height=3.1cm,valign=center]
      \tikz[baseline=-0.7ex]\node[circle,fill=white,draw=popink,line width=1.3pt,minimum size=1.6cm,inner sep=0]{\popdisplay\fontsize{24}{24}\selectfont\color{poppurple}+};%
      \hspace{8pt}\begin{minipage}[c]{0.6\linewidth}
        {\popdisplay\fontsize{11}{12.5}\selectfont\color{white}\MakeUppercase{Passe à OnBuch+}}\par\vspace{4pt}
        {\raggedright\popsemi\fontsize{8.4}{11}\selectfont\color{white}Tous les cours signés, Tuteur IA illimité, fiches PDF — onglet OnBuch+ de l'app\par}
      \end{minipage}
    \end{tcolorbox}
  \end{minipage}
  \vspace{8pt}
  \popcontenu
  \vfill
  \signatureonbuch
  \restoregeometry
  \newpage
}

% Rangée « Ce cours contient » (page de garde)
\newcommand{\poptile}[3]{% #1 couleur #2 gros texte #3 légende
  \begin{tcolorbox}[enhanced,colback=#1,colframe=popink,boxrule=2pt,arc=9pt,shadow={2.5pt}{-2.5pt}{0pt}{popink},
      left=6pt,right=6pt,top=9pt,bottom=9pt,halign=center,valign=center,height=1.75cm,width=\linewidth,nobeforeafter]
    {\popdisplay\fontsize{12.5}{13}\selectfont\color{white}\MakeUppercase{#2}}\par\vspace{3pt}
    {\centering\popsemi\fontsize{7.8}{9.5}\selectfont\color{white}#3\par}
  \end{tcolorbox}}
\newcommand{\popcontenu}{%
  \par\noindent{\popxboldls\fontsize{7.5}{7.5}\selectfont\color{popmuted}CE COURS SIGNÉ CONTIENT}\par\vspace{7pt}
  \noindent\begin{minipage}[t]{0.235\linewidth}\poptile{popblue}{Cours}{complet, illustré, pas à pas}\end{minipage}\hfill
  \begin{minipage}[t]{0.235\linewidth}\poptile{poporange}{Méthodes}{et exercices résolus}\end{minipage}\hfill
  \begin{minipage}[t]{0.235\linewidth}\poptile{poppink}{Exercices}{type examen + corrigés}\end{minipage}\hfill
  \begin{minipage}[t]{0.235\linewidth}\poptile{popgreen}{Fiche bilan}{l'essentiel à retenir}\end{minipage}\par}

% Sommaire
\newtcolorbox{sommaire}{enhanced,colback=white,colframe=popink,boxrule=2.2pt,arc=10pt,
  drop shadow=popink,left=18pt,right=18pt,top=13pt,bottom=13pt,before skip=6pt,after skip=20pt,
  before upper={\poppill{Sommaire}{poppurple}{white}\par\vspace{9pt}\popsemi\fontsize{10.6}{14.5}\selectfont\color{popink}}}

% Signature de fin de cours — poussée tout en bas de la dernière page
\newcommand{\findecours}{%
  \par\vfill\nopagebreak
  \begin{center}\popsquiggle{poporange}\end{center}\vspace{4pt}
  \signatureonbuch
}
\AtBeginDocument{\def\llbracket{\mathopen{[\mkern-3.5mu[}}\def\rrbracket{\mathclose{]\mkern-3.5mu]}}} % intervalles d'entiers (glyphe ⟦ absent de la police)
% Glyphes absents de Fira Math : redéfinis à partir de glyphes disponibles
\AtBeginDocument{%
  \def\setminus{\mathbin{\backslash}}\let\smallsetminus\setminus
  \def\vdots{\mathord{\vbox{\baselineskip3.2pt\lineskiplimit0pt\kern4pt\hbox{.}\hbox{.}\hbox{.}}}}%
  \def\ddots{\mathinner{\mkern1mu\raise6.5pt\hbox{.}\mkern2mu\raise3.6pt\hbox{.}\mkern2mu\raise0.7pt\hbox{.}\mkern1mu}}%
  \def\triangle{\mathord{\scalebox{0.9}{$\symup{\Delta}$}}}%
  \def\nabla{\mathord{\raisebox{\depth}{\scalebox{1}[-1]{$\symup{\Delta}$}}}}%
  \def\checkmark{\popcheck{popdark}}%
  \def\star{\mathbin{\ast}}\def\bigstar{\mathord{\ast}}%
  \def\diamond{\mathbin{\scalebox{0.6}{$\Diamond$}}}%
  \def\bigcup{\mathop{\mathchoice{\raisebox{-0.3em}{\scalebox{1.7}{$\cup$}}}{\scalebox{1.25}{$\cup$}}{\cup}{\cup}}\displaylimits}%
  \def\bigcap{\mathop{\mathchoice{\raisebox{-0.3em}{\scalebox{1.7}{$\cap$}}}{\scalebox{1.25}{$\cap$}}{\cap}{\cap}}\displaylimits}%
  \def\lesseqqgtr{\mathrel{\lessgtr}}\def\bigoplus{\mathop{\oplus}\displaylimits}%
}
\providecommand{\ssi}{si et seulement si} % abréviation fréquente
\AtBeginDocument{\providecommand{\wideparen}{\overparen}} % arc de cercle

%% FIGFIX-BEGIN v4 — correctifs globaux des figures (docs/onbuch-generation/tools/figfix)
\usepackage{adjustbox}\usepackage{etoolbox}
% toute figure TikZ/pgfplots plus large que la ligne est réduite à la largeur disponible
\BeforeBeginEnvironment{tikzpicture}{\begin{adjustbox}{max width=\linewidth}}
\AfterEndEnvironment{tikzpicture}{\end{adjustbox}}
% légendes pgfplots lisibles par-dessus les courbes
\pgfplotsset{every axis/.append style={legend style={fill=white,fill opacity=0.92,text opacity=1,draw=black!15,inner sep=3pt}}}
% étiquettes d'arêtes (syntaxe edge["..."]) avec fond blanc
\tikzset{every edge quotes/.append style={fill=white,inner sep=1.2pt}}
%% FIGFIX-END
\begin{document}

\pagegarde{Dissertation : produire le plan détaillé d’une dissertation}{Français}{1re Tech}{Français – classe de Première technique}{N01}{21 séances (fiche de progression harmonisée)}{Cette leçon outille l'élève de Première technique pour analyser un sujet de dissertation, formuler une problématique, rechercher et organiser les idées, puis élaborer un plan détaillé conforme aux exigences du Baccalauréat. Elle s'appuie sur des situations de communication concrètes et sur l'étude du texte ouvroir Le Lion et la perle de Wole Soyinka.
\vspace{4pt}
\begin{multicols}{2}\small
\begin{itemize}
\item À la fin de cette leçon, je sais identifier les paratextes d'une œuvre et en tirer des hypothèses de lecture
\item À la fin de cette leçon, je sais analyser un sujet de dissertation (mots-clés, type de sujet, présupposés)
\item À la fin de cette leçon, je sais formuler une problématique sous forme de question directrice
\item À la fin de cette leçon, je sais mobiliser les techniques de recherche des idées (brainstorming, questionnement, mobilisation de la culture)
\item À la fin de cette leçon, je sais distinguer et construire les trois types de plans (thématique, dialectique, analytique)
\item À la fin de cette leçon, je sais rédiger un plan détaillé complet avec arguments et exemples hiérarchisés
\end{itemize}\end{multicols}}

\begin{sommaire}
\begin{multicols}{2}
\begin{enumerate}
\item Entrée dans l'unité : paratextes et situations de communication
\item Analyse du sujet et formulation de la problématique
\item La recherche des idées
\item L'élaboration du plan détaillé
\item Outils linguistiques de la dissertation
\item Compte rendu, remédiation et projet d'études
\item Activité d'intégration
\item Méthodes et exercices résolus
\item Exercices
\item Corrigés des exercices
\item Fiche bilan
\end{enumerate}
\end{multicols}
\end{sommaire}

\begin{objectifs}
\begin{itemize}
\item À la fin de cette leçon, je sais identifier les paratextes d'une œuvre et en tirer des hypothèses de lecture
\item À la fin de cette leçon, je sais analyser un sujet de dissertation (mots-clés, type de sujet, présupposés)
\item À la fin de cette leçon, je sais formuler une problématique sous forme de question directrice
\item À la fin de cette leçon, je sais mobiliser les techniques de recherche des idées (brainstorming, questionnement, mobilisation de la culture)
\item À la fin de cette leçon, je sais distinguer et construire les trois types de plans (thématique, dialectique, analytique)
\item À la fin de cette leçon, je sais rédiger un plan détaillé complet avec arguments et exemples hiérarchisés
\item À la fin de cette leçon, je sais identifier les situations de communication et les marques explicites de l'énonciation dans un texte
\item À la fin de cette leçon, je sais employer correctement la ponctuation et distinguer dénotation et connotation pour affiner mon expression
\item À la fin de cette leçon, je sais confronter, évaluer et améliorer une production écrite à l'aide d'une grille critériée
\end{itemize}
\end{objectifs}

\begin{prerequis}
\begin{itemize}
\item Notions de Seconde sur les genres littéraires et le texte argumentatif
\item Lecture méthodique : repérage du thème et de la thèse d'un texte (classe de Seconde)
\item Grammaire de la phrase : types et formes de phrases, connecteurs logiques
\item Connaissance générale de l'œuvre au programme et de son contexte africain
\item Techniques de prise de notes et de classement d'informations
\end{itemize}
\end{prerequis}

\newpage

\coursec{Entrée dans l'unité : paratextes et situations de communication}

Au concours d'entrée dans une grande école de Yaoundé, deux candidats traitent le même sujet de dissertation : « Le travail garantit-il la liberté de l'homme ? ». Le premier écrit dix pages d'idées brillantes mais désordonnées ; le second, trois pages claires, structurées, avec un plan visible. C'est le second qui est admis. Pourquoi ? Parce qu'en dissertation, la forme du raisonnement compte autant que les idées. Comment passer d'un sujet à un plan détaillé solide, capable de convaincre un correcteur du \textbf{Probatoire} ? Cette unité vous donne les clés : lire les paratextes pour entrer dans l'œuvre, maîtriser la situation de communication pour analyser tout texte, puis construire pas à pas une argumentation rigoureuse.

\courssub{Les paratextes : définition, repérage et exploitation}

\begin{definition}[Paratexte (Gérard Genette)]
Le \cle{paratexte} est l'ensemble des éléments qui entourent le texte principal et orientent sa lecture : titre, nom d'auteur, couverture, quatrième de couverture, préface, dédicace, illustrations, mais aussi bandeaux publicitaires, entretiens d'auteur, critiques. Genette distingue le \cle{péritexte} (physiquement attaché au livre) de l'\cle{épitexte} (extérieur au livre : interviews, critiques, correspondances).
\end{definition}

\begin{methode}[Formuler des hypothèses de lecture à partir du titre et de la quatrième de couverture]
\begin{enumerate}
\item \textbf{Lire le titre} : repérer les mots-clés, les oppositions, les allusions culturelles.
\item \textbf{Lire la quatrième de couverture} : identifier le résumé, les thèmes annoncés, le ton (promesse d'aventure, de réflexion, de témoignage).
\item \textbf{Croiser les deux} : formuler 2 à 3 hypothèses sur le genre, le thème central, le contexte historique, le point de vue narratif.
\item \textbf{Noter ses attentes} : qu'attendez-vous de cette lecture ? Quelles questions le texte devra-t-il répondre ?
\end{enumerate}
\end{methode}

\begin{popfigure}
\begin{center}
\begin{tabularx}{\linewidth}{|l|X|X|}\hline
\rowcolor{popblueL}
\textbf{Paratexte} & \textbf{Péritexte (dans le livre)} & \textbf{Épitexte (hors du livre)} \\ \hline
\textbf{Titre} & \textit{Le Lion et la perle} : opposition animale / minérale, force / beauté, tradition / modernité & Entretiens de Soyinka : « Le lion est la tradition, la perle la modernité » \\ \hline
\textbf{Auteur} & Wole Soyinka (Nigeria, prix Nobel 1986) & Biographies, discours de réception du Nobel \\ \hline
\textbf{Couverture} & Souvent : masque africain, couleurs ocre/rouge, graphismes géométriques & Affiches de spectacles, éditions successives \\ \hline
\textbf{Quatrième de couv.} & Résumé : conflit entre Baroka (chef traditionnel) et Lakunle (instituteur moderne) pour Sidi & Critiques universitaires, dossiers pédagogiques \\ \hline
\textbf{Préface} & Parfois préface d'un éditeur ou d'un critique (ex. : « Une comédie satirique ») & Articles de presse, études universitaires \\ \hline
\textbf{Dédicace} & « À ma mère » ou absente selon éditions & Correspondances de l'auteur \\ \hline
\textbf{Illustrations} & Gravures, masques, schémas de scène (rares en édition de poche) & Photographies de mises en scène, costumes \\ \hline
\end{tabularx}
\end{center}
\legende{Tableau comparatif : péritexte et épitexte de \textit{Le Lion et la perle} de Wole Soyinka}
\end{popfigure}

\begin{exemplebox}[Hypothèses de lecture pour \textit{Le Lion et la perle}]
À partir du titre et de la quatrième de couverture (éd. Présence Africaine) :
\begin{itemize}
\item \textbf{Genre} : théâtre (pièce en trois actes, dialogues, didascalies).
\item \textbf{Thème central} : affrontement tradition / modernité au Nigeria colonial (années 1950-60).
\item \textbf{Contexte} : village yoruba (Ilujinle), pouvoir coutumier vs éducation occidentale.
\item \textbf{Personnages types} : Baroka (le « Lion », chef rusé, polygamie), Lakunle (le « moderne », maladroit, monogame), Sidi (la « Perle », beauté villageoise, enjeu du conflit).
\item \textbf{Ton} : comédie satirique, ironie, proverbes yoruba, chants et danses intégrés.
\end{itemize}
\end{exemplebox}

\begin{savaistu}[Wole Soyinka]
\cle{Wole Soyinka} (né en 1934 à Abeokuta, Nigeria) est le premier Africain à recevoir le \textbf{prix Nobel de littérature} (1986). Dramaturge, poète, essayiste, militant pour la démocratie, il a été emprisonné pendant la guerre du Biafra. Son œuvre mêle tradition orale yoruba et formes occidentales. \textit{Le Lion et la perle} (1959) est sa pièce la plus jouée dans le monde francophone.
\end{savaistu}

\courssub{Le texte ouvroir : rôle et première approche de \textit{Le Lion et la perle}}

\begin{aretenir}[Texte ouvroir]
Le \cle{texte ouvroir} est l'œuvre intégrale étudiée tout au long de l'unité. Il sert de \textbf{fil conducteur} : tous les exercices (analyse, dissertation, résumé, commentaire, expression orale) s'y réfèrent. Il permet d'ancrer les notions abstraites (paratexte, énonciation, figures de style, argumentation) dans un matériau concret et unique.
\end{aretenir}

\begin{experience}[Lecture des paratextes de \textit{Le Lion et la perle} — Travail guidé]
\textbf{Matériel} : édition de poche (Présence Africaine ou Hatier) ou photocopies couverture + quatrième de couverture + page de titre + préface s'il y a.
\textbf{Protocole} :
\begin{enumerate}
\item En binôme, remplir le tableau ci-dessous (à reproduire dans le cahier).
\item Échanger avec un autre binôme, confronter les hypothèses.
\item Mise en commun au tableau : construire une « fiche d'identité » collective de l'œuvre.
\end{enumerate}
\textbf{Tableau à compléter} :
\begin{center}
\begin{tabularx}{\linewidth}{|l|X|X|}\hline
\rowcolor{popblueL}
\textbf{Paratexte} & \textbf{Observations (ce que je vois/lis)} & \textbf{Hypothèses (ce que j'en déduis)} \\ \hline
Titre & & \\ \hline
Auteur / nationalité / date & & \\ \hline
Couverture (images, couleurs) & & \\ \hline
Quatrième de couverture (résumé, extraits) & & \\ \hline
Préface (si présente) & & \\ \hline
Dédicace & & \\ \hline
Illustrations intérieures & & \\ \hline
\end{tabularx}
\end{center}
\textbf{Observations attendues} : le titre métaphorique, l'auteur nigérian prix Nobel, la couverture aux motifs traditionnels, le résumé annonçant une « comédie » sur le choc des cultures, la préface situant l'œuvre dans le théâtre africain d'expression anglaise.
\textbf{Interprétation} : ces paratextes préparent le lecteur à une pièce engagée, drôle, ancrée dans la culture yoruba, questionnant l'avenir de l'Afrique face à la modernité importée.
\end{experience}

\courssub{La situation de communication : le modèle de Jakobson}

\begin{definition}[Situation de communication (Roman Jakobson)]
Tout acte de communication met en jeu six \cle{composantes} indispensables :
\begin{itemize}
\item l'\cle{émetteur} (qui parle/écrit),
\item le \cle{récepteur} (à qui s'adresse le message),
\item le \cle{message} (ce qui est transmis),
\item le \cle{code} (la langue, le système de signes utilisé),
\item le \cle{canal} (le support physique : air, papier, onde radio, écran),
\item le \cle{référent} (la réalité dont on parle, le sujet du message).
\end{itemize}
\end{definition}

\begin{popfigure}
\begin{tikzpicture}[
  func/.style={ellipse, draw=popblue, thick, fill=popblue!10, minimum width=2.4cm, minimum height=1.3cm, align=center, font=\small},
  comp/.style={rectangle, draw=poporange, thick, fill=poporange!10, minimum width=2.2cm, minimum height=1cm, align=center, font=\small},
  arrow/.style={-Stealth, thick, popdark},
  node distance=1.5cm and 1.5cm
]
% Facteurs (hexagone)
\node[comp] (emetteur) {ÉMETTEUR};
\node[comp] (recepteur) [right=of emetteur] {RÉCEPTEUR};
\node[comp] (message) [below right=of recepteur] {MESSAGE};
\node[comp] (code) [below left=of message] {CODE};
\node[comp] (canal) [left=of code] {CANAL};
\node[comp] (referent) [above left=of canal] {RÉFÉRENT};

% Fonctions (ellipses)
\node[func, align=center] (expressive) [above left=of emetteur]{EXPRESSIVE\\ (émetteur)};
\node[func, align=center] (conative) [above right=of recepteur]{CONATIVE\\ (récepteur)};
\node[func, align=center] (referentielle) [right=of message]{RÉFÉRENTIELLE\\ (référent)};
\node[func, align=center] (poetique) [below right=of message]{POÉTIQUE\\ (message)};
\node[func, align=center] (metalang) [below left=of code]{MÉTALINGUISTIQUE\\ (code)};
\node[func, align=center] (phatique) [left=of canal]{PHATIQUE\\ (canal)};

% Liens facteurs -> fonctions
\draw[arrow] (emetteur) -- (expressive);
\draw[arrow] (recepteur) -- (conative);
\draw[arrow] (referent) -- (referentielle);
\draw[arrow] (message) -- (poetique);
\draw[arrow] (code) -- (metalang);
\draw[arrow] (canal) -- (phatique);

% Hexagone facteurs
\draw[thick, popline] (emetteur) -- (recepteur) -- (message) -- (code) -- (canal) -- (referent) -- cycle;

% Légende
\node[align=center, font=\footnotesize, text=popmuted] at (0, -4.2) {Modèle de Jakobson : 6 facteurs $\rightarrow$ 6 fonctions du langage};
\end{tikzpicture}
\legende{Schéma de la communication selon Jakobson : les six facteurs et les six fonctions du langage}
\end{popfigure}

\begin{attention}[Piège fréquent : référent vs code]
Ne confondez pas le \textbf{référent} (ce dont on parle : un objet, un événement, une idée) et le \textbf{code} (la langue ou le système de signes : français, anglais, morse, langage informatique). Exemple : « Le lion rugit. » — Référent : l'animal lion / son rugissement. Code : la langue française (vocabulaire, grammaire).
\end{attention}

\begin{exoresolu}[Identification des composantes dans un extrait de \textit{Le Lion et la perle}]
\textbf{Extrait (Acte I, Lakunle à Sidi) :} « Sidi, ma perle, mon trésor, écoute-moi ! La coutume du prix de la dot est une barbarie. Je te paierai, mais pas comme on achète une bête. Nous serons égaux, comme au pays du Blanc. »
\tcblower
\textbf{Analyse} :
\begin{itemize}
\item \textbf{Émetteur} : Lakunle (instituteur, partisan de la modernité).
\item \textbf{Récepteur} : Sidi (jeune fille villageoise, belle, convoitée).
\item \textbf{Message} : refus de la dot coutumière, proposition d'un mariage « moderne » égalitaire.
\item \textbf{Code} : anglais (langue originale) / français (traduction) ; registre soutenu, vocabulaire militant (« barbarie », « égaux »).
\item \textbf{Canal} : oral (théâtre, voix de l'acteur), ondes sonores dans la salle.
\item \textbf{Référent} : la coutume de la dot (pratique sociale yoruba), le modèle occidental de mariage.
\end{itemize}
\end{exoresolu}

\courssub{Les six fonctions du langage}

\begin{propriete}[Les six fonctions du langage (Jakobson)]
Chaque composante de la communication détermine une fonction dominante du message :
\begin{itemize}
\item \textbf{Fonction expressive (émotive)} : centrée sur l'\textbf{émetteur} — exprime sentiments, émotions (ex. : « Quelle beauté ! »).
\item \textbf{Fonction conative} : centrée sur le \textbf{récepteur} — cherche à influencer son comportement (ex. : « Fermez la porte. »).
\item \textbf{Fonction référentielle} : centrée sur le \textbf{référent} — informe sur la réalité (ex. : « Il pleut à Yaoundé. »).
\item \textbf{Fonction phatique} : centrée sur le \textbf{canal} — vérifie, ouvre ou maintient le contact (ex. : « Allô ? Vous m'entendez ? »).
\item \textbf{Fonction métalinguistique} : centrée sur le \textbf{code} — parle de la langue elle-même (ex. : « "Dot" signifie "prix de la mariée". »).
\item \textbf{Fonction poétique} : centrée sur le \textbf{message} — soigne la forme pour elle-même (ex. : proverbe, rime, parallélisme : « Le lion ne perd pas son temps à chasser la souris. »).
\end{itemize}
\end{propriete}

\begin{exemplebox}[Fonctions dominantes dans l'extrait de Lakunle]
\begin{center}
\begin{tabularx}{\linewidth}{|l|X|l|}\hline
\rowcolor{popblueL}
\textbf{Segment} & \textbf{Fonction dominante} & \textbf{Justification} \\ \hline
« Sidi, ma perle, mon trésor » & Expressive & Émotion, séduction, subjectivité \\ \hline
« écoute-moi ! » & Conative & Impératif, appelle l'attention de Sidi \\ \hline
« La coutume de la dot est une barbarie » & Référentielle & Jugement sur une réalité sociale \\ \hline
« "Dot" signifie "prix de la mariée" » & Métalinguistique & Explication lexicale (si présent) \\ \hline
« Nous serons égaux, comme au pays du Blanc » & Poétique & Parallélisme, rythme, image valorisante \\ \hline
\end{tabularx}
\end{center}
\end{exemplebox}

\courssub{Application : situations de communication du quotidien camerounais}

\begin{methode}[Analyser une situation de communication]
\begin{enumerate}
\item Identifier les six composantes (émetteur, récepteur, message, code, canal, référent).
\item Déterminer la ou les fonctions dominantes.
\item Noter le registre de langue (soutenu, courant, familier), le ton, les marques d'énonciation (je, tu, indicateurs spatio-temporels).
\item Conclure sur l'intention communicative et l'efficacité du message.
\end{enumerate}
\end{methode}

\begin{exoresolu}[Analyse de quatre situations camerounaises]
\textbf{Situation 1 : Palabre villageoise (chef de village $\rightarrow$ notables, place publique, langue locale + français)} \\
Composantes : Émetteur = chef (autorité traditionnelle) ; Récepteur = notables + villageois ; Message = décision foncière, règlement de conflit ; Code = langue locale (bassa, ewondo, bamileke…) + français officiel ; Canal = oral, amplifié par porte-voix ; Référent = terre, coutume, lignage. \\
Fonctions : \textbf{Référentielle} (exposer les faits), \textbf{Conative} (ordonner, apaiser), \textbf{Phatique} (saluts rituels, « Le village vous écoute »).

\textbf{Situation 2 : Annonce à la radio CRTV (speaker $\rightarrow$ auditeurs, ondes FM, français/anglais)} \\
Composantes : Émetteur = speaker (institutionnel) ; Récepteur = public national ; Message = communiqué officiel (ex. : résultats examen, alerte météo) ; Code = français standard / anglais ; Canal = ondes radio ; Référent = événement national. \\
Fonctions : \textbf{Référentielle} dominante (information), \textbf{Phatique} (« Yaoundé, CRTV, 20 heures »).

\textbf{Situation 3 : Message WhatsApp (élève $\rightarrow$ camarade, smartphone, français familier + emojis)} \\
Composantes : Émetteur = élève ; Récepteur = ami(e) ; Message = « Wesh, t'as vu le prof ? Il a pas venu ! 😂 » ; Code = français familier, abréviations, emojis ; Canal = internet (data) ; Référent = absence du prof. \\
Fonctions : \textbf{Expressive} (émotion, humour), \textbf{Conative} (appeler une réaction), \textbf{Poétique} (jeu graphique emojis).

\textbf{Situation 4 : Discours officiel (Ministre $\rightarrow$ nation, télévision, français soutenu)} \\
Composantes : Émetteur = Ministre (fonction) ; Récepteur = citoyens ; Message = politique publique, réforme ; Code = français administratif, rhétorique ; Canal = TV + streaming ; Référent = éducation, santé, économie. \\
Fonctions : \textbf{Référentielle} (chiffres, mesures), \textbf{Conative} (mobiliser, rassurer), \textbf{Métalinguistique} (définir termes techniques).
\tcblower
\textbf{Bilan pédagogique} : La maîtrise du modèle de Jakobson permet de décoder l'intention de l'émetteur et d'adapter sa propre production (dissertation, courriel, oral) au récepteur visé (correcteur, supérieur, pair), en choisissant le code, le canal et les fonctions pertinents.
\end{exoresolu}

\begin{exercice}[Entraînement : identifier composantes et fonctions]
Pour chacune des situations ci-dessous, complétez le tableau (émetteur, récepteur, message, code, canal, référent) et indiquez la fonction dominante.
\begin{enumerate}
\item Un agent de la SNEC (Société Nationale d'Électricité) laisse un avis de coupure sur la porte d'un abonné.
\item Un griot chante les louanges d'un notable lors d'une cérémonie de funérailles à Bafoussam.
\item Un élève envoie un courriel formel à son professeur pour demander un rendez-vous.
\item Une pancarte « Interdit de stationner » plantée sur un trottoir de Douala.
\end{enumerate}
\end{exercice}

\begin{corrige}[Exercice n°1]
\begin{enumerate}
\item \textbf{Avis SNEC} : Émetteur = agent SNEC ; Récepteur = abonné ; Message = date/heure coupure, excuse ; Code = français administratif ; Canal = papier (avis collé) ; Référent = coupure électrique programmée. \textit{Fonction dominante : Référentielle (informer) + Phatique (contact institutionnel).}
\item \textbf{Griot à Bafoussam} : Émetteur = griot ; Récepteur = assistance + notable honoré ; Message = généalogie, vertus, proverbes ; Code = langue bamileke, registre poétique ; Canal = oral (voix, tam-tam) ; Référent = vie du défunt, histoire du lignage. \textit{Fonction dominante : Poétique (forme chantée) + Expressive (émotion) + Conative (louer pour honorer).}
\item \textbf{Courriel élève $\rightarrow$ prof} : Émetteur = élève ; Récepteur = professeur ; Message = demande RDV, créneaux proposés ; Code = français soutenu, formules de politesse ; Canal = internet (email) ; Référent = entretien pédagogique. \textit{Fonction dominante : Conative (obtenir RDV) + Phatique (formules « Monsieur, je vous prie d'agréer... »).}
\item \textbf{Pancarte « Interdit de stationner »} : Émetteur = autorité municipale ; Récepteur = automobilistes ; Message = interdiction, sanction possible ; Code = français signalétique, pictogramme ; Canal = panneau métallique (visuel) ; Référent = stationnement illicite. \textit{Fonction dominante : Conative (interdire) + Référentielle (règle).}
\end{enumerate}
\end{corrige}

\coursec{Analyse du sujet et formulation de la problématique}

Après avoir exploré les paratextes du \emph{Lion et la perle} et identifié les situations de communication qui structurent nos échanges quotidiens, nous entrons dans le cœur de la production écrite : la dissertation. Comprendre un sujet, c'est déjà le résoudre à moitié. Cette section vous donne la méthode rigoureuse pour ne jamais être \og hors-sujet \fg{}, la première cause d'échec au Probatoire.

\courssub{Nature d'un sujet de dissertation}

Un sujet de dissertation n'est jamais une simple invitation à raconter ou à réciter un cours. C'est une \cle{question}, une \cle{citation} ou une \cle{affirmation à discuter} qui contient une tension intellectuelle. Votre tâche n'est pas de dire ce que vous savez sur le thème, mais de répondre à l'interrogation précise que le sujet pose.

\begin{definition}[Les trois formes du sujet]
\begin{itemize}
    \item \textbf{La question directe} : \og L'école est-elle la seule voie de la réussite ? \fg{} Elle attend une réponse nuancée (oui, mais / non, car).
    \item \textbf{L'affirmation à discuter} : \og Le progrès technique rend l'homme libre. \fg{} Elle suppose un accord partiel, une discussion des limites.
    \item \textbf{La citation} : \og La tradition, c'est la garde du feu, non la vénération des cendres \fg{} (Gustav Mahler). Elle demande d'analyser la portée de la pensée de l'auteur et de la confronter à d'autres points de vue.
\end{itemize}
\end{definition}

\courssub{Les cinq étapes de l'analyse du sujet}

L'analyse suit un cheminement logique, de la surface du texte (les mots) vers sa profondeur (la problématique). Voici le schéma de cette montée en abstraction :

\begin{popfigure}
\begin{tikzpicture}[
    node distance=1.2cm and 2cm,
    stepbox/.style={rectangle, draw=popblue, thick, fill=popblueL, rounded corners, minimum width=3.5cm, minimum height=1.2cm, align=center, font=\sffamily\bfseries},
    arrow/.style={-Stealth, thick, popdark}
]
\node[stepbox, align=center] (e1){\textbf{Étape 1}\\\emph{Lecture attentive}\\\& Soulignement des mots-clés};
\node[stepbox, below of=e1, align=center] (e2){\textbf{Étape 2}\\\emph{Définition des mots-clés}\\\& (Dénotation / Connotation)};
\node[stepbox, below of=e2, align=center] (e3){\textbf{Étape 3}\\\emph{Dégagement du thème}\\\& des présupposés};
\node[stepbox, below of=e3, align=center] (e4){\textbf{Étape 4}\\\emph{Repérage de la tension}\\\& ou du paradoxe};
\node[stepbox, below of=e4, align=center] (e5){\textbf{Étape 5}\\\emph{Formulation de la}\\\emph{problématique}};
\draw[arrow] (e1) -- (e2);
\draw[arrow] (e2) -- (e3);
\draw[arrow] (e3) -- (e4);
\draw[arrow] (e4) -- (e5);
\end{tikzpicture}
\legende{Les 5 marches de l'analyse du sujet : du mot-clé à la question directrice.}
\end{popfigure}

\textbf{Étape 1 : Lecture attentive et soulignement des mots-clés}\par
Prenez le temps. Lisez trois fois. Soulignez les \cle{termes du sujet} : les noms centraux, les verbes d'action (discuter, analyser, faut-il), les connecteurs logiques (seule, seulement, aussi), les mots qui portent la charge polémique.
\begin{methode}[Repérer les mots-clés]
\begin{enumerate}
    \item Isoler les \textbf{termes substantifs} (ex: \og école \fg{}, \og réussite \fg{}, \og voie \fg{}).
    \item Repérer les \textbf{termes restrictifs ou exclusifs} (ex: \og seule \fg{}, \og unique \fg{}, \og principal \fg{}).
    \item Identifier le \textbf{verbe opératoire} (ex: \og est-elle \fg{} $\rightarrow$ jugement de valeur ; \og faut-il \fg{} $\rightarrow$ obligation morale/pratique).
\end{enumerate}
\end{methode}

\textbf{Étape 2 : Définition des mots-clés (dénotation) et repérage de leurs connotations}\par
C'est l'étape linguistique cruciale. Un mot a un sens littéral (dénotation) et des sens suggérés, culturels, affectifs (connotations). Ne vous contentez pas du dictionnaire : le contexte du sujet oriente le sens.

\begin{methode}[Définir un mot-clé : Dictionnaire + Contexte + Champ lexical]
\begin{enumerate}
    \item \textbf{Dénotation} : Définition littérale (ex: \og École \fg{} = Établissement d'enseignement).
    \item \textbf{Connotations} : Valeurs associées (ex: \og École \fg{} $\rightarrow$ discipline, diplôme, conformisme, mais aussi émancipation, socialisation).
    \item \textbf{Champ lexical} : Mots de la même famille, synonymes, antonymes (ex: \og Réussite \fg{} $\rightarrow$ réussite sociale, professionnelle, personnelle, épanouissement, échec, argent, honneur).
\end{enumerate}
\end{methode}

\begin{popfigure}
\begin{tikzpicture}[
    mindmap,
    concept color=popblue,
    level 1/.style={level distance=3.5cm, sibling angle=90},
    level 2/.style={level distance=2.5cm, sibling angle=45},
    every node/.style={concept, execute at begin node=\sffamily\small},
    root concept/.append style={font=\sffamily\bfseries\large, concept color=poporange},
    child concept/.append style={font=\sffamily\small}
]
\node[root concept] {Tradition}
    child[concept color=popgreen] {node[align=center]{Coutume\\Rituel}
        child {node {Héritage}}
        child {node {Transmission}}
    }
    child[concept color=poppurple] {node[align=center]{Ancêtres\\Mémoire}
        child {node {Identité}}
        child {node {Racines}}
    }
    child[concept color=poppink] {node[align=center]{Stabilité\\Ordre social}
        child {node {Cohésion}}
        child {node {Normes}}
    }
    child[concept color=popgold] {node[align=center]{Modernité\\Rupture}
        child {node {Progrès}}
        child {node {Innovation}}
        child {node {Conflit générations}}
    };
\end{tikzpicture}
\legende{Explosion du mot-clé \og tradition \fg{} en champs lexicaux (lien avec \emph{Le Lion et la perle} : Baroka vs Lakunle).}
\end{popfigure}

\textbf{Étape 3 : Dégagement du thème et des présupposés}\par
Le \cle{thème} est le sujet général (l'éducation, la tradition, le progrès). Les \cle{présupposés} sont ce que le sujet \emph{suppose vrai} sans le dire explicitement.
\begin{itemize}
    \item Sujet : \og L'école est-elle la seule voie de la réussite ? \fg{}
    \item Thème : Éducation et réussite sociale.
    \item Présupposés : 1) L'école \emph{est} une voie de réussite. 2) Il existe \emph{d'autres} voies possibles. 3) La \og réussite \fg{} est un concept univoque et mesurable.
\end{itemize}
Mettre au jour les présupposés, c'est déjà trouver vos axes de discussion.

\textbf{Étape 4 : Repérage de la tension ou du paradoxe}\par
Une dissertation sans tension est un exposé. La tension naît de l'opposition entre deux pôles légitimes.
\begin{itemize}
    \item \textbf{Pôle A (Thèse)} : L'école apporte diplômes, compétences, ascenseur social $\rightarrow$ voie royale.
    \item \textbf{Pôle B (Antithèse)} : Entrepreneuriat, talent artistique, héritage, autodidaxie $\rightarrow$ réussites hors école.
    \item \textbf{Paradoxe} : L'école forme pour le salariat, mais l'économie crée peu d'emplois salariés ; l'école standardise, mais la réussite demande de la singularité.
\end{itemize}
C'est cette friction qui produit la problématique.

\textbf{Étape 5 : Formulation de la problématique}\par
La problématique n'est \textbf{pas} le sujet reformulé. C'est une \cle{question directrice} (ou un faisceau de questions imbriquées) qui \textbf{naît de la tension} et \textbf{commande l'ensemble du devoir}. Elle doit être :
\begin{itemize}
    \item \textbf{Ouverte} : ne pas appeler un oui/non.
    \item \textbf{Complexe} : articuler les termes du sujet (souvent par \og dans quelle mesure \fg{}, \og comment \fg{}, \og jusqu'à quel point \fg{}).
    \item \textbf{Opératoire} : elle dicte votre plan (I, II, III).
\end{itemize}

\begin{definition}[Problématique]
Question directrice qui naît de la tension contenue dans le sujet et qui commande l'ensemble du devoir. Elle transforme un thème en objet d'étude problématisé.
\end{definition}

\begin{attention}[Piège mortel : La paraphrase]
\og \textbf{Erreur :} \og Il faut voir si l'école est la seule voie de la réussite. \fg{} (C'est le sujet avec \og il faut voir \fg{} devant).
\og \textbf{Erreur :} \og Quelles sont les voies de la réussite ? \fg{} (Trop large : catalogue, pas de tension).
\og \textbf{Bonne problématique :} \og \textbf{Dans quelle mesure l'institution scolaire détient-elle le monopole de la légitimation sociale de la réussite, face aux parcours alternatifs qui la contestent ?} \fg{} (Articule école / réussite / monopole / alternatives).
\end{attention}

\courssub{Distinction entre problématiques : trop large, trop étroite, pertinente}

\begin{center}
\begin{tabularx}{\linewidth}{|l|X|X|X|}\hline
\textbf{Type} & \textbf{Caractéristique} & \textbf{Exemple (Sujet : Faut-il préserver les traditions ?)} & \textbf{Conséquence au Bac} \\ \hline
Trop large & Dépasse le sujet, ignore la tension & \og Qu'est-ce que la tradition dans le monde ? \fg{} & Hors-sujet (note $\le$ 5/20) \\ \hline
Trop étroite & Se focalise sur un détail, empêche le plan & \og Faut-il garder le port du pagne au Cameroun ? \fg{} & Sujet traité partiellement (note $\approx$ 8-10/20) \\ \hline
\rowcolor{popgreenL}
\textbf{Pertinente} & \textbf{Articule les termes, ouvre le débat, guide le plan} & \og \textbf{Comment arbitrer entre la nécessaire transmission du patrimoine culturel et l'indispensable adaptation aux mutations du monde moderne ?} \fg{} & \textbf{Sujet maîtrisé, plan cohérent (note $\ge$ 14/20)} \\ \hline
\end{tabularx}
\end{center}

\courssub{Entraînement sur sujets types Bac technique}

\begin{exoresolu}[Sujet : « Le progrès technique rend-il l'homme plus libre ? »]
\textbf{Étape 1 -- Mots-clés :} \og Progrès technique \fg{}, \og homme \fg{}, \og libre \fg{}, \og plus \fg{} (comparatif), \og rend \fg{} (causalité).
\textbf{Étape 2 -- Définitions :}
\begin{itemize}
    \item \og Progrès technique \fg{} : Dénotation = amélioration des outils/méthodes. Connotation = émancipation (médecine, transport) MAIS aliénation (écrans, surveillance, dépendance énergétique).
    \item \og Libre \fg{} : Dénotation = qui n'est pas esclave. Connotation = autonomie, choix, responsabilité, mais aussi isolement, anxiété du choix.
\end{itemize}
\textbf{Étape 3 -- Thème \& Présupposés :} Thème = Technique et condition humaine. Présupposé : le progrès technique \emph{change} la condition de l'homme ; la liberté est un bien mesurable.
\textbf{Étape 4 -- Tension :} Libération des contraintes naturelles (faim, distance, maladie) VS Nouvelles servitudes (addiction, surveillance algorithmique, rythme imposé par la machine).
\textbf{Étape 5 -- Problématique :}
\begin{center}
\og \textbf{Dans quelle mesure le progrès technique, en affranchissant l'homme de la nécessité naturelle, ne le soumet-il pas à de nouvelles formes de dépendance technique, économique et psychique qui redéfinissent les contours de sa liberté ?} \fg{}
\end{center}
\textbf{Plan induit :}
I. L'affranchissement par la technique (maîtrise de la nature, temps libre, santé).
II. Les nouvelles chaînes de la société technicienne (aliénation au travail, surveillance, consumérisme).
III. Vers une liberté reconquise : éducation critique, éthique du numérique, low-tech.
\tcblower
\textbf{Solution détaillée} : La problématique utilise \og Dans quelle mesure \fg{} + structure binaire (affranchissement / soumission) + ouverture (redéfinition). Elle intègre \og plus \fg{} (comparatif historique) et \og rend \fg{} (causalité complexe). Le plan en trois parties suit exactement la dialectique de la question.
\end{exoresolu}

\begin{exoresolu}[Sujet : « Faut-il préserver les traditions ? » (Lien \emph{Le Lion et la perle})]
\textbf{Analyse rapide :} \og Faut-il \fg{} (obligation morale), \og préserver \fg{} (conserver intact vs adapter), \og traditions \fg{} (coutumes, rites, valeurs).
\textbf{Tension :} Baroka (le Lion) incarne la tradition vivante, adaptative, séductrice ; Lakunle incarne le moderne rejeton, caricatural, qui refuse l'héritage. Sidi hésite.
\textbf{Problématique pertinente :}
\og \textbf{La préservation des traditions exige-t-elle leur figement muséal ou leur réinvention créative pour répondre aux défis de la modernité ?} \fg{}
\textbf{Plan :} I. Préserver = transmettre l'identité (rôle du \og feu \fg{}). II. Figer = tuer la tradition (le piège des \og cendres \fg{}, Lakunle). III. La tradition vivante : Baroka, synthèse dynamique (le timbre-poste, la danse).
\tcblower
\textbf{Solution détaillée} : Cette problématique cite Mahler (paratexte culturel), utilise les personnages de la pièce comme arguments littéraires, et structure le devoir sur le mouvement \og conservation / adaptation \fg{}.
\end{exoresolu}

\begin{savaistu}[Le hors-sujet : ennemi n°1 au Probatoire]
Au Baccalauréat technique, \textbf{plus de 60\% des notes inférieures à 10/20 en dissertation} sanctionnent un hors-sujet. Dans 9 cas sur 10, l'élève a \textbf{sauté l'étape 2} (définition/connotation) ou l'\textbf{étape 4} (tension). Il a écrit \og tout ce qu'il savait sur l'école \fg{} au lieu de répondre \og l'école est-elle la \textbf{seule} voie ? \fg{}. \textbf{Astuce :} Avant de rédiger l'introduction, relisez votre problématique et demandez-vous : \og Si je réponds à cette question, est-ce que je traite \textbf{exactement} le sujet posé ? \fg{} Si non, reformulez.
\end{savaistu}

\begin{aretenir}[Check-list avant de passer au plan]
\begin{enumerate}
    \item Ai-je \textbf{souligné et défini} chaque mot-clé (dénotation + connotation) ?
    \item Ai-je formulé les \textbf{présupposés} cachés du sujet ?
    \item Ai-je écrit \textbf{noir sur blanc la tension} (Thèse vs Antithèse) ?
    \item Ma problématique est-elle une \textbf{question ouverte, complexe, opératoire} ?
    \item Mon plan (I, II, III) répond-il \textbf{point par point} à cette problématique ?
\end{enumerate}
Si une réponse est \og non \fg{}, revenez en arrière. Le temps de l'analyse (20-25 min sur 4h) est du temps gagné sur la rédaction.
\end{aretenir}

\coursec{La recherche des idées}

Dans la section précédente, nous avons appris à analyser un sujet de dissertation : repérer les mots-clés, définir les termes, délimiter le champ et formuler une problématique sous forme de question. Mais une problématique, aussi bien formulée soit-elle, reste une question sans réponse. Pour y répondre, il faut de la \emph{matière} : des idées, des arguments, des exemples. C'est précisément l'objet de cette nouvelle étape, trop souvent négligée par les candidats pressés : la \cle{recherche des idées}. C'est le moment où l'on remplit son « sac » avant de ranger ses affaires : le rangement (le plan détaillé) viendra dans la section suivante.

\courssub{Qu'est-ce que chercher des idées ?}

\begin{definition}[Recherche des idées]
La \cle{recherche des idées} est l'étape de la dissertation qui consiste, après l'analyse du sujet, à rassembler au brouillon le maximum d'\cle{arguments} (raisons qui répondent à la problématique) et d'\cle{exemples} (faits précis qui illustrent ces raisons), avant de les trier et de les classer.
\end{definition}

L'intuition est simple : devant un sujet comme « Faut-il préserver les traditions ? », le cerveau ne produit pas d'emblée un plan en trois parties. Il produit d'abord des étincelles désordonnées : « les traditions unissent les familles », « mais certaines sont dépassées », « Sidi dans \emph{Le Lion et la perle} », « le mariage coutumier », « la fête Ngondo »... Le travail de l'élève consiste à \textbf{capter toutes ces étincelles}, puis à les organiser. On distingue donc deux gestes successifs : \textbf{produire} (sans censure), puis \textbf{sélectionner et classer} (avec rigueur).

\begin{definition}[Argument et exemple]
Un \cle{argument} est une raison avancée pour défendre ou réfuter une thèse. Il répond à la question « pourquoi ? ». Un \cle{exemple} est un fait précis (œuvre, événement historique, fait de société, proverbe, expérience vécue) qui illustre et prouve l'argument. L'argument est le squelette ; l'exemple est la chair qui le rend visible et convaincant.
\end{definition}

\courssub{Les techniques de production des idées}

\textbf{La technique du brainstorming (remue-méninges).} Le principe est de libérer totalement la pensée pendant un temps court : on note \emph{tout} ce qui vient à l'esprit à propos du sujet, sous forme de mots ou de courtes phrases, sans juger, sans raturer, sans chercher à organiser. La censure est l'ennemie de cette étape : une idée qui semble faible peut en appeler une meilleure. Le tri viendra après.

\begin{methode}[Conduire un brainstorming efficace en 10 minutes]
\begin{enumerate}
\item Recopier le sujet en haut du brouillon et entourer les mots-clés (2 minutes).
\item Régler mentalement un chronomètre : écrire pendant 5 minutes \textbf{toutes} les idées, questions, exemples, souvenirs de lecture qui se rapportent au sujet, en vrac, sans rature ni hiérarchie. Règle d'or : \textbf{la quantité d'abord}.
\item Relire la liste et souligner d'une couleur les idées directement liées à la problématique (2 minutes).
\item Regrouper par petits paquets les idées voisines et barrer les hors-sujets (1 minute). Le tri vient \textbf{ensuite}, jamais pendant la production.
\end{enumerate}
\end{methode}

\textbf{La technique du questionnement.} Pour éviter la page blanche, on interroge systématiquement le sujet à l'aide des questions fondamentales : \textbf{qui ? quoi ? où ? quand ? comment ? pourquoi ?} Appliquées au sujet « Faut-il préserver les traditions ? », ces questions deviennent un véritable moteur d'idées :

\begin{center}
\begin{tabularx}{\linewidth}{|l|X|}
\hline
\rowcolor{popblueL} \textbf{Question} & \textbf{Idées dégagées} \\
\hline
Qui ? & Les dépositaires des traditions : anciens, chefs traditionnels, familles ; les jeunes, souvent attirés par la modernité. \\
\hline
Quoi ? & Ce qu'on appelle « traditions » : langues, rites, danses, mariage coutumier, funérailles, contes, proverbes. \\
\hline
Où ? & Villages et villes du Cameroun, d'Afrique ; confrontation village/ville (Sidi et Lakunle à Ilujinle). \\
\hline
Quand ? & Traditions anciennes face à la modernité actuelle : école, photographie, progrès technique. \\
\hline
Comment ? & Transmission orale, cérémonies, initiation ; menaces : urbanisation, mépris, oubli des langues. \\
\hline
Pourquoi ? & Pourquoi préserver : identité, cohésion, sagesse. Pourquoi transformer : certaines pratiques freinent ou blessent. \\
\hline
\end{tabularx}
\end{center}

Chaque case du tableau est une piste d'argument potentielle. Le questionnement transforme un sujet intimidant en une série de petites questions faciles.

\courssub{Mobiliser ses connaissances : les réservoirs d'idées}

D'où viennent les idées ? De quatre \cle{réservoirs} que tout élève de Première technique possède déjà :

\begin{itemize}
\item \textbf{Les œuvres étudiées en classe.} C'est le réservoir le plus sûr, car il est frais et précis. Pour notre sujet, \emph{Le Lion et la perle} de Wole Soyinka (écrite en 1959, publiée en 1963) est une mine : Sidi, la « perle » du village, hésite entre le mariage traditionnel et la vanité nouvelle ; Baroka, le « Lion », chef de 62 ans, incarne une tradition rusée qui sait feindre la modernité ; Lakunle, l'instituteur, défend un progrès maladroit et caricatural (refus de la dot, chaussures, « civilisation »). La pièce montre que ni la tradition figée ni la modernité copiée ne suffisent.
\item \textbf{Les lectures personnelles.} Romans, contes, articles de journaux lus hors de la classe : ils prouvent une culture personnelle.
\item \textbf{L'actualité et l'histoire camerounaises et africaines.} Le mariage coutumier qui précède ou accompagne le mariage civil, la fête \cle{Ngondo} des Sawa à Douala, les funérailles traditionnelles, la place des langues nationales dans l'éducation, la chefferie traditionnelle reconnue par l'État : autant d'exemples vivants et proches.
\item \textbf{Les proverbes et la sagesse orale.} « Ce que le vieillard voit assis, le jeune ne le voit pas debout » : un proverbe bien choisi peut introduire ou conclure un argument avec force.
\end{itemize}

\begin{exemplebox}[Une mobilisation concrète et chiffrée]
Pour le sujet « Faut-il préserver les traditions ? », un bon brouillon peut mobiliser en dix minutes : \textbf{trois personnages} d'une œuvre au programme (Sidi, Baroka, Lakunle dans \emph{Le Lion et la perle}, 1959), \textbf{deux réalités camerounaises} (le mariage coutumier avec la remise de la dot, la fête Ngondo célébrée chaque année par les Sawa au bord du Wouri), \textbf{un fait historique} (la colonisation a imposé des langues et des modèles étrangers), \textbf{un proverbe} africain. Soit \textbf{sept exemples précis}, largement suffisants pour alimenter trois parties de développement.
\end{exemplebox}

La carte mentale ci-dessous visualise cette recherche d'idées autour du sujet.

\begin{popfigure}
\begin{center}
\begin{tikzpicture}[mindmap, grow cyclic,
  every node/.style={concept, font=\small},
  concept color=popblue!60,
  level 1/.append style={level distance=3.6cm, sibling angle=90, font=\small},
  level 2/.append style={level distance=2.4cm, sibling angle=50, font=\scriptsize}]
\node[concept color=popdark!80, text=white] {Faut-il préserver les traditions ?}
  child[concept color=popgreen!60] { node {Oui : identité}
    child { node {Cohésion familiale et sociale} }
    child { node {Sagesse des proverbes, langues} }
    child { node {Ngondo, mariage coutumier} }
  }
  child[concept color=poporange!70] { node {Oui : équilibre}
    child { node {Baroka : la tradition rusée s'adapte} }
    child { node {Modernité sans racines = dérive} }
  }
  child[concept color=poppink!70] { node {Non : limites}
    child { node {Certaines pratiques dépassées} }
    child { node {Lakunle : le refus de la dot} }
    child { node {Progrès : école, santé, droits} }
  }
  child[concept color=popgold!70] { node {Exemples}
    child { node {\emph{Le Lion et la perle} (1959)} }
    child { node {Histoire : colonisation} }
    child { node {Proverbes africains} }
  };
\end{tikzpicture}
\end{center}
\legende{Carte mentale de la recherche d'idées pour le sujet « Faut-il préserver les traditions ? » : au centre le sujet, autour les grandes pistes d'arguments (en couleurs) et, au bout de chaque branche, les exemples précis qui les illustrent.}
\end{popfigure}

\courssub{Sélectionner : éliminer et regrouper}

Une fois la matière produite, il faut trier. Deux gestes s'imposent.

\textbf{Éliminer les hors-sujets.} Une idée est hors sujet si elle ne répond pas à la problématique. Par exemple, raconter toute l'intrigue du \emph{Lion et la perle} ou décrire la ville de Douala ne répond pas à la question « faut-il préserver les traditions ? ». Critère simple : si l'on ne peut pas rattacher l'idée à la problématique par un « parce que » ou un « donc », on la barre.

\textbf{Regrouper les idées voisines.} « Les traditions unissent les familles » et « les funérailles rassemblent le village » sont deux formulations d'une même idée : la cohésion sociale. On les fusionne en un seul argument, le second élément devenant l'exemple du premier. Ce regroupement prépare directement le plan : chaque paquet d'idées correspondra à une partie ou à un paragraphe.

\begin{center}
\begin{tabularx}{\linewidth}{|X|X|X|}
\hline
\rowcolor{popblueL} \textbf{Idée du brainstorming} & \textbf{Sort} & \textbf{Justification} \\
\hline
Les traditions fondent l'identité d'un peuple & Retenue & Répond directement au « pourquoi préserver ». \\
\hline
\emph{Le Lion et la perle} : Baroka défend les coutumes & Retenue & Exemple littéraire précis, au programme. \\
\hline
La fête Ngondo rassemble les Sawa chaque année & Retenue & Exemple camerounais concret de cohésion. \\
\hline
Certaines coutumes sont contraires aux droits humains & Retenue & Alimente la thèse adverse (nuance). \\
\hline
Sidi est la plus belle fille d'Ilujinle & Écartée & Anecdote : ne répond pas à la problématique. \\
\hline
Douala est une grande ville portuaire & Écartée & Hors sujet : aucun lien avec les traditions. \\
\hline
J'aime la danse traditionnelle de ma région & Écartée (ou transformée) & Trop personnel et vague ; utilisable seulement s'il est rattaché à un argument. \\
\hline
\end{tabularx}
\end{center}

\courssub{Classer : du constat à l'approfondissement}

Les idées retenues doivent enfin être \textbf{ordonnées}. Deux principes simples guident ce classement :

\begin{itemize}
\item \textbf{Du plus simple au plus complexe} : on commence par ce que le lecteur comprend immédiatement (les traditions donnent une identité) avant d'aborder ce qui demande plus de réflexion (comment concilier tradition et progrès).
\item \textbf{Du constat à l'approfondissement} : on décrit d'abord la réalité (les traditions structurent encore la vie sociale au Cameroun), puis on l'explique (pourquoi sont-elles précieuses ?), puis on la discute (quelles limites ? quel avenir ?).
\end{itemize}

Ce classement épouse naturellement la logique du plan dialectique (thèse, antithèse, synthèse) ou du plan thématique, que nous construirons dans la section suivante.

\begin{attention}[Erreur fréquente : l'accumulation d'exemples sans argument]
Beaucoup de copies alignent les exemples comme des perles sans fil : « Par exemple Ngondo, par exemple le mariage coutumier, par exemple Baroka... ». Or \textbf{un exemple ne vaut que s'il est expliqué} et rattaché à un argument qui le porte. La bonne séquence est toujours : \textbf{argument} (les traditions renforcent la cohésion sociale) $\rightarrow$ \textbf{exemple} (la fête Ngondo réunit chaque année les Sawa autour de leurs rites) $\rightarrow$ \textbf{explication} (cette célébration transmet aux jeunes les valeurs et l'histoire du groupe). Un exemple sans argument est une preuve sans procès.
\end{attention}

\begin{savaistu}[Ce que les correcteurs valorisent vraiment]
Les correcteurs du Probatoire et du Baccalauréat technique valorisent les exemples \textbf{précis} : un auteur nommé, un titre exact, un fait daté, un lieu identifiable. « Dans \emph{Le Lion et la perle} (1959), Wole Soyinka oppose Baroka à Lakunle » impressionne infiniment plus que « on constate que dans certains livres, on voit que... ». Les formules vagues (« on constate que », « de tout temps », « dans la vie ») signalent une culture approximative ; les références exactes signalent un candidat sérieux. Chaque exemple précis peut faire gagner des points.
\end{savaistu}

\courssub{Se constituer une banque d'arguments et d'exemples}

La recherche des idées ne commence pas le jour de l'examen : elle se prépare toute l'année. L'outil idéal est le \cle{carnet de culture} : un cahier (ou un fichier) où l'élève consigne régulièrement, classées par grands thèmes (tradition et modernité, amour, pouvoir, éducation, femme et société, guerre et paix), les références rencontrées : œuvres étudiées, faits historiques, actualités, proverbes. Pour chaque entrée, on note : la référence précise, l'idée qu'elle illustre, et les sujets où elle pourrait servir. Ainsi, la fiche « \emph{Le Lion et la perle} » pourra servir pour des sujets sur la tradition, sur l'éducation, sur la condition féminine ou sur le pouvoir.

\begin{exoresolu}[Du brainstorming au classement]
\textbf{Énoncé.} Voici le brouillon d'un élève pour le sujet « Faut-il préserver les traditions ? » : (1) les traditions donnent une identité ; (2) Lakunle veut moderniser Ilujinle ; (3) mon quartier est animé le soir ; (4) certaines coutumes violent les droits des femmes ; (5) le Ngondo rassemble les Sawa ; (6) il faut concilier tradition et progrès ; (7) Baroka triomphe par la ruse ; (8) les proverbes transmettent la sagesse. Éliminez les hors-sujets, regroupez les idées voisines et classez-les du constat à l'approfondissement.
\tcblower
\textbf{Solution détaillée.}
\begin{enumerate}
\item \textbf{Élimination.} L'idée (3) est un hors-sujet : l'animation d'un quartier ne répond pas à la question des traditions. On la barre.
\item \textbf{Regroupement.} Les idées (1), (5) et (8) forment un paquet : la valeur des traditions (identité, cohésion, sagesse). Les idées (2) et (4) forment un paquet : la critique de certaines traditions et l'aspiration au progrès. L'idée (6) forme à elle seule la voie de dépassement. L'idée (7) est un exemple à ranger dans le premier paquet (la tradition est vivante et forte), à condition de l'expliquer.
\item \textbf{Classement.} Du constat à l'approfondissement : d'abord le constat de la valeur des traditions (1 + 5 + 8 + 7) ; ensuite la critique et les limites (2 + 4) ; enfin l'approfondissement, c'est-à-dire la synthèse (6) : préserver l'essentiel tout en s'ouvrant au progrès, comme le suggère la victoire de Baroka, qui utilise la photographie — un objet moderne — au service de la tradition.
\end{enumerate}
Ce classement dessine déjà, en germe, un plan dialectique en trois parties.
\end{exoresolu}

\begin{exercice}[À vous de jouer]
Sujet : « L'école détruit-elle les valeurs traditionnelles ? »
\begin{enumerate}
\item Conduisez un brainstorming de cinq minutes et notez au moins huit idées ou exemples.
\item Appliquez la technique du questionnement (qui ? quoi ? où ? quand ? comment ? pourquoi ?) pour enrichir votre liste.
\item Éliminez les hors-sujets, regroupez les idées voisines et classez-les du constat à l'approfondissement.
\item Pour chaque argument retenu, associez un exemple précis (œuvre, fait camerounais ou africain, proverbe) et rédigez la phrase qui l'explique.
\end{enumerate}
\end{exercice}

\begin{corrige}[Exercice : L'école détruit-elle les valeurs traditionnelles ?]
Éléments de correction attendus.
\begin{enumerate}
\item Brainstorming possible : l'école occupe le temps des enfants ; les langues françaises et anglaises dominent à l'école ; Lakunle, l'instituteur, méprise les coutumes d'Ilujinle ; mais l'école transmet aussi l'histoire ; les programmes incluent les langues nationales ; l'école ouvre l'esprit et lutte contre les pratiques néfastes ; les grands-parents transmettent les contes le soir ; certains élèves ne parlent plus leur langue maternelle.
\item Questionnement : qui ? les élèves, les maîtres, les familles ; quoi ? langues, rites, respect des anciens ; où ? écoles urbaines et rurales du Cameroun ; quand ? depuis la colonisation ; comment ? par l'horaire, la langue d'enseignement, les programmes ; pourquoi ? l'école forme pour la modernité, ce qui peut éloigner des coutumes.
\item Classement : constat (l'école éloigne parfois des langues et des rites) ; limites du reproche (l'école lutte contre les pratiques néfastes et peut transmettre la culture) ; approfondissement (école et tradition peuvent se compléter : enseignement des langues nationales).
\item Exemple de phrase d'explication : « L'école peut détourner les jeunes des coutumes (argument) : dans \emph{Le Lion et la perle}, Lakunle, l'instituteur, méprise la dot et les rites d'Ilujinle (exemple), ce qui montre qu'une éducation mal assimilée peut couper l'homme de ses racines (explication). »
\end{enumerate}
\end{corrige}

\begin{aretenir}[L'essentiel de la recherche des idées]
\begin{itemize}
\item La recherche des idées suit l'analyse du sujet et précède le plan détaillé : on produit d'abord, on trie ensuite.
\item Deux techniques complémentaires : le \cle{brainstorming} (quantité d'abord, sans censure) et le \cle{questionnement} (qui ? quoi ? où ? quand ? comment ? pourquoi ?).
\item Quatre réservoirs d'idées : œuvres étudiées, lectures personnelles, actualité et histoire camerounaises et africaines, proverbes.
\item Trois gestes de mise en ordre : éliminer les hors-sujets, regrouper les idées voisines, classer du constat à l'approfondissement.
\item Règle d'or : tout exemple doit être précis et rattaché à un argument par une explication.
\item Le carnet de culture prépare cette étape toute l'année.
\end{itemize}
\end{aretenir}

Nos idées sont désormais produites, triées et classées. Il reste à leur donner une architecture solide et définitive : c'est l'objet de la section suivante, consacrée à l'élaboration du plan détaillé.

\coursec{L'élaboration du plan détaillé}

Après avoir exploré les techniques de \textbf{recherche des idées} (brainstorming, mots-clés, questions QQOQCP, lecture d'ouvroir), nous disposons désormais d'une matière brute : arguments, exemples, citations, définitions. Cette matière doit maintenant être \textbf{architecturée}. Le plan détaillé est le squelette de la dissertation ; sans lui, les idées les plus brillantes s'effondrent en un flot désordonné. C'est l'étape où la pensée devient \textbf{structure logique}, où chaque argument trouve sa place, sa fonction et son lien avec les autres.

\courssub{Les trois types de plans : choisir l'architecture adaptée}

Tout sujet de dissertation appelle une structure spécifique. Trois architectures majeures existent ; le choix ne relève pas du goût personnel, mais de la \textbf{nature du sujet} et de la \textbf{problématique} formulée.

\begin{definition}[Plan détaillé]
Le \textbf{plan détaillé} est la présentation écrite, hiérarchisée et numérotée de l'ensemble de la dissertation \emph{avant sa rédaction}. Il comporte : les titres des parties (I, II, III) et des sous-parties (A, B), rédigés en phrases affirmatives ; les arguments principaux numérotés (1, 2, 3) pour chaque sous-partie ; les exemples précis (références littéraires, faits de société, données) attachés à chaque argument ; les formules de transition entre les parties. Il sert de feuille de route au rédacteur et de grille d'évaluation pour le correcteur.
\end{definition}

\begin{popfigure}
\begin{tikzpicture}[
  level distance=1.8cm,
  sibling distance=3.5cm,
  every node/.style={font=\small, align=center, inner sep=3pt},
  edge from parent/.style={draw=popdark, thick, -{Stealth[length=3pt]}},
  planA/.style={fill=popblueL, draw=popblue, rounded corners, minimum width=2.8cm},
  planB/.style={fill=poporangeL, draw=poporange, rounded corners, minimum width=2.8cm},
  planC/.style={fill=popgreenL, draw=popgreen, rounded corners, minimum width=2.8cm},
  part/.style={fill=white, draw=popmuted, rectangle, minimum width=2.2cm},
  sub/.style={fill=popcream, draw=popline, rectangle, minimum width=1.8cm, font=\tiny}
]
% Plan Thématique
\begin{scope}[xshift=-11cm]
\node[planA] (T0) {PLAN THÉMATIQUE};
\node[part, below of=T0] (T1) {I. Aspect 1};
\node[part, below of=T1] (T2) {II. Aspect 2};
\node[part, below of=T2] (T3) {III. Aspect 3};
\draw[popblue, thick] (T0) -- (T1) (T1) -- (T2) (T2) -- (T3);
\node[sub, right of=T1, xshift=3.2cm] (T1A) {A. Sous-aspect};
\node[sub, right of=T1A, xshift=2.2cm] (T1B) {B. Sous-aspect};
\node[sub, right of=T2, xshift=3.2cm] (T2A) {A. Sous-aspect};
\node[sub, right of=T2A, xshift=2.2cm] (T2B) {B. Sous-aspect};
\node[sub, right of=T3, xshift=3.2cm] (T3A) {A. Sous-aspect};
\node[sub, right of=T3A, xshift=2.2cm] (T3B) {B. Sous-aspect};
\draw[popblue, dashed] (T1) -- (T1A) (T1A) -- (T1B);
\draw[popblue, dashed] (T2) -- (T2A) (T2A) -- (T2B);
\draw[popblue, dashed] (T3) -- (T3A) (T3A) -- (T3B);
\node[align=center, font=\tiny, text width=3.5cm, below of=T3, yshift=-0.8cm] {\textbf{Sujets types :} « Les fonctions du langage », « Les formes du théâtre », « La place de la femme dans \emph{Le Lion et la perle} »};
\end{scope}
% Plan Dialectique
\begin{scope}[xshift=0cm]
\node[planB] (D0) {PLAN DIALECTIQUE};
\node[part, below of=D0] (D1) {I. Thèse};
\node[part, below of=D1] (D2) {II. Antithèse};
\node[part, below of=D2] (D3) {III. Synthèse};
\draw[poporange, thick] (D0) -- (D1) (D1) -- (D2) (D2) -- (D3);
\node[sub, right of=D1, xshift=3.2cm] (D1A) {A. Argument pour};
\node[sub, right of=D1A, xshift=2.2cm] (D1B) {B. Argument pour};
\node[sub, right of=D2, xshift=3.2cm] (D2A) {A. Argument contre};
\node[sub, right of=D2A, xshift=2.2cm] (D2B) {B. Argument contre};
\node[sub, right of=D3, xshift=3.2cm] (D3A) {A. Dépassement};
\node[sub, right of=D3A, xshift=2.2cm] (D3B) {B. Nouvelle perspective};
\draw[poporange, dashed] (D1) -- (D1A) (D1A) -- (D1B);
\draw[poporange, dashed] (D2) -- (D2A) (D2A) -- (D2B);
\draw[poporange, dashed] (D3) -- (D3A) (D3A) -- (D3B);
\node[align=center, font=\tiny, text width=3.5cm, below of=D3, yshift=-0.8cm] {\textbf{Sujets types :} « Faut-il préserver les traditions ? », « Le progrès est-il toujours bénéfique ? », « La polygamie : droit ou obstacle ? »};
\end{scope}
% Plan Analytique
\begin{scope}[xshift=11cm]
\node[planC] (A0) {PLAN ANALYTIQUE};
\node[part, below of=A0] (A1) {I. Constat / Problème};
\node[part, below of=A1] (A2) {II. Causes / Analyse};
\node[part, below of=A2] (A3) {III. Solutions / Conséquences};
\draw[popgreen, thick] (A0) -- (A1) (A1) -- (A2) (A2) -- (A3);
\node[sub, right of=A1, xshift=3.2cm] (A1A) {A. Manifestation 1};
\node[sub, right of=A1A, xshift=2.2cm] (A1B) {B. Manifestation 2};
\node[sub, right of=A2, xshift=3.2cm] (A2A) {A. Cause interne};
\node[sub, right of=A2A, xshift=2.2cm] (A2B) {B. Cause externe};
\node[sub, right of=A3, xshift=3.2cm] (A3A) {A. Solution court terme};
\node[sub, right of=A3A, xshift=2.2cm] (A3B) {B. Solution long terme};
\draw[popgreen, dashed] (A1) -- (A1A) (A1A) -- (A1B);
\draw[popgreen, dashed] (A2) -- (A2A) (A2A) -- (A2B);
\draw[popgreen, dashed] (A3) -- (A3A) (A3A) -- (A3B);
\node[align=center, font=\tiny, text width=3.5cm, below of=A3, yshift=-0.8cm] {\textbf{Sujets types :} « Causes et conséquences de l'exode rural », « Comment lutter contre la corruption ? », « L'échec scolaire : diagnostic et remèdes »};
\end{scope}
\end{tikzpicture}
\legende{Comparaison visuelle des trois architectures de plan : thématique (aspects juxtaposés), dialectique (opposition puis dépassement), analytique (processus causalité $\rightarrow$ résolution).}
\end{popfigure}

\courssub{Critères de choix du plan selon le type de sujet}

Le tableau ci-dessous synthétise les indices linguistiques et logiques qui orientent le choix. \textbf{La règle d'or :} la problématique dicte le plan. Si votre problématique pose une tension (« Dans quelle mesure... », « Faut-il... »), le plan dialectique s'impose. Si elle demande une analyse de composantes (« Quels sont les aspects... »), le plan thématique convient. Si elle cherche un processus (« Pourquoi... Comment... »), le plan analytique est requis.

\begin{center}
\begin{tabularx}{\linewidth}{|l|X|X|X|}
\hline
\rowcolor{popblueL}
\textbf{Critère} & \textbf{Plan thématique} & \textbf{Plan dialectique} & \textbf{Plan analytique} \\
\hline
\textbf{Mots-clés du sujet} & « Les formes de... », « Les aspects de... », « Étudier... », « Caractéristiques de... » & « Faut-il... », « Pour ou contre... », « Dans quelle mesure... », « Est-il juste de... » & « Causes de... », « Conséquences de... », « Comment expliquer... », « Solutions pour... » \\
\hline
\textbf{Problématique type} & Quels sont les différents aspects de X ? & X est-il souhaitable / légitime / vrai ? & Pourquoi X ? Quelles en sont les causes / solutions ? \\
\hline
\textbf{Logique interne} & \textbf{Parallélisme} : les parties sont de même nature, indépendantes mais complémentaires. & \textbf{Confrontation} : thèse et antithèse s'opposent ; la synthèse les \textbf{dépasse} (niveau supérieur). & \textbf{Chronologie / Causalité} : I $\rightarrow$ II $\rightarrow$ III (enchaînement nécessaire). \\
\hline
\textbf{Piège à éviter} & Liste de courses sans fil conducteur ; parties étanches. & « Synthèse » qui n'est qu'un compromis mou (« les deux ont raison ») ; thèse/antithèse déséquilibrées. & Confondre causes et conséquences ; solutions sans lien avec les causes identifiées. \\
\hline
\textbf{Exemple (Le Lion et la perle)} & « Les figures du pouvoir dans la pièce » : I. Pouvoir traditionnel (Baroka), II. Pouvoir moderne (Lakunlé), III. Pouvoir féminin (Sidi, Sadiku). & « Faut-il préserver la tradition yoruba face à la modernité ? » : I. La tradition garantit l'identité, II. La tradition freine l'émancipation, III. Une modernité enracinée (ni rejet, ni mimétisme). & « Causes et conséquences du conflit générationnel » : I. Constat : affrontement Baroka/Lakunlé, II. Causes : choc des valeurs, rapport au temps, III. Conséquences : marginalisation des jeunes, rigidité du pouvoir. \\
\hline
\end{tabularx}
\end{center}

\courssub{Structure d'une partie et rédaction du plan détaillé}

Une partie de dissertation n'est pas un simple conteneur ; c'est une \textbf{unité de démonstration}. Chaque partie (I, II, III) défend une \textbf{idée directrice} (thèse partielle) qui répond à un pan de la problématique.

\begin{propriete}[Règle d'or du plan dialectique]
Dans un plan dialectique, la \textbf{synthèse (III) doit dépasser réellement l'opposition thèse/antithèse}, sans se réduire à un simple compromis ou à un « mélange des deux ». Elle doit apporter un \textbf{niveau de compréhension supérieur} : montrer que les termes de l'opposition étaient mal posés, proposer une conciliation dialectique (ex. « tradition vivante » vs « tradition fossilisée »), ou ouvrir une perspective nouvelle qui résout la tension. \emph{Exigible au Probatoire : la synthèse est notée sur sa capacité à dépasser l'alternative binaire.}
\end{propriete}

\begin{methode}[Rédiger un plan détaillé complet — 6 étapes]
1. \textbf{Formulez les trois titres de parties (I, II, III)} en phrases affirmatives, complètes, qui énoncent l'idée directrice (pas de simples noms : « I. La tradition » $\rightarrow$ « I. La tradition fonde l'identité collective »).\\
2. \textbf{Découpez chaque partie en deux sous-parties (A, B)} : chaque sous-partie = un argument distinct qui soutient l'idée directrice. Titres aussi en phrases affirmatives.\\
3. \textbf{Numérotez les arguments (1, 2, 3...)} sous chaque sous-partie : ce sont les raisonnements logiques (définition, cause, conséquence, comparaison, autorité).\\
4. \textbf{Attachez un exemple précis à chaque argument} : citation courte (auteur, œuvre, acte/scène), fait historique daté, statistique, référence culturelle (proverbe, rite). \emph{Au moins un exemple par sous-partie}.\\
5. \textbf{Rédigez les transitions} entre I/II et II/III : bilan de la partie qui s'achève + annonce de la partie qui vient (voir méthode transitions ci-dessous).\\
6. \textbf{Vérifiez l'équilibre} : même nombre de sous-parties par partie, même densité d'arguments, exemples variés (littéraires + société).
\end{methode}

\begin{exemplebox}[Plan détaillé complet commenté — Sujet : « Faut-il préserver les traditions ? » (Plan dialectique)]
\textbf{Problématique :} La tradition est-elle un héritage à conserver intact ou un fardeau à dépasser pour construire l'avenir ?

\textbf{I. La tradition est le socle de l'identité et de la cohésion sociale} (Thèse)\\
A. Elle transmet une mémoire collective et des repères symboliques\\
\quad 1. Rites d'initiation et proverbes structurent l'appartenance (ex. \emph{Le Lion et la perle} : le mime de la danse du voyageur, p. 45, montre la transmission orale)\\
\quad 2. La chefferie traditionnelle incarne la continuité historique (ex. Baroka, le Bale, « le renard » qui connaît « les sentiers de la brousse »)\\
B. Elle régule la vie collective par des normes partagées\\
\quad 1. Le mariage coutumier organise les alliances entre lignages (ex. négociation du prix de la mariée, Sadiku comme intermédiaire)\\
\quad 2. La justice traditionnelle (palabre) privilégie la réconciliation sur la peine (ex. scène du marché, règlement des différends)\\
\textbf{Transition I $\rightarrow$ II :} \emph{« Si la tradition cimente le groupe, elle peut aussi figer les rapports de force et exclure ceux qui ne correspondent pas aux normes établies. »}

\textbf{II. La tradition peut devenir un obstacle à l'émancipation individuelle et au progrès} (Antithèse)\\
A. Elle légitime des hiérarchies injustes (âge, sexe, naissance)\\
\quad 1. La gérontocratie muselle la parole des jeunes (ex. Lakunlé, bien que moderne, n'est pas écouté aux palabres)\\
\quad 2. La condition féminine : polygamie, lévirat, prix de la mariée = marchandisation (ex. Sidi refusant le prix, « je ne suis pas une bête de marché »)\\
B. Elle freine l'innovation technique et sociale par sacralisation de l'usage\\
\quad 1. Rejet de l'école occidentale vue comme corruption (ex. Lakunlé moqué pour ses « mots de livres »)\\
\quad 2. Refus des infrastructures modernes (ex. Baroka détourne le chemin de fer pour préserver son pouvoir)\\
\textbf{Transition II $\rightarrow$ III :} \emph{« Mais opposer tradition figée et modernité destructrice est un faux dilemme : les sociétés africaines ont toujours su adapter leurs coutumes. »}

\textbf{III. Vers une tradition vivante : réinventer l'héritage sans le rompre} (Synthèse)\\
A. La tradition n'est pas un bloc monolithique : distinguer l'essence (valeurs) de la forme (pratiques datées)\\
\quad 1. L'ubuntu / le « vivre-ensemble » peut s'exprimer hors de la polygamie (ex. Sadiku qui, tout en servant Baroka, manœuvre pour Sidi)\\
\quad 2. Les rites d'initiation peuvent intégrer l'éducation formelle (ex. projet d'école du village géré par les anciens)\\
B. La modernité africaine se construit par \textbf{relecture critique} et non par mimétisme\\
\quad 1. « Faire du neuf avec du vieux » : théâtre populaire yoruba (ex. la pièce elle-même, \emph{Le Lion et la perle}, est une forme moderne portée par la tradition orale)\\
\quad 2. Exemples contemporains : code de la famille au Cameroun (1981/2007) alliant droit moderne et coutume ; start-up « made in Africa » utilisant savoirs endogènes.
\end{exemplebox}

\courssub{L'introduction et la conclusion : les portes de la dissertation}

L'introduction et la conclusion ne sont pas des accessoires ; ce sont les \textbf{moments forts} où le correcteur évalue votre compréhension du sujet et votre maîtrise de la démonstration.

\begin{popfigure}
\begin{tikzpicture}[
  funnel/.style={draw=popblue, thick, fill=popblueL},
  txt/.style={font=\small, align=center, text width=5cm},
  arr/.style={-Stealth, thick, popdark}
]
% Entonnoir : 5 niveaux
\foreach \i/\lab/\w/\y in {
  1/{Amorce (accroche)\\ Général, actualité, citation}/7/4.5,
  2/{Présentation du sujet\\ Délimitation, contexte}/6/3.2,
  3/{Définition des termes\\ Clés du sujet}/5/1.9,
  4/{Problématique\\ Question tension}/4/0.6,
  5/{Annonce du plan\\ I, II, III}/3/-0.7} {
  \node[funnel, minimum width=\w cm, minimum height=0.9cm] (L\i) at (0,\y) {};
  \node[txt] at (0,\y) {\lab};
  \ifnum\i>1 \draw[arr] (L\the\numexpr\i-1\relax.south) -- (L\i.north); \fi
}
% Annotations latérales
\node[txt, xshift=5.5cm, yshift=4.5cm, align=center]{\textbf{But :} capter l'attention,\\montrer la pertinence};
\node[txt, xshift=5.5cm, yshift=3.2cm, align=center]{\textbf{But :} montrer qu'on a\\compris le sujet};
\node[txt, xshift=5.5cm, yshift=1.9cm, align=center]{\textbf{But :} lever l'ambiguïté,\\fixer le sens};
\node[txt, xshift=5.5cm, yshift=0.6cm, align=center]{\textbf{But :} poser LA question\\qui guide tout};
\node[txt, xshift=5.5cm, yshift=-0.7cm, align=center]{\textbf{But :} donner la carte,\\rassurer le correcteur};
\end{tikzpicture}
\legende{Structure en entonnoir de l'introduction : du général (amorce) au particulier (annonce du plan). Chaque niveau remplit une fonction rhétorique précise.}
\end{popfigure}

\begin{aretenir}[Structure type de l'introduction (5 mouvements obligatoires)]
1. \textbf{Amorce (accroche) :} 1--2 phrases. Citation d'auteur, fait d'actualité, proverbe, paradoxe, statistique choc. \emph{Éviter :} « De tout temps... », « Depuis la nuit des temps... » (banalités).\\
2. \textbf{Présentation du sujet :} Reprise du libellé, délimitation (espace, époque, corpus), mise en contexte.\\
3. \textbf{Définition des termes clés :} Sens précis dans le sujet (dictionnaire + contextualisation). Ex. « Tradition » : coutumes transmises, mais aussi « transmission » (étymologie).\\
4. \textbf{Problématique :} Question unique, tension, qui ne se répond pas par oui/non. « Dans quelle mesure... », « Comment concilier... ».\\
5. \textbf{Annonce du plan :} « Nous verrons d'abord que... (I), puis que... (II), pour enfin montrer que... (III). » \emph{Pas de numéros romains seuls : phrases complètes.}
\end{aretenir}

\begin{aretenir}[Structure de la conclusion (2 mouvements + ouverture)]
1. \textbf{Bilan (réponse à la problématique) :} 3--4 phrases. Reprendre les idées forces de I, II, III pour formuler la réponse synthétique. \emph{Ne pas ajouter d'argument nouveau.}\\
2. \textbf{Ouverture :} Une piste de réflexion élargie : autre œuvre, autre époque, autre discipline, actualité brûlante, question éthique. \emph{Ex. :} « Cette réflexion sur la tradition rejoint le débat actuel sur la restitution du patrimoine africain : comment préserver sans figer ? »
\end{aretenir}

\courssub{Les transitions : l'articulation fluide de la démonstration}

La transition n'est pas une politesse ; c'est la \textbf{preuve que votre pensée progresse}. Elle se place \textbf{à la fin de la partie} (dernier paragraphe de I, puis de II), jamais au début de la suivante.

\begin{methode}[Rédiger une transition efficace — Formule en 2 temps]
\textbf{Temps 1 : Bilan (rétrospectif).} Résumez en une phrase l'apport essentiel de la partie qui s'achève.\\
\textbf{Temps 2 : Annonce (prospectif).} Montrez l'insuffisance ou la limite de ce bilan pour justifier la partie qui vient.\\
\textbf{Modèle :} « Si [idée force de la partie achevée], [limite / angle mort / conséquence] ; c'est pourquoi il faut désormais examiner [idée directrice de la partie suivante]. »\\
\textbf{Exemple (I $\rightarrow$ II, sujet traditions) :} \emph{« Si la tradition fonde l'identité collective par des rites et des normes partagés, elle peut aussi figer des hiérarchies injustes et freiner l'émancipation individuelle ; c'est pourquoi il faut désormais analyser les limites et les dérives de l'ordre coutumier. »}
\end{methode}

\begin{attention}[Erreurs fréquentes sur les transitions]
\begin{itemize}
\item \textbf{Transition « liste » :} « Maintenant, passons à la deuxième partie. » (Zéro point : aucune logique).
\item \textbf{Transition « résumé seulement » :} « Nous avons vu que la tradition est importante. » (Pas d'appel à la suite).
\item \textbf{Transition « contradiction brute » :} « Mais la tradition est mauvaise. » (Pas de nuance, pas de lien logique).
\item \textbf{Oubli de transition :} Partie I se termine, partie II commence sans lien. (Rupture de cohérence, pénalité forte).
\end{itemize}
\end{attention}

\courssub{Confrontation et amélioration : grille d'auto-évaluation et d'évaluation croisée}

Avant de rédiger, le plan détaillé doit passer au crible d'une grille rigoureuse. En classe, on pratique l'\textbf{évaluation croisée} : chaque élève échange son plan avec un pair, applique la grille, puis revient vers son auteur pour discussion. Le professeur valide in fine.

\begin{center}
\begin{tabularx}{\linewidth}{|l|X|X|>{\columncolor{popgreenL}}X|}
\hline
\rowcolor{popblueL}
\textbf{Critère} & \textbf{Indicateurs de réussite} & \textbf{Signaux d'alerte} & \textbf{Action corrective} \\
\hline
\textbf{Respect du sujet} & Tous les termes du sujet sont traités ; aucun hors-sujet ; problématique issue du sujet. & Mots du sujet ignorés ; problématique hors champ ; exemples sans lien. & Re-formuler la problématique ; rayer les arguments hors-sujet ; en chercher de nouveaux. \\
\hline
\textbf{Cohérence logique} & Plan adapté au type de sujet ; enchaînement I $\rightarrow$ II $\rightarrow$ III nécessaire ; synthèse qui dépasse (dialectique). & Plan thématique pour un sujet « Faut-il... » ; synthèse = compromis ; causes/conséquences inversées. & Changer de type de plan ; réécrire la synthèse ; réordonner les parties. \\
\hline
\textbf{Équilibre des parties} & 3 parties $\times$ 2 sous-parties ; densité arguments/exemples équivalente ; \textasciitilde 300--400 mots par partie à la rédaction. & I : 4 sous-parties, 10 arguments ; II : 1 sous-partie, 2 arguments ; III : inexistante. & Fusionner/scinder les sous-parties ; transférer des arguments ; compléter le vide. \\
\hline
\textbf{Qualité des exemples} & Au moins 1 exemple précis par sous-partie ; variés (littéraire, historique, social, personnel) ; correctement référencés. & Exemples vagues (« on dit que... »), tous littéraires, ou absents ; erreurs de référence (acte, page). & Rechercher 2 exemples concrets par sous-partie ; vérifier les citations à l'ouvroir. \\
\hline
\textbf{Transitions} & 2 transitions rédigées (I/II, II/III) ; formule bilan + annonce ; vocabulaire de liaison logique. & Transitions absentes, ou « passons à... », ou simples résumés. & Rédiger les 2 transitions selon la méthode 2 temps ; les insérer dans le plan. \\
\hline
\textbf{Présentation / Lisibilité} & Hiérarchie claire (I, A, 1) ; titres en phrases affirmatives ; police lisible ; pas de ratures ; annexé proprement. & Numérotation confuse ; titres nominaux (« I. Tradition ») ; plan illisible, brouillon. & Recopier au propre sur feuille d'annexe ; respecter la mise en page type (marges, interlignes). \\
\hline
\end{tabularx}
\end{center}

\begin{savaistu}[Le plan détaillé au Probatoire / Baccalauréat technique]
Lors de l'épreuve écrite de Français (Probatoire), \textbf{le plan détaillé peut être exigé en annexe} et fait l'objet d'une notation spécifique (souvent 2 à 4 points sur 20). Le correcteur vérifie : (1) la pertinence du type de plan par rapport au sujet, (2) la cohérence interne (thèse/antithèse/synthèse ou causes/solutions), (3) la présence d'exemples précis et variés, (4) la rédaction des transitions. \textbf{Conseil stratégique :} consacrez 20--25 minutes sur les 3h--3h30 à l'élaboration et à la recopie soignée du plan détaillé sur la feuille d'annexe fournie. Un plan propre, complet et logique « rassure » le correcteur avant même la lecture de la dissertation ; un plan bâclé ou absent crée un a priori négatif difficile à rattraper.
\end{savaistu}

\begin{exoresolu}[Exercice : Choisir et justifier le type de plan]
\textbf{Consigne :} Pour chacun des sujets ci-dessous, indiquez le type de plan le plus adapté (thématique, dialectique, analytique) et justifiez en une phrase en vous appuyant sur les mots-clés du sujet.

1. « Les fonctions du langage dans \emph{Le Lion et la perle} de Wole Soyinka. »\\
2. « Faut-il interdire la polygamie au Cameroun ? »\\
3. « Causes et conséquences de l'exode des jeunes vers les villes. »\\
4. « Étudiez le personnage de Baroka comme figure du pouvoir traditionnel. »\\
5. « La modernité est-elle une chance pour l'Afrique ? »

\tcblower
\textbf{Solution détaillée :}\\
1. \textbf{Thématique.} Mots-clés : « Les fonctions » (pluriel, énumération d'aspects) $\rightarrow$ étude de composantes indépendantes (fonction phatique, référentielle, poétique, performative dans la pièce).\\
2. \textbf{Dialectique.} « Faut-il interdire » = question de valeur, tension éthique $\rightarrow$ thèse (arguments pour : droits femmes, égalité) / antithèse (arguments contre : liberté culturelle, régulation sociale) / synthèse (cadre légal protecteur sans interdiction brute).\\
3. \textbf{Analytique.} « Causes et conséquences » = processus causal explicite $\rightarrow$ I. Constat (ampleur exode), II. Causes (pauvreté rurale, attractivité urbaine, éducation), III. Conséquences (bidonvilles, vieillissement campagnes, main-d'œuvre urbaine) + solutions.\\
4. \textbf{Thématique (ou analytique personnage).} « Étudiez le personnage » $\rightarrow$ analyse de facettes : I. Baroka stratège politique, II. Baroka gardien de la tradition, III. Baroka séducteur manipulateur. (Plan thématique : aspects du personnage).\\
5. \textbf{Dialectique.} « Est-elle une chance » = évaluation, tension $\rightarrow$ I. Chance (technique, santé, droits), II. Risque (aliénation, dépendance, perte identitaire), III. Condition : appropriation critique / modernité endogène (ex. Silicon Valley africain, Nollywood).
\end{exoresolu}

\begin{exercice}[À vous de jouer : Plan détaillé à construire]
\textbf{Sujet :} « Le théâtre est-il le miroir de la société ? »\\
\textbf{Consignes :}
\begin{enumerate}
\item Formulez une problématique tension (une seule phrase, point d'interrogation).
\item Choisissez le type de plan justifié (1 phrase).
\item Rédigez les titres des 3 parties (I, II, III) en phrases affirmatives.
\item Pour la partie I seulement : écrivez les 2 sous-parties (A, B), 2 arguments numérotés par sous-partie, et 1 exemple précis par argument (référence \emph{Le Lion et la perle} ou autre).
\item Rédigez la transition I $\rightarrow$ II selon la méthode « bilan + annonce ».
\end{enumerate}
\textit{Temps conseillé : 25 minutes. Échangez ensuite avec un pair pour évaluation croisée (grille ci-dessus).}
\end{exercice}

\begin{corrige}[Exercice n°1 — Éléments de correction]
\textbf{Problématique type :} « Dans quelle mesure le théâtre reflète-t-il les réalités sociales tout en les transformant par la représentation ? »\\
\textbf{Type de plan :} \textbf{Dialectique} (question « est-il » = évaluation, tension reflet / transformation).\\
\textbf{I. Le théâtre reflète fidèlement les structures, conflits et langages de son époque.}\\
\textbf{II. Le théâtre déforme, grossit ou réinvente le réel par les codes de la scène (fiction, catharsis, esthétique).}\\
\textbf{III. Le théâtre devient ainsi un laboratoire critique où la société se voit et se pense pour mieux se transformer.}\\
\textbf{Détail partie I :}\\
A. La scène reproduit les hiérarchies et rapports de force réels\\
\quad 1. Représentation du pouvoir traditionnel (ex. Baroka, le Bale, incarne l'autorité gérontocratique yoruba réelle).\\
\quad 2. Mise en scène des rapports de genre (ex. Sadiku, femme senior, manipule mais reste soumise au Bale).\\
B. Le théâtre intègre la langue, les proverbes, les rites du vécu quotidien\\
\quad 1. Emploi de l'anglais pidgin / yoruba translittéré (ex. chants, mimes, « \emph{Il n'est pas bon de laisser le lion affamé} »).\\
\quad 2. Reconstitution de l'espace villageois (marché, case, place du palabre) comme microcosme social.\\
\textbf{Transition I $\rightarrow$ II :} \emph{« Si le théâtre offre un reflet saisissant des structures sociales par ses personnages, sa langue et son décor, il ne saurait se réduire à une simple photographie : la mise en scène, par ses artifices, transforme le réel en le portant à une portée universelle. »}
\end{corrige}

\coursec{Outils linguistiques de la dissertation}

Vous savez désormais analyser un sujet, formuler une problématique, rassembler vos idées et les organiser en un plan détaillé. Mais un plan, même excellent, ne suffit pas : il faut encore le \emph{rédiger}. Or, entre l'idée que vous avez dans la tête et la phrase que lit votre correcteur, il y a des outils de langue précis : la \cle{ponctuation}, les marques d'\cle{énonciation}, le choix des mots (\cle{dénotation} et \cle{connotation}), la \cle{description} et les \cle{connecteurs logiques}. Cette section vous donne ces outils, un par un, avec leurs fonctions et leurs pièges.

\courssub{La ponctuation dans l'écriture argumentative}

\textbf{Intuition.} Lisez à voix haute cette phrase : « L'école forme les citoyens elle doit donc être accessible à tous. » Vous avez buté ? Normal : sans ponctuation, le lecteur ne sait pas où s'arrête une idée et où commence la suivante. La ponctuation est le \emph{code de la route} de votre raisonnement : elle guide le lecteur, rythme votre pensée et hiérarchise vos idées.

\begin{definition}[La ponctuation]
La \cle{ponctuation} est l'ensemble des signes graphiques (point, virgule, point-virgule, deux-points, guillemets, parenthèses\ldots) qui structurent la phrase, marquent les pauses et précisent les rapports logiques entre les idées.
\end{definition}

\begin{popfigure}
\begin{center}
\begin{tikzpicture}[
  signe/.style={rectangle, draw=popblue, fill=popblueL, rounded corners=2pt, minimum width=2.2cm, minimum height=0.7cm, font=\bfseries, align=center},
  fonc/.style={rectangle, draw=popline, fill=popcream, rounded corners=2pt, minimum width=3.6cm, minimum height=0.7cm, font=\small, align=center},
  ex/.style={rectangle, draw=popline, fill=white, rounded corners=2pt, minimum width=6.6cm, minimum height=0.7cm, font=\small\itshape, align=center}
]
\node[signe] (s1) at (0,0) {Le point ( . )};
\node[fonc] (f1) at (3.4,0) {Clôt une idée complète ; marque une pause forte};
\node[ex] (e1) at (8.9,0) {L'école forme les citoyens. Elle doit donc être accessible à tous.};

\node[signe] (s2) at (0,-1.1) {La virgule ( , )};
\node[fonc] (f2) at (3.4,-1.1) {Sépare les éléments d'une énumération, isole une incise};
\node[ex] (e2) at (8.9,-1.1) {Certes, l'école coûte cher, mais elle enrichit la nation.};

\node[signe] (s3) at (0,-2.2) {Le point-virgule ( ; )};
\node[fonc] (f3) at (3.4,-2.2) {Relie deux idées proches et parallèles};
\node[ex] (e3) at (8.9,-2.2) {L'instruction libère l'individu ; elle libère aussi la société.};

\node[signe] (s4) at (0,-3.3) {Les deux-points ( : )};
\node[fonc] (f4) at (3.4,-3.3) {Annoncent une explication, une citation, une liste};
\node[ex] (e4) at (8.9,-3.3) {Une évidence s'impose : sans éducation, pas de développement.};

\node[signe] (s5) at (0,-4.4) {Les guillemets ( « » )};
\node[fonc] (f5) at (3.4,-4.4) {Encadrent une citation ou un mot employé au figuré};
\node[ex] (e5) at (8.9,-4.4) {Soyinka écrit : « La jeunesse est l'avenir du village. »};

\node[signe] (s6) at (0,-5.5) {Les parenthèses ( ( ) )};
\node[fonc] (f6) at (3.4,-5.5) {Insèrent une précision secondaire};
\node[ex] (e6) at (8.9,-5.5) {Le théâtre (art total) touche tous les publics.};
\end{tikzpicture}
\end{center}
\legende{Les principaux signes de ponctuation : fonction et exemple argumentatif.}
\end{popfigure}

\textbf{Ponctuation et articulation du raisonnement.} Deux signes sont particulièrement précieux dans une dissertation :

\begin{itemize}
\item \textbf{Les deux-points pour annoncer.} Ils créent une attente chez le lecteur et introduisent ce qui justifie votre affirmation : explication, preuve, citation. Exemple : « La tradition joue un rôle essentiel dans le roman : elle garantit la cohésion du village. »
\item \textbf{Le point-virgule pour relier des idées proches.} Il unit deux phrases complètes qui expriment un parallélisme ou un équilibre : « Baroka défend la tradition ; Sadiku incarne la ruse féminine. » Il est très utile dans les développements symétriques (thèse / antithèse).
\end{itemize}

\begin{attention}[Erreurs fréquentes de ponctuation]
\begin{itemize}
\item \textbf{Jamais de virgule entre le sujet et son verbe.} On écrit « La jeunesse du village rêve de modernité », et non « La jeunesse du village, rêve de modernité ». La virgule n'est admise que si une incise complète s'intercale : « La jeunesse du village, séduite par la ville, rêve de modernité. »
\item \textbf{Le point-virgule n'est pas une virgule renforcée.} Il ne peut séparer que deux \emph{propositions complètes} (chacune avec sujet et verbe conjugué). « Lakunle veut moderniser le village ; sans renoncer à ses idéaux » est fautif : la seconde partie n'est pas une phrase complète.
\end{itemize}
\end{attention}

\courssub{L'énonciation : les marques explicites}

\textbf{Intuition.} Comparez : « On constate que la tradition résiste » et « Je pense que la tradition résiste ». Dans les deux cas, quelqu'un parle ; mais ce « quelqu'un » ne se place pas de la même manière face à son lecteur. Toute phrase porte la trace de celui qui l'énonce : c'est l'\cle{énonciation}.

\begin{definition}[L'énonciation]
L'\cle{énonciation} est l'ensemble des marques linguistiques par lesquelles l'énonciateur (celui qui parle ou écrit) s'inscrit dans son discours et situe son message par rapport au destinataire, au temps et au lieu.
\end{definition}

\textbf{Les marques explicites.}

\begin{itemize}
\item \textbf{Marques de la personne} : \emph{je} (engagement personnel : « je soutiens que\ldots »), \emph{nous} (association du lecteur : « nous devons reconnaître que\ldots »), \emph{on} (généralisation, effet d'objectivité : « on observe que\ldots »).
\item \textbf{Marques de temps et de lieu} : \emph{aujourd'hui, désormais, ici, dans notre société, de nos jours}. Elles ancrent le propos dans l'actualité : « Aujourd'hui, dans nos villes camerounaises, la question se pose avec acuité. »
\item \textbf{Les modalisateurs} : mots et tournures qui expriment le degré de certitude de l'énonciateur : \emph{sans doute, certes, il est clair que, il semble que, peut-être, assurément}. « Il est clair que » affirme avec force ; « sans doute » laisse une ouverture à la discussion.
\end{itemize}

\begin{popfigure}
\begin{center}
\begin{tikzpicture}[
  actant/.style={rectangle, draw=popblue, fill=popblueL, rounded corners=3pt, minimum width=3.4cm, minimum height=1.1cm, align=center, font=\small\bfseries},
  marque/.style={font=\small\itshape, text=popdark, align=center},
  fleche/.style={-{Stealth[length=3mm]}, thick, poporange}
]
\node[actant, align=center] (enonc) at (0,0){\'Enonciateur\\ (je / nous / on)};
\node[actant, align=center] (dest) at (10,0){Destinataire\\ (le lecteur, le correcteur)};
\node[actant, fill=popgreenL, draw=popgreen, align=center] (sit) at (5,-3){Situation de communication\\ (temps : aujourd'hui ; lieu : ici)};

\draw[fleche] (enonc) -- node[marque, above, align=center]{message argumenté\\ modalisateurs : \emph{certes, sans doute,}\\ \emph{il est clair que}} (dest);
\draw[fleche, popblue] (enonc) -- node[marque, left=2pt, align=center]{s'engage\\ (\emph{je pense que})} (sit);
\draw[fleche, popblue] (dest) -- node[marque, right=2pt, align=center]{est pris en compte\\ (\emph{nous constatons})} (sit);
\end{tikzpicture}
\end{center}
\legende{Schéma de l'énonciation : l'énonciateur adresse un message au destinataire dans une situation marquée par le temps et le lieu ; les flèches indiquent les marques explicites.}
\end{popfigure}

\textbf{Effets des marques d'énonciation.} Le choix de la personne n'est jamais neutre :

\begin{itemize}
\item \emph{Je} implique fortement l'énonciateur : il assume sa thèse. À utiliser avec mesure, surtout dans l'introduction (annonce du plan) et la conclusion.
\item \emph{Nous} crée une complicité avec le lecteur : il l'associe au raisonnement et le rend complice de la conclusion. C'est la marque la plus sûre dans le développement.
\item \emph{On} et les tournures impersonnelles (\emph{il apparaît que}) donnent au propos une allure objective, comme si l'idée s'imposait d'elle-même.
\end{itemize}

\begin{exemplebox}[Trois formulations, trois effets]
\begin{itemize}
\item « \emph{Je} trouve que Baroka est un personnage ambigu. » $\rightarrow$ opinion personnelle assumée.
\item « \emph{Nous} voyons bien que Baroka est un personnage ambigu. » $\rightarrow$ le lecteur est invité à partager le constat.
\item « \emph{Il est clair que} Baroka est un personnage ambigu. » $\rightarrow$ l'idée est présentée comme une évidence indiscutable.
\end{itemize}
\end{exemplebox}

\courssub{Dénotation et connotation : choisir ses mots}

\textbf{Intuition.} Un « vieillard » et un « ancien » désignent la même personne âgée. Pourtant, dire de Baroka qu'il est un « vieillard rusé » ou un « ancien respecté » ne produit pas le même effet sur le lecteur. Le sens objectif est identique ; le sens suggéré diffère.

\begin{definition}[Dénotation et connotation]
La \cle{dénotation} est le sens stable et objectif d'un mot, celui que donne le dictionnaire, partagé par tous les locuteurs. La \cle{connotation} est l'ensemble des valeurs affectives, culturelles ou imagées qu'un mot évoque en plus de son sens propre.
\end{definition}

\begin{center}
\begin{tabularx}{\linewidth}{|l|X|X|}
\hline
\rowcolor{popblueL} \textbf{Mot} & \textbf{Dénotation (sens objectif)} & \textbf{Connotations (sens suggéré)} \\
\hline
vieillard & homme âgé & faiblesse, déclin (péjoratif) \\
\hline
ancien / sage & homme âgé & expérience, respect, autorité (mélioratif) \\
\hline
modernité & fait d'être moderne & progrès, avenir (positive) ou rupture, perte des valeurs (négative selon le contexte) \\
\hline
tradition & transmission des coutumes & enracinement, sagesse (positive) ou immobilisme, archaïsme (négative selon le contexte) \\
\hline
\end{tabularx}
\end{center}

\textbf{Le choix du vocabulaire selon la thèse défendue.} Dans une dissertation, votre lexique doit servir votre argumentation : si vous défendez la tradition, vous parlerez d'« héritage », de « racines », de « sagesse ancestrale » ; si vous la critiquez, vous parlerez d'« immobilisme », de « poids du passé », de « coutumes désuètes ». Le fond reste le même ; les connotations orientent le jugement du lecteur. Attention toutefois à rester mesuré : un vocabulaire trop chargé donne l'impression d'un parti pris excessif et affaiblit la démonstration.

\courssub{Le texte descriptif au service de l'argumentation}

\textbf{Intuition.} Une dissertation n'est pas faite que d'idées abstraites. Pour prouver, il faut parfois \emph{montrer}. Une description bien menée transforme un exemple vague en preuve concrète : le lecteur \emph{voit} ce que vous affirmez.

\begin{definition}[Le texte descriptif]
Le \cle{texte descriptif} présente un personnage, un lieu ou un objet en le rendant visible au lecteur. Ses composantes sont : un \textbf{cadre spatio-temporel} (où ? quand ?), une \textbf{sélection des détails} (on ne décrit pas tout, on choisit ce qui sert l'effet visé) et l'\textbf{expansion du nom} (adjectifs qualificatifs, compléments du nom, propositions subordonnées relatives qui enrichissent les noms).
\end{definition}

\begin{exemplebox}[Une description qui argumente]
Pour illustrer l'idée « le village d'Ilujinle vit encore au rythme de la tradition », on peut écrire : « À l'entrée du village, sous le grand baobab \emph{qui domine la place}, les anciens, \emph{vêtus de leurs boubous aux couleurs éclatantes}, délibèrent \emph{depuis l'aube} autour du chef. » Le cadre (entrée du village, place), les détails choisis (baobab, anciens, délibération) et les expansions du nom (subordonnée relative, groupe adjectival, complément de temps) rendent la tradition visible : la description devient un \emph{argument}.
\end{exemplebox}

\textbf{Rôle de la description dans l'argumentation.} Dans une dissertation, la description n'est jamais gratuite : elle sert à \emph{développer un exemple}. Au lieu d'écrire simplement « la pauvreté pousse les jeunes à quitter l'école », vous pouvez décrire brièvement une salle de classe délabrée ou un élève vendeur à la sauvette : l'exemple gagne en force persuasive. La règle d'or : toute description doit être \emph{suivie d'une phrase d'analyse} qui la relie à l'argument (« Cette scène montre que\ldots »).

\courssub{Les connecteurs logiques de la dissertation}

\textbf{Intuition.} Vos idées sont des pierres ; les connecteurs logiques sont le ciment. Sans eux, le lecteur voit des affirmations juxtaposées ; avec eux, il suit un \emph{raisonnement}.

\begin{definition}[Connecteur logique]
Un \cle{connecteur logique} est un mot ou une expression qui explicite la relation entre deux idées : addition, opposition, cause, conséquence, concession, conclusion.
\end{definition}

\begin{methode}[Choisir le bon connecteur logique]
\begin{enumerate}
\item \textbf{Identifiez la relation} entre l'idée que vous venez d'énoncer et celle que vous allez énoncer : ajoutez-vous une idée ? vous opposez-vous ? expliquez-vous ? concluez-vous ?
\item \textbf{Reportez-vous au tableau des relations} : addition (\emph{d'abord, ensuite, de plus, en outre}), explication (\emph{en effet, car, c'est-à-dire}), opposition (\emph{cependant, mais, en revanche, pourtant}), cause (\emph{parce que, puisque}), conséquence (\emph{par conséquent, donc, ainsi, c'est pourquoi}), concession (\emph{certes\ldots mais, bien sûr\ldots cependant}), conclusion (\emph{en définitive, en somme, finalement}).
\item \textbf{Vérifiez la cohérence} : relisez les deux phrases liées par le connecteur. Si la relation annoncée n'existe pas réellement (par exemple \emph{donc} sans lien de conséquence), changez de connecteur.
\item \textbf{Variez les connecteurs} : ne répétez pas \emph{mais} ou \emph{donc} à chaque phrase ; utilisez les synonymes.
\end{enumerate}
\end{methode}

\begin{exoresolu}[Réécrire un paragraphe argumentatif fautif]
Énoncé. Voici un paragraphe d'élève, truffé de fautes de ponctuation et de connecteurs mal employés. Réécrivez-le correctement.

« Lakunle, le maître d'école veut moderniser le village, il refuse de payer le prix de la dot. Sidi refuse sa demande, donc elle est vaniteuse. Baroka est vieux ; étant le chef du village. La tradition est forte, mais elle résiste à la modernité, en définitive le conflit oppose deux visions du monde. »

\tcblower
Solution détaillée. Version corrigée :

« Lakunle, le maître d'école, veut moderniser le village : il refuse notamment de payer le prix de la dot. Sidi refuse sa demande, \emph{car} elle se sent rabaissée. Baroka, lui, est vieux ; il reste pourtant le chef du village. La tradition est forte ; \emph{cependant}, la modernité progresse. \emph{En définitive}, le conflit oppose deux visions du monde. »

Commentaires des corrections :
\begin{itemize}
\item « Lakunle, le maître d'école, veut » : l'incise « le maître d'école » est encadrée par \emph{deux} virgules ; une seule virgule séparerait fautivement le sujet du verbe.
\item « village : il refuse » : les deux-points annoncent une explication de la volonté de moderniser.
\item « \emph{car} elle se sent rabaissée » : la version fautive utilisait \emph{donc} (conséquence) là où il fallait une \emph{cause} ; de plus, « elle est vaniteuse » jugeait sans expliquer.
\item « il reste pourtant le chef du village » : le point-virgule ne peut pas introduire « étant le chef » (participe, pas de phrase complète) ; on rétablit une proposition complète.
\item « \emph{cependant} » remplace le \emph{mais} qui opposait la tradition à elle-même (contresens) ; « en définitive » est déplacé pour introduire la vraie conclusion.
\end{itemize}
\end{exoresolu}

\begin{savaistu}[« Certes\ldots mais » : l'outil de concession roi]
Dans une partie dialectique (celle où vous examinez la thèse adverse avant de la dépasser), le correcteur attend presque toujours la structure de concession « \emph{certes\ldots mais} » : « \emph{Certes}, la modernité apporte le confort et le progrès technique ; \emph{mais} elle menace les valeurs qui fondent l'identité du village. » Le premier membre de phrase reconnaît honnêtement le point de vue adverse (ce qui montre votre ouverture d'esprit) ; le second le dépasse (ce qui montre votre capacité à trancher). Maîtriser ce schéma, c'est tenir la clé de la troisième partie de la dissertation dialectique, celle de la synthèse.
\end{savaistu}

\begin{aretenir}[L'essentiel de la section]
\begin{itemize}
\item La ponctuation structure le raisonnement : les deux-points \emph{annoncent}, le point-virgule \emph{relie} des idées proches, les guillemets encadrent les citations. Jamais de virgule entre le sujet et le verbe.
\item Les marques d'énonciation (personnes, temps, lieu, modalisateurs) inscrivent l'énonciateur dans son discours et ménagent le lecteur ; \emph{nous} est la marque la plus sûre dans le développement.
\item Un mot a un sens objectif (dénotation) et un sens suggéré (connotation) : choisissez votre vocabulaire en fonction de la thèse défendue.
\item La description, dans une dissertation, sert à développer un exemple ; elle doit toujours être suivie d'une analyse.
\item Les connecteurs logiques explicite­nt les relations entre les idées : addition, explication, opposition, cause, conséquence, concession, conclusion. Choisissez-les en fonction de la relation réelle, et variez-les.
\end{itemize}
\end{aretenir}

\begin{exercice}[À vous de jouer]
\begin{enumerate}
\item Ponctuez correctement ce passage : « Baroka le chef du village incarne la tradition il est rusé puissant et respecté certes il est âgé mais son autorité demeure intacte ».
\item Dans la phrase « On constate aujourd'hui que les jeunes délaissent sans doute les coutumes », relevez toutes les marques d'énonciation et précisez leur effet.
\item Remplacez les mots soulignés par des termes à connotation opposée : « Le \emph{sage} Baroka défend l'\emph{héritage} de ses pères. »
\item Reliez les paires d'idées suivantes par le connecteur logique approprié : (a) L'école coûte cher. / Elle est un investissement pour la nation. (b) Sidi est belle. / Baroka veut l'épouser. (c) Lakunle refuse la dot. / Sidi le rejette.
\end{enumerate}
\end{exercice}

\begin{corrige}[Exercice « À vous de jouer »]
\begin{enumerate}
\item « Baroka, le chef du village, incarne la tradition : il est rusé, puissant et respecté. Certes, il est âgé, mais son autorité demeure intacte. » (Incise entre deux virgules ; deux-points annonçant l'explication ; virgules d'énumération ; structure de concession « certes\ldots mais ».)
\item Marques : \emph{on} (personne, effet d'objectivité), \emph{aujourd'hui} (marque de temps, ancrage dans l'actualité), \emph{sans doute} (modalisateur, certitude nuancée qui laisse place à la discussion).
\item Version à connotation négative : « Le \emph{vieillard} Baroka s'accroche aux \emph{vieilles coutumes} de ses pères. » (Autres choix possibles : « ringard », « immobilisme », « passéisme ».)
\item (a) « L'école coûte cher ; \emph{cependant}, elle est un investissement pour la nation » (opposition). (b) « Sidi est belle ; \emph{c'est pourquoi} Baroka veut l'épouser » (conséquence). (c) « Lakunle refuse la dot ; \emph{par conséquent}, Sidi le rejette » (conséquence) ou « \emph{Puisque} Lakunle refuse la dot, Sidi le rejette » (cause).
\end{enumerate}
\end{corrige}

Ces outils linguistiques maîtrisés, il reste à vérifier qu'ils fonctionnent dans une rédaction complète : c'est l'objet de la dernière section, consacrée au compte rendu, à la remédiation et au projet d'études autour du \emph{Lion et la perle}.

\coursec{Compte rendu, remédiation et projet d'études}

Après avoir maîtrisé les outils linguistiques de la dissertation — ponctuation, énonciation, description, dénotation/connotation — nous abordons l'étape finale de l'unité : l'évaluation de nos acquis, la correction de nos faiblesses et la projection vers l'étude approfondie de \emph{Le Lion et la perle}. Cette section transforme l'apprenant en acteur réflexif de sa progression.

\courssub{Contrôle de lecture : \emph{Le Lion et la perle}}

\begin{definition}[Contrôle de lecture]
Vérification systématique de la compréhension globale d'une œuvre littéraire à travers l'identification des personnages, le suivi de l'intrigue et la saisie des enjeux thématiques.
\end{definition}

Le contrôle de lecture ne vise pas la mémorisation de détails anecdotiques, mais la construction d'une \cle{représentation mentale} cohérente de l'œuvre. Pour \emph{Le Lion et la perle} de Wole Soyinka, trois pôles structurent cette représentation :

\begin{itemize}
    \item \textbf{Personnages principaux} : Baroka (le Bale, le « lion », chef rusé et gardien de la tradition), Lakunle (instituteur moderniste, occidentalisé), Sidi (la « perle », beauté du village convoitée par les deux hommes), Sadiku (épouse aînée de Baroka, messagère manipulée malgré elle).
    \item \textbf{Intrigue en trois temps} : (1) Le matin — la modernité arrive au village (le photographe, le magazine qui a rendu Sidi célèbre) ; (2) Le midi — la ruse de Baroka, qui feint l'impuissance pour attirer Sidi ; (3) Le soir — le triomphe de la tradition : Sidi choisit Baroka, Lakunle reste seul avec ses discours.
    \item \textbf{Enjeux} : confrontation tradition/modernité, pouvoir de la parole et de l'image, place de la femme, ruse contre naïveté, théâtre comme miroir social.
\end{itemize}

\begin{methode}[Vérifier sa compréhension globale]
\begin{enumerate}
    \item Rédiger en 10 lignes maximum le résumé de l'intrigue (qui fait quoi, dans quel ordre, avec quelles conséquences).
    \item Dresser le portrait croisé de Baroka et Lakunle : deux visions du monde, deux stratégies de séduction.
    \item Identifier trois scènes-clés où le conflit tradition/modernité s'incarne physiquement (corps, langage, objets).
    \item Formuler en une phrase la thèse implicite de l'œuvre sur le changement social.
\end{enumerate}
\end{methode}

\begin{exoresolu}[Contrôle de lecture express]
\textbf{Énoncé :} En 8 lignes, résumer l'intrigue de \emph{Le Lion et la perle} en mettant en valeur le mécanisme de la ruse de Baroka.

\textbf{Solution :}
L'instituteur Lakunle courtise Sidi selon des rites modernes (refus de la dot, promesses d'égalité), mais elle exige le respect de la coutume. Baroka, le Bale, apprend par Sadiku que Sidi a refusé sa demande en mariage. Il feint alors l'impuissance devant Sadiku, sachant qu'elle ne pourra pas garder le secret. Sidi, informée, va chez Baroka pour le narguer ; il la séduit par sa parole, son intelligence et sa virilité réelle. Sidi revient triomphante : elle a « dompté le lion » et choisit de l'épouser selon la coutume. Lakunle, déçu, se tourne vers une autre jeune fille en reprenant son discours moderniste. La ruse de Baroka (feindre la faiblesse pour provoquer la curiosité) révèle la supériorité de la sagesse traditionnelle sur l'arrogance moderniste.
\tcblower
\textbf{Analyse :} Ce résumé respecte la chronologie, identifie le pivot dramaturgique (la feinte impuissance), montre la causalité psychologique et thématique, et se termine sur la résolution du conflit central.
\end{exoresolu}

\courssub{Technique du compte rendu}

\begin{definition}[Compte rendu]
Texte bref, objectif et fidèle qui restitue l'essentiel d'une lecture ou d'une activité sans jugement personnel, en respectant la chronologie ou la logique interne du texte source.
\end{definition}

\begin{methode}[Réussir un compte rendu de lecture : la méthode QQOQCP]
\begin{enumerate}
    \item \textbf{Qui ? (source)} Auteur, titre, genre, date, contexte de publication.
    \item \textbf{Qui ? (personnages)} Protagonistes, antagonistes, figures secondaires signifiantes.
    \item \textbf{Où ? Quand ?} Espace (village d'Ilujinle, case de Baroka, école, marché) et temps (une seule journée, du matin au soir, à l'époque contemporaine de l'écriture).
    \item \textbf{Quoi ?} Intrigue principale : nœud, péripéties, dénouement.
    \item \textbf{Comment ?} Procédés marquants : théâtre dans le théâtre, chants, proverbes, structure en trois mouvements (matin, midi, soir).
    \item \textbf{Pourquoi ? (enjeu principal)} Thèse de l'œuvre : la tradition n'est pas figée, elle sait absorber la modernité par la ruse et la vitalité ; la modernité importée sans racine reste stérile.
\end{enumerate}
\end{methode}

\begin{attention}[Compte rendu vs commentaire]
Ne confondez jamais :
\begin{itemize}
    \item \textbf{Compte rendu} = \emph{objectivité, fidélité, reformulation, sélection de l'essentiel}. « L'auteur montre que... », « L'intrigue suit... », « Le personnage incarne... ».
    \item \textbf{Commentaire} = \emph{subjectivité assumée, analyse, interprétation, argumentation}. « Je pense que... », « Cette scène révèle... », « On peut voir ici une critique de... ».
\end{itemize}
Au Probatoire technique, l'exercice demandé est souvent un \textbf{compte rendu} (contrôle de lecture) ou une \textbf{dissertation} (analyse argumentée), rarement un commentaire composé pur.
\end{attention}

\begin{exemplebox}[Extrait de compte rendu modèle (environ 150 mots)]
Wole Soyinka, \emph{Le Lion et la perle} (1959), comédie en trois mouvements (matin, midi, soir), met en scène le conflit entre tradition yoruba et modernité occidentale dans le village nigérian d'Ilujinle. Baroka, le Bale (chef traditionnel), et Lakunle, jeune instituteur formé à l'occidentale, rivalisent pour épouser Sidi, la belle du village. Lakunle refuse de payer la dot, jugeant la coutume archaïque ; Sidi l'exige. Baroka, informé du refus de Sidi, feint l'impuissance devant son épouse Sadiku, certain qu'elle divulguera le secret. Sidi, venue le railler, tombe sous le charme de sa verve et de sa virilité. Elle choisit Baroka. Lakunle, inchangé, poursuit une autre jeune fille. L'œuvre montre que la tradition, loin d'être rigide, use de ruse et de vitalité pour intégrer le nouveau ; la modernité déracinée, dogmatique, échoue par manque d'ancrage. La langue mêle proverbes, chants et danses, faisant du théâtre un espace de négociation culturelle.
\end{exemplebox}

\courssub{Négociation du projet d'études}

La \cle{négociation du projet d'études} est une démarche collective : la classe, guidée par le professeur, choisit les angles d'approfondissement de l'œuvre pour les séances suivantes. C'est un exercice de \textbf{démocratie cognitive} : chaque élève propose, argue, écoute, et le groupe décide.

\begin{center}
\begin{tabularx}{\linewidth}{|l|X|X|}\hline
\rowcolor{popblueL}\textbf{Angle proposé} & \textbf{Question directrice} & \textbf{Activités envisagées} \\ \hline
Thème : tradition vs modernité & La tradition est-elle figée ou dynamique ? & Débat contradictoire, dissertation, saynète \\ \hline
Mise en scène : théâtre dans le théâtre & Comment la danse du voyageur perdu structure-t-elle la pièce ? & Analyse de la métathéâtralité, mise en voix \\ \hline
Personnages : figures du pouvoir & Baroka et Lakunle : deux modèles de masculinité, deux pouvoirs & Portraits croisés, procès simulé \\ \hline
Langue : oralité et écriture & Quel rôle pour les proverbes, les chants, les danses ? & Recueil, traduction, étude stylistique \\ \hline
Femme : Sidi, objet ou sujet ? & Sidi a-t-elle le pouvoir de choisir ? & Table ronde, réécriture du point de vue de Sidi \\ \hline
Contexte : Nigeria post-colonial & En quoi la pièce reflète-t-elle les tensions des années 1960 ? & Exposé documentaire, carte mentale \\ \hline
\end{tabularx}
\end{center}

\begin{methode}[Négocier le projet d'études en classe]
\begin{enumerate}
    \item \textbf{Remue-méninges individuel (5 min)} : chaque élève note 3 angles qui l'intéressent et 1 question brûlante.
    \item \textbf{Mise en commun (10 min)} : au tableau, regroupement des propositions par affinités thématiques.
    \item \textbf{Argumentation (10 min)} : pour chaque piste, un élève défend : « Pourquoi cet angle est-il fécond pour comprendre l'œuvre et progresser en dissertation ? »
    \item \textbf{Vote pondéré (5 min)} : 3 voix par élève à répartir. Les 2 ou 3 angles majoritaires forment le projet.
    \item \textbf{Contractualisation (5 min)} : écriture au tableau du \emph{contrat d'études} : angles retenus, calendrier, productions attendues, critères d'évaluation.
\end{enumerate}
\end{methode}

\courssub{Grille de remédiation : diagnostic des erreurs types}

La remédiation n'est pas une punition ; c'est une \cle{seconde chance d'apprentissage} prévue par la pédagogie par compétences. Elle part d'un diagnostic précis.

\begin{savaistu}[La remédiation dans la pédagogie par compétences]
Dans l'approche par compétences (APC) adoptée par le MINESEC, l'évaluation est \emph{formative} avant d'être certificative. L'erreur n'est pas une faute à sanctionner, mais un \emph{indicateur} du chemin cognitif de l'élève. La remédiation est l'acte pédagogique qui transforme l'erreur en apprentissage : on réexplique, on remodélise, on fait refaire différemment. C'est un droit de l'élève, pas une faveur.
\end{savaistu}

\begin{center}
\begin{popfigure}
\begin{tikzpicture}[
    crit/.style={rectangle, draw=popblue, fill=popblueL, thick, rounded corners=4pt, minimum height=1.1cm, text width=2.8cm, align=center, font=\small\bfseries},
    pts/.style={rectangle, draw=popgreen, fill=popgreenL, thick, rounded corners=4pt, minimum height=1.1cm, text width=1.2cm, align=center, font=\small\bfseries},
    desc/.style={rectangle, draw=poporange, fill=poporangeL, thick, rounded corners=4pt, minimum height=1.1cm, text width=4.5cm, align=left, font=\small},
    header/.style={rectangle, draw=popdark, fill=popdark, thick, rounded corners=4pt, minimum height=0.9cm, text width=2.8cm, align=center, font=\small\bfseries, text=white},
    headerpts/.style={rectangle, draw=popdark, fill=popdark, thick, rounded corners=4pt, minimum height=0.9cm, text width=1.2cm, align=center, font=\small\bfseries, text=white},
    headerdesc/.style={rectangle, draw=popdark, fill=popdark, thick, rounded corners=4pt, minimum height=0.9cm, text width=4.5cm, align=center, font=\small\bfseries, text=white},
]
% En-têtes
\node[header] (h1) at (0,0) {Critère};
\node[headerpts] (h2) at (2.2,0) {Points};
\node[headerdesc] (h3) at (5.4,0) {Descripteurs de performance (niveau attendu 1re Tech)};

% Ligne 1
\node[crit] (c1) at (0,-1.2) {Analyse du sujet};
\node[pts] (p1) at (2.2,-1.2) {4 pts};
\node[desc] (d1) at (5.4,-1.2) {Mots-clés identifiés, sujet reformulé, type de dissertation nommé, tension repérée};

% Ligne 2
\node[crit] (c2) at (0,-2.4) {Problématique};
\node[pts] (p2) at (2.2,-2.4) {4 pts};
\node[desc] (d2) at (5.4,-2.4) {Question ouverte, issue du sujet, porteuse d'enjeu, ni trop large ni trop étroite, guide le plan};

% Ligne 3
\node[crit] (c3) at (0,-3.6) {Cohérence du plan};
\node[pts] (p3) at (2.2,-3.6) {6 pts};
\node[desc] (d3) at (5.4,-3.6) {3 parties équilibrées, 2-3 sous-parties logiques, progression dialectique (thèse/antithèse/synthèse), transitions annoncées};

% Ligne 4
\node[crit] (c4) at (0,-4.8) {Exemples / arguments};
\node[pts] (p4) at (2.2,-4.8) {4 pts};
\node[desc] (d4) at (5.4,-4.8) {Au moins 2 exemples pertinents par sous-partie (texte étudié, culture générale, expérience), explicités, liés à la thèse};

% Ligne 5
\node[crit] (c5) at (0,-6.0) {Langue / expression};
\node[pts] (p5) at (2.2,-6.0) {2 pts};
\node[desc] (d5) at (5.4,-6.0) {Phrases complètes, vocabulaire précis, connecteurs logiques, ponctuation correcte, orthographe soignée};

% Total
\node[crit] (ct) at (0,-7.2) {\textbf{TOTAL}};
\node[pts] (pt) at (2.2,-7.2) {\textbf{20 pts}};
\node[desc] (dt) at (5.4,-7.2) {\textbf{Seuil de maîtrise : 12/20. En dessous : remédiation ciblée sur le(s) critère(s) faible(s).}};
\end{tikzpicture}
\legende{Grille critériée d'évaluation d'un plan détaillé de dissertation (sur 20 points) — adaptée aux attendus du Probatoire technique.}
\end{popfigure}
\end{center}

\begin{aretenir}[Grille de remédiation : diagnostic rapide]
\begin{itemize}
    \item \textbf{Hors-sujet (analyse 0-1/4)} : mots-clés non identifiés, sujet non reformulé, type de dissertation inconnu $\rightarrow$ atelier « Décortiquer un sujet ».
    \item \textbf{Problématique absente ou fermée (problématique 0-1/4)} : question fermée, hors-sujet, binaire, ou simple reformulation $\rightarrow$ atelier « De la tension à la question ».
    \item \textbf{Plan confus, déséquilibré, linéaire (plan 0-3/6)} : parties inégales, pas de progression, répétitions, pas de transitions $\rightarrow$ atelier « Architecturer sa pensée ».
    \item \textbf{Exemples vides ou hors-sujet (exemples 0-1/4)} : citations non analysées, culture générale absente, exemples personnels non généralisables $\rightarrow$ atelier « Argumenter avec le texte ».
    \item \textbf{Langue défaillante (langue 0-1/2)} : phrases nominales, ponctuation aléatoire, vocabulaire pauvre, fautes récurrentes $\rightarrow$ atelier « Écrire en phrases, ponctuer pour lire ».
\end{itemize}
\end{aretenir}

\courssub{Ateliers de remédiation différenciés}

Chaque atelier dure 30 à 45 minutes, en petits groupes (3-4 élèves) ou individuellement, pendant que le reste de la classe avance sur le projet d'études.

\begin{center}
\begin{tabularx}{\linewidth}{|l|X|X|}\hline
\rowcolor{popblueL}\textbf{Atelier} & \textbf{Objectif} & \textbf{Dispositif} \\ \hline
A1 : Décortiquer un sujet & Identifier mots-clés, reformuler, typer, trouver la tension & 5 sujets types $\rightarrow$ fiche « Analyse » commune $\rightarrow$ comparaison et correction collective \\ \hline
A2 : De la tension à la problématique & Transformer une opposition en question ouverte féconde & Jeu de cartes : « Tension » $\rightarrow$ « Questions candidates » $\rightarrow$ choix justifié $\rightarrow$ test : « Cette question guide-t-elle un plan en 3 parties ? » \\ \hline
A3 : Architecturer sa pensée & Construire un plan dialectique équilibré & Plan « squelette » à trous $\rightarrow$ remplissage guidé $\rightarrow$ vérification : 3 parties ? progression ? transitions ? $\rightarrow$ réécriture propre \\ \hline
A4 : Argumenter avec le texte & Sélectionner, citer, expliciter, lier à la thèse & Corpus de 10 citations de \emph{Le Lion et la perle} $\rightarrow$ classement par thèse possible $\rightarrow$ rédaction de 2 paragraphes types (thèse + exemple + commentaire) \\ \hline
A5 : Écrire en phrases, ponctuer pour lire & Maîtriser la phrase complexe, la ponctuation forte et faible & Dictée raisonnée de 3 phrases modèles $\rightarrow$ analyse syntaxique $\rightarrow$ réécriture de 5 phrases « télégraphiques » en phrases complètes ponctuées \\ \hline
\end{tabularx}
\end{center}

\begin{experience}[Atelier type A3 : Architecturer sa pensée — déroulement pas à pas]
\textbf{Matériel} : sujet « La tradition est-elle un obstacle au progrès ? », plan squelette (A3), feutres de 3 couleurs (thèse / antithèse / synthèse), affichage des connecteurs logiques.

\textbf{Protocole} :
\begin{enumerate}
    \item \textbf{Modélisation (5 min)} : le professeur construit au tableau le plan du sujet avec la classe (thèse : tradition frein ; antithèse : tradition ressource ; synthèse : tradition vivante = innovation ancrée).
    \item \textbf{Manipulation (15 min)} : en binôme, les élèves reçoivent 12 cartes « idées/arguments » mélangées. Ils les trient en 3 colonnes, les ordonnent à l'intérieur de chaque colonne, puis ajoutent 2 connecteurs par transition.
    \item \textbf{Verbalisation (10 min)} : chaque binôme présente son plan à l'oral : « Partie 1 : ... parce que ... ; transition : cependant / en revanche ; partie 2 : ... » La classe valide ou conteste la cohérence.
    \item \textbf{Transfert (10 min)} : individuellement, chaque élève rédige le plan détaillé (introduction, 3 parties avec 2-3 sous-parties, conclusion) pour un nouveau sujet proche : « Faut-il tourner le dos aux coutumes pour se moderniser ? »
\end{enumerate}

\textbf{Observations attendues} : les élèves qui peinaient sur la progression dialectique (souvent un plan « catalogue » ou « chronologique ») produisent un plan à tension interne. Le codage couleur rend visible la structure argumentative.

\textbf{Interprétation} : la remédiation par manipulation (tri, ordonnancement, verbalisation) externalise la pensée logique avant l'écrit. Elle réduit la charge cognitive de la rédaction finale.
\end{experience}

\courssub{Bilan de l'unité : auto-évaluation et contrat de progrès}

\begin{center}
\begin{popfigure}
\begin{tikzpicture}[
    node/.style={circle, draw=popdark, thick, minimum size=1.8cm, font=\small\bfseries, align=center},
    arrow/.style={->, >=Stealth, thick, popdark},
    label/.style={font=\small, align=center, text width=2.2cm},
    plus/.style={rectangle, draw=popgreen, fill=popgreenL, rounded corners=3pt, font=\tiny, minimum width=1.5cm, minimum height=0.6cm},
    moins/.style={rectangle, draw=poporange, fill=poporangeL, rounded corners=3pt, font=\tiny, minimum width=1.5cm, minimum height=0.6cm},
]
% Nœuds du cycle
\node[node, fill=popblueL, align=center] (prod) at (0,0){PRODUCTION\\Dissertation\\Plan détaillé};
\node[node, fill=poporangeL, align=center] (eval) at (5,0){ÉVALUATION\\Grille 20 pts\\Auto / pair / prof};
\node[node, fill=poppurpleL, align=center] (remed) at (5,-4){RÉMÉDIATION\\Ateliers ciblés\\A1--A5};
\node[node, fill=popgreenL, align=center] (newprod) at (0,-4){NOUVELLE\\PRODUCTION\\Sujet voisin\\Transfert};

% Flèches principales
\draw[arrow] (prod.east) -- node[above, label, align=center]{Grille\\critériée} (eval.west);
\draw[arrow] (eval.south) -- node[right, label, align=center]{Diagnostic\\critère par critère} (remed.north);
\draw[arrow] (remed.west) -- node[below, label, align=center]{Compétence\\retravaillée} (newprod.east);
\draw[arrow] (newprod.north) -- node[left, label, align=center]{Progression\\mesurée} (prod.south);

% Boucles de différenciation
\draw[arrow, popgreen, dashed, bend left=20] (eval.north east) to node[plus, pos=0.5, above right, align=center]{Maîtrisé \\$\geq$ 12/20} (prod.north east);
\draw[arrow, poporange, dashed, bend right=20] (eval.south east) to node[moins, pos=0.5, below right, align=center]{Fragile \\$<$ 12/20} (remed.north east);

% Ancrage
\node[label] at (2.5, 2.2) {\textbf{CYCLE D'APPRENTISSAGE PAR COMPÉTENCES}};
\node[label] at (2.5, -5.8) {Chaque tour de boucle = une compétence consolidée. La remédiation n'est pas une sortie de route, c'est le virage nécessaire.};
\end{tikzpicture}
\legende{Schéma du cycle d'apprentissage : production $\rightarrow$ évaluation $\rightarrow$ remédiation $\rightarrow$ nouvelle production. Les boucles vertes (maîtrise) et oranges (fragilité) montrent la différenciation.}
\end{popfigure}
\end{center}

\begin{methode}[Bilan personnel : auto-évaluation et contrat de progrès]
Chaque élève remplit cette fiche en fin d'unité (15 min, en silence, avec sincérité) :

\begin{enumerate}
    \item \textbf{Compétences visées} (entourer la mention qui convient : \emph{acquis} / \emph{en cours} / \emph{fragile}) :
    \begin{itemize}
        \item Analyser un sujet de dissertation (mots-clés, reformulation, type, tension) : acquis / en cours / fragile
        \item Formuler une problématique féconde : acquis / en cours / fragile
        \item Construire un plan détaillé dialectique en 3 parties équilibrées : acquis / en cours / fragile
        \item Sélectionner et exploiter des exemples pertinents (texte + culture générale) : acquis / en cours / fragile
        \item Rédiger en phrases correctes, ponctuées, connectées : acquis / en cours / fragile
        \item Rendre un compte rendu objectif et structuré d'une œuvre lue : acquis / en cours / fragile
        \item Négocier un projet d'études collectif (écoute, argumentation, engagement) : acquis / en cours / fragile
    \end{itemize}
    \item \textbf{Mon point fort dans cette unité} : \underline{\hspace{8cm}}
    \item \textbf{Mon point de vigilance (le critère le plus bas sur la grille)} : \underline{\hspace{8cm}}
    \item \textbf{Mon action concrète pour progresser d'ici la prochaine dissertation} (choisir 1 atelier A1-A5, 1 exercice précis, 1 relecture ciblée) : \underline{\hspace{8cm}}
    \item \textbf{Date de réévaluation visée} : \underline{\hspace{3cm}} \quad Signature : \underline{\hspace{4cm}}
\end{enumerate}
\end{methode}

\begin{exemplebox}[Contrat de progrès type — Élève K.]
\begin{itemize}
    \item \textbf{Point fort} : « Je sais maintenant trouver la tension dans un sujet et écrire une problématique ouverte. » (Problématique : 4/4)
    \item \textbf{Point de vigilance} : « Mon plan reste déséquilibré : partie 1 trop longue, partie 3 inexistante. » (Plan : 2/6)
    \item \textbf{Action} : « Atelier A3 (Architecturer) + refaire le plan détaillé du sujet "La tradition est-elle un obstacle au progrès ?" en respectant le gabarit 3 parties / 2 sous-parties / transitions explicites. »
    \item \textbf{Réévaluation} : « Vendredi prochain, sujet d'entraînement noté sur grille. Objectif : plan $\geq$ 4/6. »
\end{itemize}
\end{exemplebox}

\courssub{Synthèse de l'unité : de la réception à la production experte}

\begin{propriete}[Parcours complet de l'unité « Dissertation : produire un plan détaillé »]
L'unité a suivi une progression spiralée, du texte support (\emph{Le Lion et la perle}) aux outils linguistiques, jusqu'à la production autonome évaluée et remédiée :
\begin{enumerate}
    \item \textbf{Réception} : paratextes $\rightarrow$ lecture ouvroir $\rightarrow$ contrôle de lecture $\rightarrow$ compte rendu $\rightarrow$ négociation du projet.
    \item \textbf{Langue} : situations de communication $\rightarrow$ énonciation (marques explicites) $\rightarrow$ ponctuation $\rightarrow$ texte descriptif $\rightarrow$ dénotation/connotation.
    \item \textbf{Production} : analyse du sujet $\rightarrow$ problématique $\rightarrow$ recherche des idées $\rightarrow$ plan détaillé $\rightarrow$ confrontation et amélioration.
    \item \textbf{Métacognition} : grille critériée $\rightarrow$ auto-évaluation $\rightarrow$ remédiation différenciée $\rightarrow$ contrat de progrès.
\end{enumerate}
Chaque étape a nourri la suivante : la lecture de Soyinka a fourni le réservoir d'exemples ; les outils linguistiques ont armé l'expression ; la grille a objectivé l'exigence ; la remédiation a bouclé la boucle.
\end{propriete}

\begin{attention}[Pièges de fin d'unité à éviter]
\begin{itemize}
    \item \textbf{Zapper l'auto-évaluation} : « Je n'ai pas le temps. » $\rightarrow$ C'est le moment le plus rentable : 15 minutes ici économisent 2 heures de répétition d'erreurs.
    \item \textbf{Remédiation « punitive »} : « Tu as eu 8, refais. » $\rightarrow$ Sans diagnostic critérié ni atelier ciblé, l'élève refait la même erreur.
    \item \textbf{Projet d'études « subi »} : le professeur impose les angles. $\rightarrow$ Perte d'engagement, surface au lieu de profondeur.
    \item \textbf{Compte rendu « commentaire »} : l'élève glisse son avis personnel. $\rightarrow$ Perte de points au Probatoire (consigne non respectée).
\end{itemize}
\end{attention}

\begin{exoresolu}[Bilan intégratif : du plan à la dissertation — transfert]
\textbf{Énoncé :} À partir du plan détaillé élaboré dans la section 4 (sujet : « La tradition est-elle un obstacle au progrès ? »), rédiger l'introduction complète et le premier paragraphe de la partie I (thèse : la tradition frein).

\textbf{Plan rappelé :}
\begin{itemize}
    \item \textbf{Introduction} : accroche (proverbe) $\rightarrow$ présentation du sujet $\rightarrow$ problématique $\rightarrow$ annonce du plan.
    \item \textbf{I. La tradition comme frein au changement} (thèse)
        \begin{itemize}
            \item A. Poids des coutumes figées (ex. : dot, polygamie, gérontocratie)
            \item B. Résistance aux innovations techniques et sociales (ex. : école, santé, droits des femmes)
            \item C. Risque de repli identitaire et de refus du dialogue interculturel
        \end{itemize}
    \item \textbf{II. La tradition comme ressource pour un progrès ancré} (antithèse) ...
    \item \textbf{III. Vers une modernité négociée : la tradition vivante} (synthèse) ...
    \item \textbf{Conclusion} : bilan $\rightarrow$ ouverture (autres sociétés, rôle de l'éducation).
\end{itemize}

\textbf{Production attendue (extrait) :}

\textbf{Introduction.} « Quand le vent du changement souffle, les uns construisent des murs, les autres des moulins. » Ce proverbe résonne avec la pièce de Wole Soyinka, \emph{Le Lion et la perle}. Dans le village d'Ilujinle, la confrontation entre Baroka, gardien de la tradition, et Lakunle, héraut de la modernité, pour conquérir Sidi, met en scène le dilemme de toute société en mutation : faut-il rompre avec son passé pour avancer ? La tradition, ensemble de coutumes, de valeurs et de savoirs transmis, est-elle structurellement un obstacle au progrès, entendu comme amélioration matérielle, sociale et morale ? Ou peut-elle en être le ferment ? Nous montrerons d'abord que la tradition peut freiner le changement par son inertie et ses interdits (I), avant de démontrer qu'elle offre aussi des ressources symboliques et sociales pour un progrès endogène (II), pour aboutir à l'idée d'une modernité négociée, où la tradition vivante innove sans s'arracher de ses racines (III).

\textbf{Premier paragraphe (I.A).} La première forme d'obstacle réside dans le poids des coutumes figées, érigées en normes intangibles. Dans \emph{Le Lion et la perle}, la dot en est l'exemple central : Lakunle la refuse au nom de l'égalité moderne, mais Sidi l'exige comme garantie de sa valeur et de son honneur. Cette coutume, conçue pour sceller l'alliance entre lignages, devient un blocage quand elle est brandie comme condition absolue, sans adaptation aux nouvelles aspirations. De même, la gérontocratie — le pouvoir des aînés — empêche la parole jeune : Sadiku, malgré son rang d'épouse aînée, est instrumentalisée par Baroka ; les jeunes filles n'ont pas voix au chapitre. Hors du texte, on observe dans de nombreuses sociétés traditionnelles que la polygamie, le lévirat ou certains interdits peuvent entraver l'accès des femmes à l'éducation, à la santé, à l'autonomie économique. La tradition, lorsqu'elle se fige en dogme, devient alors un frein structurel au progrès social.
\tcblower
\textbf{Analyse didactique} : l'introduction suit le canevas (accroche contextualisée, sujet posé, problématique, plan annoncé). Le paragraphe I.A respecte la structure « thèse de la sous-partie + exemple textuel analysé + exemple hors-texte généralisé + mini-conclusion ». La ponctuation (deux-points, point-virgule, tirets) structure la pensée. Le connecteur « De même » assure la cohésion interne. C'est le niveau attendu en fin de Première technique.
\end{exoresolu}

\begin{aretenir}[Essentiel à retenir pour le Probatoire]
\begin{itemize}
    \item \textbf{Compte rendu} = objectivité, fidélité, méthode QQOQCP, enjeu principal. Pas d'« à mon avis ».
    \item \textbf{Remédiation} = droit, pas sanction. Diagnostic par grille $\rightarrow$ atelier ciblé $\rightarrow$ nouvelle production.
    \item \textbf{Projet d'études} = négociation collective, contrat écrit, angles choisis = motivation + profondeur.
    \item \textbf{Auto-évaluation} = grille critériée + contrat de progrès (1 action concrète + 1 date).
    \item \textbf{Cycle} : produire $\rightarrow$ évaluer $\rightarrow$ remédier $\rightarrow$ reproduire (mieux). C'est le moteur de la compétence.
\end{itemize}
\end{aretenir}

Cette unité vous a fait parcourir le chemin complet : de la lecture d'une œuvre africaine majeure à la production d'un plan de dissertation évalué, remédié, transféré. Les outils linguistiques (ponctuation, énonciation, description, dénotation/connotation) ne sont pas des exercices de grammaire déconnectés : ce sont les \textbf{armes rhétoriques} de votre pensée argumentative. Le projet d'études sur \emph{Le Lion et la perle} vous appartient désormais : c'est vous qui en tracez les pistes, vous qui en portez la parole. Le Probatoire n'est pas une épreuve de mémoire, mais de \textbf{compétence en action}. Vous avez la méthode, la grille, le contrat. À vous de jouer.

\newpage

\coursec{Activité d'intégration}

\coursec{Leçon 01 : Dissertation -- Produire le plan détaillé d'une dissertation}

\courssub{Objectifs de l'activité}
\begin{itemize}
    \item Mobiliser de manière autonome et articulée l'ensemble des savoir-faire de l'unité : analyse de sujet, formulation de la problématique, recherche d'idées (brainstorming, carte mentale), choix et justification du type de plan, rédaction d'un plan détaillé complet (introduction, développement structuré, conclusion).
    \item Exploiter des ressources culturelles (étude de \emph{Le Lion et la perle}), contextuelles (actualité camerounaise : artisanat, numérique, sport) et chiffrées (données d'insertion professionnelle) pour nourrir l'argumentation.
    \item Développer l'esprit critique et la compétence évaluative par la co-évaluation croisée à l'aide d'une grille critériée, suivie d'une remédiation immédiate et d'un compte rendu oral structuré.
\end{itemize}

\courssub{Situation}
\label{sit:integration}
Dans le cadre de la préparation à l'épreuve de dissertation du Probatoire (série technique), votre établissement organise un \emph{Concours de l'Excellence Argumentative} inter-classes. Le jury, composé de professeurs de lettres, de conseillers d'orientation et de professionnels du secteur privé (entrepreneurs, artisans, responsables RH), a retenu le sujet suivant :

\begin{center}
\textbf{« L'école est-elle la seule voie de la réussite sociale au Cameroun ? »}
\end{center}

\vspace{0.5em}
\noindent
\textbf{Contexte camerounais et données fournies aux candidats :}
\begin{itemize}
    \item \textbf{Données statistiques (source : INS, Rapport 2023 sur l'insertion des jeunes) :}
    \begin{itemize}
        \item Taux d'insertion professionnelle à 3 ans : Filières techniques professionnelles (BT, BTS) : \textbf{68\%} ; Filières générales (Baccalauréat général) : \textbf{42\%} ; Non diplômés (apprentissage informel, auto-formation) : \textbf{35\%}.
        \item Taux de chômage des jeunes diplômés de l'enseignement supérieur (Licence/Master) : \textbf{23\%} (dont une forte proportion en attente d'emploi public).
    \end{itemize}
    \item \textbf{Référence littéraire obligatoire :} \emph{Le Lion et la perle} de Wole Soyinka. Le conflit entre Baroka (le Bale, tradition, pouvoir coutumier, réussite sociale sans école occidentale) et Lakunle (instituteur, modernité scolaire, échec social) éclaire la tension entre \emph{instruction formelle} et \emph{savoir-faire endogène}.
    \item \textbf{Actualité locale (exemples à disposition) :}
    \begin{itemize}
        \item \textit{Artisanat d'excellence :} Les bronziers de Foumban, les tisseurs de Baham, les sculpteurs de la région du Nord (insertion via coopératives, export, reconnaissance UNESCO).
        \item \textit{Numérique \& Tech :} Jeunes développeurs auto-formés (Yaoundé, Douala, Bafoussam) employés par des startups (KmerTech, GiftedMom) ou freelances internationaux, sans diplôme universitaire initial.
        \item \textit{Sport :} Parcours de Francis Ngannou (de la carrière de sable au titre mondial MMA) ou de jeunes footballeurs formés en académies locales (Kadji Sports, APEJES) signant en Europe.
    \end{itemize}
\end{itemize}

\vspace{0.5em}
\noindent
Votre classe est divisée en îlots de 4 à 5 élèves. Chaque îlot constitue une « équipe de rédaction » qui doit produire, en temps limité, un \textbf{plan détaillé prêt à être développé}, évalué par les pairs selon la grille officielle du Probatoire (cohérence, pertinence, exploitation des connaissances, qualité de la langue).

\courssub{Consignes}
\label{consignes:integration}
\emph{Travail en îlots -- Durée totale : 2h30 (répartie sur 2 séances).}

\begin{enumerate}
    \item \textbf{Analyse lexicale et formulation de la problématique (30 min).}
    \begin{enumerate}
        \item Identifiez les \cle{mots-clés} du sujet (au moins 4) et donnez pour chacun sa \cle{dénotation} et une \cle{connotation} pertinente dans le contexte camerounais.
        \item Reformulez le sujet en question simple.
        \item Proposez \textbf{trois problématiques candidates} (de type \emph{oui/mais}, \emph{dans quelle mesure}, \emph{comment}) et sélectionnez la plus féconde en justifiant votre choix (ouverture, tension, faisabilité argumentative).
    \end{enumerate}

    \item \textbf{Brainstorming collectif et carte mentale (15 min).}
    \begin{enumerate}
        \item Sur une feuille A3 (ou outil numérique), construisez une \cle{carte mentale} centrée sur la problématique retenue. Branches obligatoires : \emph{Arguments POUR l'école (thèse)}, \emph{Arguments CONTRE / Voies alternatives (antithèse)}, \emph{Synthèse / Articulation}.
        \item Mobilisez explicitement : \emph{Le Lion et la perle} (2 références précises minimum), les \emph{données chiffrées} (2 citations minimum), l'\emph{actualité} (2 exemples distincts minimum).
    \end{enumerate}

    \item \textbf{Choix et justification du type de plan (15 min).}
    \begin{enumerate}
        \item Parmi les types de plan vus en cours (dialectique, thématique, analytique/problématique, comparatif), choisissez celui qui structure le mieux votre démonstration.
        \item Rédigez une \textbf{justification argumentée (5 à 7 lignes)} : pourquoi ce plan et pas un autre ? En quoi répond-il à la problématique ?
    \end{enumerate}

    \item \textbf{Rédaction du plan détaillé complet (1h).}
    Produisez une fiche unique (format A4 recto) contenant :
    \begin{itemize}
        \item \textbf{Introduction rédigée intégralement} (amorce, présentation du sujet, définition des termes, problématique, annonce du plan).
        \item \textbf{Développement :} 2 ou 3 parties principales (I, II, [III]), chacune divisée en 2 sous-parties (A, B). Pour chaque sous-partie : \cle{idée directrice} (phrase nominale ou verbale), \cle{argument} (1 phrase), \cle{exemple/preuve} (référence précise : texte, chiffre, fait d'actualité), \cle{transition} vers la sous-partie suivante ou la partie suivante.
        \item \textbf{Conclusion rédigée intégralement} (bilan, réponse à la problématique, ouverture).
    \end{itemize}

    \item \textbf{Évaluation croisée, remédiation et compte rendu oral (30 min).}
    \begin{enumerate}
        \item Échangez les productions avec un autre îlot. Évaluez avec la \textbf{grille critériée} (voir Ressources). Notez les forces (✓) et faiblesses (✗) sur la marge.
        \item Récupérez votre production. En 10 min, corrigez/améliorez les points faibles identifiés (encre verte).
        \item \textbf{Compte rendu oral (2 min par groupe) :} Un porte-parole présente : la problématique choisie, le type de plan, l'argument « choc » (le plus original), la difficulté rencontrée et la correction apportée.
    \end{enumerate}
\end{enumerate}

\courssub{Ressources à mobiliser}
\label{ressources:integration}

\begin{methode}[Analyse de sujet -- Étapes obligatoires]
\begin{enumerate}
    \item \textbf{Cerner le sujet :} Souligner les mots-clés ; définir (dénotation/connotation) ; reformuler.
    \item \textbf{Problématiser :} Transformer le sujet en question problématique (tension, débat). Éviter les questions fermées (Oui/Non) ou trop vastes.
    \item \textbf{Tester la problématique :} Vérifier qu'elle appelle une démonstration en plusieurs temps (thèse/antithèse/synthèse ou causes/conséquences/solutions).
\end{enumerate}
\end{methode}

\begin{methode}[Types de plan de dissertation -- Choix stratégique]
\begin{itemize}
    \item \textbf{Plan dialectique (Thèse / Antithèse / Synthèse) :} Idéal pour un sujet « Est-ce que... ? » ou « Dans quelle mesure... ? » portant sur une opposition nette (École vs Hors-école).
    \item \textbf{Plan thématique (Aspects / Dimensions) :} Si le sujet invite à explorer des facettes (économique, culturel, civique) sans opposition binaire.
    \item \textbf{Plan analytique / Problématique (Causes / Conséquences / Solutions ou Constat / Analyse / Perspectives) :} Pour une question « Comment ? » ou « Pourquoi ? ».
    \item \textbf{Plan comparatif :} Pour comparer deux objets/modèles (Ex: Formation initiale vs Formation continue).
\end{itemize}
\emph{Règle d'or :} Le plan doit être la \textbf{réponse visuelle} à la problématique.
\end{methode}

\begin{propriete}[Structure du plan détaillé -- Exigences Probatoire]
\begin{itemize}
    \item \textbf{Introduction :} Amorce (accroche) $\rightarrow$ Présentation du sujet $\rightarrow$ Définitions $\rightarrow$ Problématique $\rightarrow$ Annonce du plan (mots de liaison : \emph{Dans un premier temps... Ensuite... Enfin...}).
    \item \textbf{Développement :} Parties (I, II, III) = grandes étapes de la démonstration. Sous-parties (A, B) = arguments distincts. \textbf{Chaque sous-partie = Idée directrice + Argument + Exemple/Preuve + Transition.}
    \item \textbf{Conclusion :} Bilan synthétique (reprise des grandes idées) $\rightarrow$ Réponse directe à la problématique $\rightarrow$ Ouverture (perspective, autre sujet, citation).
    \item \textbf{Transitions :} Obligatoires entre I et II, II et III. Formulation : \emph{« Si [idée I] est vrai, toutefois [limite]... d'où la nécessité d'examiner [idée II] »}.
\end{itemize}
\end{propriete}

\begin{center}
\begin{tabularx}{\linewidth}{|l|X|}\hline
\rowcolor{popblueL}
\textbf{Critère (Grille Probatoire)} & \textbf{Indicateurs de réussite (Niveau « Maîtrise »)} \\ \hline
\cle{Cohérence / Structure} & Plan logique, parties équilibrées, transitions fluides, adéquation plan/problématique. \\ \hline
\cle{Pertinence / Contenu} & Arguments variés, exemples précis (texte, chiffres, actualité), pas de hors-sujet, nuance (pas de caricature). \\ \hline
\cle{Exploitation des connaissances} & \emph{Le Lion et la perle} utilisé comme levier d'argumentation (pas simple résumé), données chiffrées citées correctement. \\ \hline
\cle{Qualité de la langue} & Syntaxe correcte, vocabulaire précis (register formel), connecteurs logiques variés, ponctuation rigoureuse. \\ \hline
\cle{Présentation} & Lisibilité, respect du formalisme (chiffres romains, majuscules, indentation). \\ \hline
\end{tabularx}
\end{center}

\courssub{Démarche et corrigé commenté}
\label{demarche:integration}

\begin{exoresolu}[Production type attendue -- Sujet : « L'école est-elle la seule voie de la réussite sociale au Cameroun ? »]
\textbf{Étape 1 -- Analyse lexicale et Problématique}

\begin{itemize}
    \item \textbf{Mots-clés :}
    \begin{itemize}
        \item \cle{École} : Dénotation = Institution d'enseignement formel, cursus certifié (primaire, secondaire, supérieur). Connotation = Légitimité institutionnelle, ascenseur social « officiel », mais aussi rigidité, déconnexion du terrain, chômage des diplômés.
        \item \cle{Seule voie} : Dénotation = Unique chemin, exclusivité. Connotation = Exclusion, dogmatisme, déni des autres formes de mérite.
        \item \cle{Réussite sociale} : Dénotation = Ascension statutaire, reconnaissance, insertion professionnelle, revenus stables. Connotation = Subjective (argent, honneur, utilité collective, épanouissement) ; au Cameroun : « faire quelque chose de soi », « bâtir », « soutenir sa famille ».
        \item \cle{Cameroun} : Contexte = Jeunesse nombreuse, économie informelle dominante (90\% emplois), valorisation culturelle du diplôme (fonction publique) vs réalités du marché.
    \end{itemize}
    \item \textbf{Problématiques candidates :}
    \begin{enumerate}
        \item L'école garantit-elle la réussite sociale au Cameroun ? (Trop binaire, fermé).
        \item Quelles sont les voies de la réussite sociale au Cameroun ? (Trop descriptif, pas de tension).
        \item \textbf{\emph{Dans quelle mesure l'institution scolaire, bien que voie royale légitime, suffit-elle à elle seule à assurer la réussite sociale au Cameroun, au regard des parcours alternatifs porteurs ?}} (Retenue : tension dialectique, ouverture aux alternatives, ancrage contextuel).
    \end{enumerate}
\end{itemize}

\textbf{Étape 2 -- Brainstorming (Extrait carte mentale)}

\begin{popfigure}
\begin{tikzpicture}[mindmap, grow cyclic, every node/.style=concept, concept color=popblue!60,
    level 1/.append style={level distance=4.5cm, sibling angle=120},
    level 2/.append style={level distance=3cm, sibling angle=60}]
\node[concept color=popdark] {Problématique}
    child[concept color=popgreen] {node[concept] {I. École : voie royale légitime (Thèse)}
        child {node[concept, scale=0.7] {A. Certification \& Accès emploi formel (Stats: 68\% tech)}}
        child {node[concept, scale=0.7] {B. Citoyenneté, culture commune, mobilité (Lakunle : idéal modernité)}}
    }
    child[concept color=poporange] {node[concept] {II. Limites \& Voies alternatives (Antithèse)}
        child {node[concept, scale=0.7] {A. Chômage diplômés (23\% sup) ; inadéquation formation/emploi}}
        child {node[concept, scale=0.7] {B. Réussites hors école : Artisanat (Foumban), Numérique (auto-formés), Sport (Ngannou) ; Baroka : pouvoir coutumier}}
    }
    child[concept color=poppurple] {node[concept] {III. Complémentarité \& Redéfinition (Synthèse)}
        child {node[concept, scale=0.7] {A. École comme socle (compétences de base) + Spécialisation pro (BTS, apprentissage)}}
        child {node[concept, scale=0.7] {B. Reconnaissance officielle des acquis (VAE, certification métiers) ; École « hors les murs »}}
    };
\end{tikzpicture}
\legende{Carte mentale structurée autour de la problématique retenue : Thèse (vert), Antithèse (orange), Synthèse (violet).}
\end{popfigure}

\textbf{Étape 3 -- Choix du plan : Plan Dialectique (Thèse / Antithèse / Synthèse)}
\textit{Justification :} Le sujet pose une alternative exclusive (« seule voie »). Le plan dialectique permet d'abord de reconnaître la force de la thèse (l'école reste le cadre légitime majeur), puis d'en montrer les limites factuelles (chômage, inadéquation) et la vitalité des alternatives (antithèse), pour enfin proposer une articulation dynamique : l'école nécessaire mais non suffisante, devant s'ouvrir aux compétences réelles (synthèse). Un plan thématique noierait la tension ; un plan analytique ne rendrait pas justice au débat de société.

\textbf{Étape 4 -- Plan Détaillé Complet}

\noindent\textbf{INTRODUCTION (Rédigée)}
« "L'école est la clé de l'avenir", proclame l'adage. Pourtant, au carrefour de la gare de Nsam à Yaoundé, des ingénieurs diplômés conduisent des motos-taxis pendant que, dans l'atelier mitoyen, un soudeur autodidacte exporte ses œuvres en Europe. Au Cameroun, pays où 75\% de la population a moins de 35 ans (INS 2023), la question de la réussite sociale brûle les lèvres. \textbf{Sujet :} L'école, institution certifiante par excellence, est-elle l'unique chemin vers cette réussite ? \textbf{Définitions :} L'école = système éducatif formel sanctionné par des diplômes. Réussite sociale = insertion durable, reconnaissance, autonomie financière. \textbf{Problématique :} \emph{Dans quelle mesure l'institution scolaire, bien que voie royale légitime, suffit-elle à elle seule à assurer la réussite sociale au Cameroun, au regard des parcours alternatifs porteurs ?} \textbf{Annonce du plan :} Nous montrerons d'abord que l'école demeure la voie structurante et légitime (I), avant d'analyser ses limites face à la vitalité des parcours hors cadre formel (II), pour enfin proposer une redéfinition de la réussite par la complémentarité des savoirs (III). »

\noindent\textbf{DÉVELOPPEMENT}

\noindent\textbf{I. L'école : une voie royale structurante et légitime (Thèse)}

\noindent\textbf{A. La certification scolaire, passeport pour l'emploi formel et la protection sociale.}
\begin{itemize}
    \item \textit{Idée directrice :} Le diplôme reste le signal principal pour l'employeur formel et l'accès à la fonction publique.
    \item \textit{Argument :} Les filières techniques (BT, BTS) affichent un taux d'insertion à 3 ans de 68\% (INS 2023), bien supérieur au général (42\%) et au non-diplômé (35\%).
    \item \textit{Exemple :} Concours administratifs (ENAM, ENIEG) : diplôme exigé. Statut de fonctionnaire = sécurité, retraite, prestige social.
    \item \textit{Transition :} Cependant, ce sésame ne garantit plus l'emploi pour tous, ni l'adéquation aux besoins réels de l'économie.
\end{itemize}

\noindent\textbf{B. L'école comme creuset de la citoyenneté, de la culture commune et de la mobilité intellectuelle.}
\begin{itemize}
    \item \textit{Idée directrice :} Au-delà du diplôme, l'école forge le citoyen et les compétences transversales (langues, sciences, pensée critique).
    \item \textit{Argument :} Elle offre un capital culturel partagé, condition de la cohésion nationale dans un pays à 250 ethnies.
    \item \textit{Exemple :} \emph{Le Lion et la perle} : Lakunle, malgré son ridicule, incarne l'idéal modernisateur : il veut transformer Ilujinle par l'instruction, la langue anglaise, la monogamie. L'école porte un projet de société.
    \item \textit{Transition :} Mais cet idéal se heurte à la réalité du marché : l'école forme-t-elle encore aux besoins du pays ?
\end{itemize}

\noindent\textbf{II. Les limites de l'exclusivité scolaire : la vitalité des parcours alternatifs (Antithèse)}

\noindent\textbf{A. Le décrochage de l'école : chômage des diplômés, inadéquation formation/emploi, coût d'opportunité.}
\begin{itemize}
    \item \textit{Idée directrice :} L'inflation diplômante dévalue le titre ; l'université produit pour un marché public saturé.
    \item \textit{Argument :} 23\% de chômage chez les diplômés du supérieur (INS 2023) ; beaucoup acceptent des emplois infra-qualifiés (Bac+5 vendeur en boutique).
    \item \textit{Exemple :} Lakunle (\emph{Le Lion et la perle}) : instruit, il perd Sidi et le respect du village. Son savoir livresque est inopérant face au pragmatisme de Baroka. L'école peut déconnecter du réel.
    \item \textit{Transition :} Face à cette impasse, la société camerounaise valorise d'autres formes de mérite, concrètes et rémunératrices.
\end{itemize}

\noindent\textbf{B. La réussite par les savoir-faire endogènes, le numérique et le talent brut : des preuves par l'exemple.}
\begin{itemize}
    \item \textit{Idée directrice :} L'artisanat d'excellence, la tech auto-apprise et le sport créent de la richesse et du statut sans diplôme initial.
    \item \textit{Argument :} Ces voies reposent sur la compétence observable, l'entrepreneuriat, la résilience -- valeurs techniques chères à votre formation.
    \item \textit{Exemples :}
    \begin{itemize}
        \item \textit{Artisanat :} Bronziers de Foumban (label UNESCO), tisseurs de Baham : revenus élevés, transmission par apprentissage (maître/compagnon), fierté identitaire. Baroka (\emph{Le Lion et la perle}) : pouvoir coutumier, richesse, épouses, sans école occidentale.
        \item \textit{Numérique :} Développeurs formés sur YouTube/GitHub, recrutés par GiftedMom, KmerTech ou en freelance international (Upwork). Compétence \emph{coding} > Diplôme.
        \item \textit{Sport :} Francis Ngannou (carrière de sable $\rightarrow$ champion UFC) ; académies de foot (Kadji, APEJES) $\rightarrow$ contrats pro Europe. Corps et mental comme capital.
    \end{itemize}
    \item \textit{Transition :} Opposer école et hors-école est stérile. L'avenir appartient à leur hybridation.
\end{itemize}

\noindent\textbf{III. Vers une complémentarité féconde : redéfinir la réussite et l'école (Synthèse)}

\noindent\textbf{A. L'école comme socle commun indispensable + spécialisation professionnalisante (voie technique/apprentissage).}
\begin{itemize}
    \item \textit{Idée directrice :} L'école ne doit pas disparaître mais muter : tronc commun solide (français, maths, numérique, soft skills) + orientation précoce vers filières techniques/pro (BTS, apprentissage dual).
    \item \textit{Argument :} Les 68\% d'insertion des filières techniques prouvent l'efficacité du lien école-entreprise.
    \item \textit{Exemple :} Modèle allemand/dual adapté au Cameroun (partenariats GIZ, entreprises locales). L'élève technicien (vous) : théorie en classe + pratique en stage = employabilité.
    \item \textit{Transition :} Il faut aussi reconnaître officiellement les compétences acquises hors les murs.
\end{itemize}

\noindent\textbf{B. La reconnaissance des acquis (VAE) et l'école « tout au long de la vie » : décloisonner les parcours.}
\begin{itemize}
    \item \textit{Idée directrice :} Instaurer la Validation des Acquis de l'Expérience (VAE) pour certifier l'artisan, le codeur, l'athlète, leur permettant d'accéder à des grades, la fonction publique, la formation continue.
    \item \textit{Argument :} Transformer la « seule voie » en « multiples entrées, même certification ».
    \item \textit{Exemple :} Projet de loi sur la VAE au Cameroun (MINESUP/MINEFOP) ; passerelles entre formation informelle et certifications nationales (CQP -- Certificats de Qualification Professionnelle).
    \item \textit{Transition (vers conclusion) :} La réussite sociale n'est pas un lieu (l'école) mais un processus : transformer un potentiel en compétence reconnue.
\end{itemize}

\noindent\textbf{CONCLUSION (Rédigée)}
« Au terme de cette analyse, il apparaît que l'école camerounaise, si elle reste la \textbf{voie structurante majeure} (certification, citoyenneté, insertion technique prouvée à 68\%), ne saurait être la \textbf{seule voie}. Le chômage des diplômés (23\% au supérieur) et l'éclatante réussite d'artisans de Foumban, de codeurs autodidactes ou de champions comme Ngannou démontrent que la compétence, l'audace et l'utilité sociale s'acquièrent aussi hors les murs. \textbf{Réponse :} L'école est une condition \emph{presque} nécessaire pour l'emploi formel, mais non \emph{suffisante} ni \emph{exclusive} pour la réussite sociale globale. \textbf{Ouverture :} Le vrai défi pour le Cameroun n'est pas de choisir entre école et terrain, mais de bâtir une \emph{"École de la Compétence"} poreuse, qui certifie les acquis de l'expérience (VAE), généralise l'apprentissage dual et valorise autant le brevet de soudeur que le doctorat en lettres. Comme le suggère la réconciliation finale entre tradition et modernité dans \emph{Le Lion et la perle}, la réussite sociale camerounaise naîtra du dialogue entre le savoir livresque et le savoir-faire ancestral. »

\tcblower
\textbf{Commentaire du corrigé (pour l'enseignant) :}
\begin{itemize}
    \item \textbf{Problématique :} Bien formulée (\emph{Dans quelle mesure...}), elle contient la thèse (voie royale légitime), l'antithèse (insuffisance, alternatives) et annonce la synthèse (complémentarité). Elle est \cle{féconde}.
    \item \textbf{Plan dialectique :} Respecté rigoureusement (I = Thèse, II = Antithèse, III = Synthèse). Les transitions explicites (pourtant, cependant, mais, vers une complémentarité) assurent la cohérence.
    \item \textbf{Exploitation des ressources :} \emph{Le Lion et la perle} utilisé deux fois comme levier d'argumentation (Lakunle = idéal scolaire déconnecté ; Baroka = pouvoir/succès coutumier), pas en résumé. Données INS citées précisément (68\%, 42\%, 35\%, 23\%). Trois exemples actualité distincts (Artisanat, Numérique, Sport).
    \item \textbf{Plan détaillé :} Chaque sous-partie (A, B) contient les 4 éléments requis : Idée directrice (phrase nominale), Argument (explicatif), Exemple (précis, daté/sourcé), Transition (lien logique). Introduction et conclusion \textbf{rédigées intégralement} (exigence Probatoire).
    \item \textbf{Langue :} Registre soutenu, connecteurs logiques variés (\emph{cependant, de surcroît, en revanche, partant de là, au terme}), ponctuation soignée (deux-points, points-virgules pour structurer les énumérations).
\end{itemize}
\end{exoresolu}

\courssub{Bilan de l'activité}
\label{bilan:integration}
Cette activité d'intégration a permis de vérifier l'acquisition effective des compétences cibles de l'unité dans une situation authentique, complexe et chronométrée, simulant l'examen :

\begin{aretenir}[Acquis validés par l'activité]
\begin{itemize}
\item \textbf{Analyse de sujet :} Les élèves distinguent désormais dénotation/connotation et produisent des problématiques \emph{ouvertes} et \emph{tendues}, non des questions fermées.
\item \textbf{Recherche d'idées :} La carte mentale impose la structuration visuelle (Thèse/Antithèse/Synthèse) et l'indexation explicite des preuves (texte, chiffres, actualité) -- fin du « vide » ou du hors-sujet.
\item \textbf{Choix du plan :} La justification écrite force la métacognition : l'élève \emph{pense} son architecture avant d'écrire. Le plan dialectique s'est imposé naturellement pour ce type de sujet.
\item \textbf{Plan détaillé :} La contrainte « Introduction/Conclusion rédigées + 4 éléments par sous-partie » ancre le formalisme attendu au Probatoire. Les transitions ne sont plus optionnelles.
\item \textbf{Co-évaluation :} La grille critériée partagée (cohérence, pertinence, connaissances, langue) objectifie la notation. La remédiation immédiate (encre verte) ancre la correction.
\item \textbf{Expression orale :} Le compte rendu de 2 min (problématique, plan, argument choc, difficulté, correction) entraîne la synthèse et la prise de parole publique.
\end{itemize}
\end{aretenir}

\begin{attention}[Points de vigilance pour la remédiation (séance suivante)]
\begin{itemize}
\item \textbf{Problématique trop descriptive :} Certains groupes ont proposé « Quelles sont les voies de la réussite ? » $\rightarrow$ Rappel : une problématique contient une \cle{tension}, un \cle{débat}.
\item \textbf{Exemples « Le Lion et la perle » :} Quelques fiches résumaient l'intrigue au lieu d'argumenter (ex: « Baroka a réussi sans école » vs « Baroka incarne la légitimité du pouvoir coutumier, alternative au modèle scolaire occidental »).
\item \textbf{Données chiffrées :} Ne pas écrire « selon l'INS » sans citer le chiffre précis (68\% vs 42\%). La précision est une preuve.
\item \textbf{Transitions faibles :} « Ensuite... » « De plus... » $\rightarrow$ Exiger la transition \emph{logique} (concession, opposition, conséquence) : « Si l'école certifie, toutefois elle ne garantit plus l'emploi... ».
\item \textbf{Ponctuation :} Relire les introductions/conclusions pour corriger les phrases trop longues (usage du point-virgule, des deux-points pour l'explication).
\end{itemize}
\end{attention}

\begin{savaistu}[Prolongement -- Projet d'études négocié]
Cette activité lance le \textbf{Projet d'études} de l'unité : chaque élève (ou binôme) choisit un sujet de dissertation en lien avec sa spécialité technique (ex: \emph{« L'intelligence artificielle menace-t-elle l'emploi technique au Cameroun ? »}, \emph{« L'apprentissage dual est-il l'avenir de la formation professionnelle ? »}) et produit, sur 3 semaines, un \textbf{dossier complet} : Analyse de sujet $\rightarrow$ 3 Problématiques $\rightarrow$ Carte mentale $\rightarrow$ Plan détaillé $\rightarrow$ Introduction et Conclusion rédigées $\rightarrow$ Auto-évaluation grillée. Ce dossier sera évalué en contrôle continu (note de devoir) et présenté à l'oral (5 min) en fin de séquence.
\end{savaistu}

\newpage

\coursec{Méthodes et exercices résolus}

\courssub{Méthodes et exercices résolus}

\begin{methode}[Analyser un sujet de dissertation et formuler la problématique]
    \textbf{Objectif :} Comprendre l'exigence du sujet, en délimiter le champ et poser la question centrale qui guidera tout le devoir.
    \begin{enumerate}
        \item \textbf{Dégager les termes du sujet :} souligner les mots-clés, les définir (dénotation/connotation) et les relier entre eux.
        \item \textbf{Identifier le type de sujet :}
            \begin{itemize}
                \item Sujet \emph{citation} (citation + « Qu'en pensez-vous ? ») : il faut discuter l'affirmation.
                \item Sujet \emph{question} directe : il faut y répondre en nuancant.
                \item Sujet \emph{thème} (nominal) : il faut le transformer en question problématisée.
            \end{itemize}
        \item \textbf{Repérer les présupposés et la tension :} quelle est l'évidence ? Quel est le paradoxe ? Où se situe le débat ?
        \item \textbf{Formuler la problématique :} c'est une \cle{question ouverte}, précise, qui contient la tension du sujet et appelle une réponse structurée (thèse/antithèse/synthèse ou constat/causes/solutions). Elle ne doit pas être une simple reformulation du sujet.
        \item \textbf{Rédiger l'accroche et l'annonce du plan :} l'accroche (citation, fait d'actualité, définition) introduit le thème. L'annonce du plan présente les \textbf{deux ou trois grandes parties} (I, II, III) en lien direct avec la problématique.
    \end{enumerate}
    \textbf{Réflexe Probatoire :} la problématique doit être écrite \textbf{en toutes lettres} à la fin de l'introduction (souvent après « Nous nous demanderons... » ou « La question qui se pose est... »).
\end{methode}

\begin{exoresolu}[Analyse de sujet et problématique — Type Probatoire]
\textbf{Sujet :} « La tradition est-elle un obstacle au progrès ? Vous répondrez à cette question en vous appuyant sur vos connaissances littéraires et votre expérience personnelle. »

\tcblower
\textbf{Solution détaillée}

\textbf{1. Dégagement des termes et définitions}
\begin{itemize}
    \item \cle{Tradition} : ensemble de coutumes, croyances, valeurs, rites transmis de génération en génération (héritage culturel). Connotation : stabilité, identité, sagesse des ancêtres, mais aussi immobilisme, routine.
    \item \cle{Progrès} : amélioration continue des conditions de vie, des connaissances, des techniques. Connotation : modernité, évolution, bien-être, mais aussi rupture, perte de repères.
    \item \cle{Obstacle} : ce qui empêche, ralentit, entrave une marche en avant. Relation dialectique : frein ou garde-fou ?
\end{itemize}

\textbf{2. Type de sujet}
Sujet \emph{question directe} (interrogation qui appelle une réponse nuancée : « Oui, mais... » / « Non, cependant... »). Il porte sur un rapport dialectique entre deux concepts abstraits.

\textbf{3. Présupposés et tension}
\begin{itemize}
    \item \emph{Présupposé 1 :} la tradition regarde vers le passé, le progrès vers l'avenir $\rightarrow$ opposition naturelle.
    \item \emph{Présupposé 2 :} le progrès suppose la rupture $\rightarrow$ la tradition serait un frein.
    \item \emph{Tension (le « mais ») :} la tradition peut être le socle indispensable au progrès (identité, éthique). Le progrès sans tradition risque d'être aveugle. La tradition vivante \emph{est} un progrès (adaptation).
\end{itemize}

\textbf{4. Formulation de la problématique}
\emph{« Dans quelle mesure la tradition, loin de s'opposer systématiquement au progrès, peut-elle en constituer le fondement nécessaire et le garde-fou éthique, à condition d'accepter de se réinventer ? »}
(Autre possibilité : « La tradition freine-t-elle le progrès ou en assure-t-elle la direction humaine ? »)

\textbf{5. Esquisse d'introduction (Accroche + Problématique + Annonce du plan)}
\textit{Accroche :} « La tradition n'est pas le culte des cendres, mais la conservation de la flamme » (formule attribuée à Jean Jaurès). Cette maxime résume le paradoxe : la tradition semble figée (cendres) mais porte une énergie vitale (flamme) pour l'avenir.
\textit{Problématique :} celle formulée ci-dessus.
\textit{Annonce du plan :}
\begin{itemize}
    \item \textbf{I.} Nous analyserons d'abord les formes de conflictualité où la tradition apparaît comme un frein au changement technique et social.
    \item \textbf{II.} Nous montrerons ensuite que la tradition fournit les repères identitaires et éthiques sans lesquels le progrès perd son sens humain.
    \item \textbf{III.} Enfin, nous démontrerons que le véritable progrès naît d'une \emph{tradition vivante}, capable d'intégrer la modernité par une réinterprétation critique.
\end{itemize}

\textbf{Conclusion de l'exercice :} l'analyse rigoureuse des termes et de la tension dialectique permet de sortir du binaire « pour/contre » pour construire une problématique dynamique, gage d'un plan dialectique réussi.
\end{exoresolu}

\begin{methode}[Rechercher, sélectionner et organiser les idées]
    \textbf{Objectif :} passer de la problématique à une matière brute structurée (arguments, exemples, références) avant le plan.
    \begin{enumerate}
        \item \textbf{Remue-méninges / carte mentale :} noter \emph{tout} ce qui vient (mots-clés, auteurs, exemples personnels, faits d'actualité, scènes du \emph{Lion et la Perle}) sans trier. Utiliser un schéma radial : problématique au centre, branches = grandes directions de réponse (I, II, III pressentis).
        \item \textbf{Classification par \cle{grandes parties} :} ranger chaque idée sous la partie correspondante (I, II, III). Vérifier l'équilibre : 2 à 3 arguments par partie.
        \item \textbf{Structuration interne de chaque partie :} pour chaque grande partie, ordonner les arguments :
            \begin{itemize}
                \item Argument 1 (le plus fort / le plus général) $\rightarrow$ exemple précis (littérature, histoire, actualité, expérience vécue).
                \item Argument 2 (nuance / contre-exemple / approfondissement) $\rightarrow$ exemple.
                \item Argument 3 (ouverture / conséquence) $\rightarrow$ exemple.
            \end{itemize}
        \item \textbf{Sélection rigoureuse (élagage) :} supprimer le hors-sujet, les doublons, les exemples trop vagues (« on dit que... »). Privilégier la \cle{précision} : « Dans \emph{Le Lion et la Perle}, Lakunle incarne... » vaut mieux que « Au théâtre africain... ».
        \item \textbf{Vérification de la cohérence :} chaque argument répond-il à la problématique ? La progression I $\rightarrow$ II $\rightarrow$ III est-elle logique (thèse $\rightarrow$ antithèse $\rightarrow$ synthèse) ?
    \end{enumerate}
    \textbf{Réflexe Probatoire :} prévoir une \cle{colonne « Références »} dans le brouillon : auteur / œuvre / personnage / scène précise. C'est la traçabilité exigée à l'examen.
\end{methode}

\begin{exoresolu}[Recherche d'idées pour le sujet : « Tradition et progrès »]
\textbf{Consigne :} à partir de la problématique et de l'annonce de plan de l'exercice précédent, établir la fiche de recherche d'idées (brouillon structuré) avec arguments et références précises (dont \emph{Le Lion et la Perle}).

\tcblower
\textbf{Solution détaillée : fiche de recherche d'idées (brouillon)}

\paragraph{Problématique :} Dans quelle mesure la tradition peut-elle être le fondement et le garde-fou du progrès ?

\begin{center}
\begin{tabularx}{\linewidth}{|l|X|X|}\hline
\rowcolor{popblueL}\textbf{Partie / Argument} & \textbf{Développement / Idées forces} & \textbf{Références précises} \\ \hline
\textbf{I. Tradition = frein au progrès ? (Thèse)} & & \\ \hline
A. Poids du passé / immobilisme & Respect aveugle des coutumes, refus du nouveau, peur du changement. & \emph{Le Lion et la Perle} : Sadiku et les femmes du village moquent la machine à timbrer ; proverbe : « On ne change pas une équipe qui gagne ». \\ \hline
B. Pratiques néfastes / droits humains & Mariages forcés, exclusion, patriarcat rigide. Le progrès = émancipation individuelle. & \emph{Le Lion et la Perle} : la dot (\emph{bride-price}) : Lakunle refuse de payer la dot au nom de l'égalité moderne. \\ \hline
C. Obstacle technique / scientifique & Croyances contre méthode scientifique ; rejet de l'école ou de la médecine moderne. & Chinua Achebe, \emph{Le Monde s'effondre} : le conflit entre Okonkwo et les missionnaires. \\ \hline
\textbf{II. Tradition = socle et garde-fou (Antithèse)} & & \\ \hline
A. Identité et mémoire collective & Le progrès sans racines = déracinement, aliénation culturelle (mimétisme). & \emph{Le Lion et la Perle} : la danse du « voyage de l'étranger » : la tradition \emph{intègre} la modernité par le mime. \\ \hline
B. Éthique et solidarité & Humanisme communautaire. Le progrès technique sans éthique mène à la barbarie. & Cheikh Hamidou Kane, \emph{L'Aventure ambiguë} : la nécessité des racines pour l'assise de l'homme. \\ \hline
C. Savoirs endogènes & Pharmacopée, agriculture durable, architecture adaptée au climat. & Exemples camerounais : médecine traditionnelle ; cases obus mousgoums (pouvoir isolant). \\ \hline
\textbf{III. Tradition vivante = progrès réel (Synthèse)} & & \\ \hline
A. Réinterprétation critique & Tri sélectif : garder l'esprit, changer la lettre. « Tradition vivante » (Amadou Hampâté Bâ). & \emph{Le Lion et la Perle} : Sidi choisit Baroka (tradition) mais a conscience de sa valeur (modernité : les photos). Baroka utilise le timbre et la machine. \\ \hline
B. Modernité endogène & Progrès adapté au contexte local ; ne pas copier-coller les modèles étrangers. & Exemple : le mobile money (Orange Money) adapte la pratique traditionnelle de la tontine à la technologie. \\ \hline
C. Dialogue des cultures & Métissage culturel ; la tradition s'enrichit de l'apport extérieur. & Léopold Sédar Senghor : « enracinement et ouverture » ; la danse finale du \emph{Lion et la Perle} unit tradition et modernité. \\ \hline
\end{tabularx}
\end{center}

\textbf{Conclusion de l'exercice :} cette fiche montre un équilibre (3 arguments par partie), des références variées (œuvre au programme, auteurs africains, exemples camerounais) et une progression dialectique claire. Elle est prête pour le plan détaillé.
\end{exoresolu}

\begin{methode}[Élaborer le plan détaillé]
    \textbf{Objectif :} transformer la fiche d'idées en un squelette rigoureux, numéroté, qui servira de guide unique à la rédaction. Le plan détaillé est \textbf{obligatoire} au brouillon et souvent noté.
    \begin{enumerate}
        \item \textbf{Choisir la structure adaptée à la problématique :}
            \begin{itemize}
                \item \cle{Plan dialectique (thèse / antithèse / synthèse)} : standard pour les sujets « Est-ce que... », « Faut-il... » (I. Oui / II. Non / III. Dépassement).
                \item \cle{Plan analytique (constat / causes / conséquences / solutions)} : pour les sujets « Quelles sont les causes de... », « Quelles solutions pour... ».
                \item \cle{Plan thématique (aspects : social, politique, culturel)} : pour les sujets larges « Étudiez le thème de... ». \textbf{À éviter} pour une dissertation argumentative classique (risque de simple liste).
            \end{itemize}
        \item \textbf{Rédiger les titres des parties (I, II, III) :} phrases nominales ou verbales, parallèles, synthétiques. Elles résument l'argument principal de la partie. \textbf{Pas de « I. Les arguments pour ».}
        \item \textbf{Détailler chaque partie en 2 ou 3 sous-parties (A, B, C) :} chaque sous-partie = \textbf{un argument unique} + \textbf{son exemple précis}. Formulation : « A. [Argument] : [exemple/référence] ».
        \item \textbf{Prévoir les transitions (liaisons logiques) :} entre I et II (opposition : « Cependant », « En revanche »), entre II et III (dépassement : « Dès lors », « C'est pourquoi »). Les écrire dans le plan.
        \item \textbf{Vérifier la « tenue » du plan :}
            \begin{itemize}
                \item équilibre visuel (longueur similaire de I, II, III) ;
                \item cohérence interne (A, B, C démontrent bien le titre de la partie) ;
                \item réponse directe à la problématique (la lecture du plan seul doit permettre de deviner la problématique).
            \end{itemize}
    \end{enumerate}
    \textbf{Réflexe Probatoire :} le plan détaillé se rédige \textbf{au brouillon}, lisible et encadré. Il ne doit pas être recopié sur la copie finale (sauf demande explicite), mais il structure la rédaction.
\end{methode}

\begin{exoresolu}[Élaboration du plan détaillé — Sujet « Tradition et progrès »]
\textbf{Consigne :} produire le plan détaillé complet (titres I, II, III, sous-parties A, B, C, transitions) à partir de la fiche d'idées précédente. Respecter la mise en page type Probatoire.

\tcblower
\textbf{Solution détaillée : plan détaillé (brouillon)}

\paragraph{Structure choisie :} dialectique (thèse / antithèse / synthèse) — adaptée à la tension « frein ou fondement ».

\begin{center}
\begin{tabularx}{\linewidth}{|X|}\hline
\rowcolor{popblueL}\textbf{PLAN DÉTAILLÉ} \\ \hline
\textbf{INTRODUCTION} \\
\hspace{0.5cm} Accroche : citation « Flamme / Cendres ». \\
\hspace{0.5cm} Présentation du sujet / définitions (tradition, progrès). \\
\hspace{0.5cm} \textbf{Problématique :} \emph{« Dans quelle mesure la tradition, loin de s'opposer systématiquement au progrès, peut-elle en constituer le fondement nécessaire et le garde-fou éthique, à condition d'accepter de se réinventer ? »} \\
\hspace{0.5cm} Annonce du plan. \\ \hline
\textbf{I. La tradition peut apparaître comme un frein au progrès technique et social.} \\
\hspace{0.5cm} \textbf{A. L'immobilisme coutumier : le refus du changement au nom du respect des ancêtres.} \\
\hspace{1cm} (Ex. : \emph{Le Lion et la Perle} — Sadiku et les femmes du village hostiles à la machine à timbrer ; proverbe « On ne change pas une équipe qui gagne »). \\
\hspace{0.5cm} \textbf{B. La perpétuation de pratiques contraires à l'émancipation individuelle.} \\
\hspace{1cm} (Ex. : \emph{Le Lion et la Perle} — la dot comme marchandage de la femme ; Lakunle la refuse au nom de l'égalité). \\
\hspace{0.5cm} \textbf{C. Le refus de la raison scientifique et de l'école.} \\
\hspace{1cm} (Ex. : Achebe, \emph{Le Monde s'effondre} — conflit entre Okonkwo et les missionnaires). \\
\hspace{0.5cm} \textit{\textbf{[Transition I $\rightarrow$ II :]}} \emph{« Cependant, réduire la tradition à ses formes figées serait ignorer sa fonction vitale de cohésion et de sens. »} \\ \hline
\textbf{II. La tradition est le socle identitaire et le garde-fou éthique indispensable à tout progrès humain.} \\
\hspace{0.5cm} \textbf{A. L'ancrage identitaire : condition de la dignité face à l'uniformisation.} \\
\hspace{1cm} (Ex. : \emph{Le Lion et la Perle} — la danse du « voyage » mime la voiture en langage traditionnel ; Senghor : « enracinement et ouverture »). \\
\hspace{0.5cm} \textbf{B. L'humanisme communautaire : rempart contre la barbarie technicienne.} \\
\hspace{1cm} (Ex. : \emph{L'Aventure ambiguë} — la nécessité des racines ; Baroka, malgré ses défauts, protège son village). \\
\hspace{0.5cm} \textbf{C. La valorisation des savoirs endogènes : une ressource pour l'innovation durable.} \\
\hspace{1cm} (Ex. : pharmacopée traditionnelle ; architecture bioclimatique mousgoum ; agriculture paysanne). \\
\hspace{0.5cm} \textit{\textbf{[Transition II $\rightarrow$ III :]}} \emph{« Dès lors, le progrès véritable ne naît pas de la rupture brutale, mais de la capacité de la tradition à se réinventer par un dialogue critique avec la modernité. »} \\ \hline
\textbf{III. Le progrès réel suppose une tradition vivante, capable d'intégration critique et de métissage fécond.} \\
\hspace{0.5cm} \textbf{A. L'innovation par réinterprétation : trier, adapter, réinventer (tradition vivante).} \\
\hspace{1cm} (Ex. : \emph{Le Lion et la Perle} — Sidi : choix traditionnel (Baroka) mais conscience moderne de sa valeur ; Baroka : chef traditionnel utilisant le timbre et la machine). \\
\hspace{0.5cm} \textbf{B. La modernité endogène : adapter le progrès au contexte local.} \\
\hspace{1cm} (Ex. : Orange Money = tontine numérique ; adaptation des pratiques d'épargne collective à la technologie). \\
\hspace{0.5cm} \textbf{C. Le dialogue des cultures : l'enrichissement mutuel comme horizon.} \\
\hspace{1cm} (Ex. : la danse finale du \emph{Lion et la Perle} ; Senghor : le « rendez-vous du donner et du recevoir »). \\
\hspace{0.5cm} \textit{\textbf{[Ouverture :]}} \emph{« Et si le défi africain était d'inventer un progrès enraciné dans sa propre mémoire ? »} \\ \hline
\textbf{CONCLUSION} \\
\hspace{0.5cm} Bilan synthétique (reprise de I, II, III). \\
\hspace{0.5cm} Réponse définitive à la problématique : la tradition n'est pas l'ennemie du progrès, mais sa boussole, à condition d'être vivante. \\
\hspace{0.5cm} Ouverture sur l'éducation et la transmission intergénérationnelle. \\ \hline
\end{tabularx}
\end{center}

\textbf{Conclusion de l'exercice :} ce plan est équilibré (3 parties, 3 sous-parties), les titres sont des phrases affirmatives, les exemples sont précis et variés (œuvre au programme, auteurs, contexte camerounais), les transitions sont explicites. Il est prêt pour la rédaction.
\end{exoresolu}

\begin{methode}[Maîtriser les outils linguistiques de la dissertation]
    \textbf{Objectif :} assurer la cohérence, la fluidité et la tonalité argumentative attendue (registre soutenu, objectivité, nuance).
    \begin{enumerate}
        \item \textbf{La ponctuation (syntaxe du sens) :}
            \begin{itemize}
                \item \cle{Point-virgule (;)} : relie deux propositions indépendantes liées logiquement (cause/conséquence, opposition). \emph{Ex. : « La tradition unit ; le progrès isole. »}
                \item \cle{Deux-points (:)} : annonce une explication, une énumération, une citation. \emph{Ex. : « Trois arguments militent pour cette thèse : ... »}
                \item \cle{Parenthèses / tirets} : incises pour une précision, une concession, un exemple. \emph{Ex. : « Baroka — le Baalé du village — incarne... »}
            \end{itemize}
        \item \textbf{Les connecteurs logiques (opérateurs argumentatifs) :} les classer par fonction et les varier.
            \begin{center}
            \begin{tabular}{|l|l|}\hline
            \rowcolor{popblueL}\textbf{Fonction} & \textbf{Connecteurs (variété attendue)} \\ \hline
            \textbf{Ajout / renforcement} & de surcroît, qui plus est, en outre, par ailleurs, non seulement... mais encore \\ \hline
            \textbf{Opposition / concession} & cependant, pourtant, en revanche, toutefois, certes... mais \\ \hline
            \textbf{Cause / conséquence} & c'est pourquoi, par conséquent, dès lors, du fait que, eu égard à \\ \hline
            \textbf{Illustration / exemple} & notamment, à titre d'exemple, c'est le cas de, ainsi \\ \hline
            \textbf{Conclusion / synthèse} & en définitive, au total, en somme, force est de constater que \\ \hline
            \end{tabular}
            \end{center}
        \item \textbf{L'énonciation (marques du locuteur) :}
            \begin{itemize}
                \item \cle{Impersonnalisation / objectivité :} « Il convient de noter que... », « Il est indéniable que... », « On observe que... » (préférer « on » ou le passif à « je pense que »).
                \item \cle{Modalisation (nuance) :} adverbes (probablement, sans doute, apparemment), verbes modaux (pouvoir, devoir, sembler), conditionnel. \emph{Éviter les certitudes absolues :} « C'est certain que » $\rightarrow$ « Il semble que ».
                \item \cle{Marques explicites :} « Nous allons voir... », « Signalons que... », « Rappelons que... » (elles guident le correcteur).
            \end{itemize}
        \item \textbf{Le texte descriptif (pour les exemples littéraires) :}
            \begin{itemize}
                \item vocabulaire précis (spatial : \emph{au premier plan, en arrière-plan} ; sensoriel : \emph{claquant, suave, rugueux}) ;
                \item figures de style : comparaison, métaphore, personnification (pour faire « voir » la scène) ;
                \item temps : imparfait (description statique), présent de narration (action au théâtre).
            \end{itemize}
    \end{enumerate}
    \textbf{Réflexe Probatoire :} relire \textbf{uniquement pour la ponctuation et les connecteurs} en fin d'épreuve. Une phrase mal ponctuée change le sens logique.
\end{methode}

\begin{exoresolu}[Application linguistique : rédaction d'un paragraphe type (Partie I, A)]
\textbf{Consigne :} rédiger le paragraphe correspondant à la sous-partie \textbf{I.A} du plan détaillé : « L'immobilisme coutumier : le refus du changement au nom du respect des ancêtres ». Intégrer : connecteurs, énonciation objective, description précise d'une scène du \emph{Lion et la Perle}, ponctuation soignée.

\tcblower
\textbf{Solution détaillée : paragraphe type}

\emph{« En premier lieu, la tradition se fige parfois en un immobilisme coutumier qui érige le respect des ancêtres en refus systématique de toute innovation. En effet, l'attachement aux usages ancestraux devient alors un principe de conservation exclusif, rejetant par principe toute forme de modernité perçue comme une profanation. C'est ce que met en lumière Wole Soyinka dans \emph{Le Lion et la Perle} : lorsque les femmes du village, conduites par Sadiku, miment la scène de la machine à timbrer, elles transforment l'outil moderne en objet de dérision grotesque ; le bruit de la machine devient le symbole d'une modernité vide de sens, que l'on oppose à la parole vivante et au rythme ancestral. De surcroît, le proverbe « On ne change pas une équipe qui gagne » résume cette logique circulaire : puisque la tradition a assuré la survie du groupe hier, elle garantirait l'avenir demain, dispensant ainsi de toute réflexion critique. Dès lors, c'est au nom de la fidélité au passé que l'on hypothèque l'avenir, en confondant tradition et traditionalisme. »}

\textbf{Analyse linguistique du paragraphe :}
\begin{itemize}
    \item \textbf{Connecteurs :} \emph{en premier lieu} (ordre), \emph{en effet} (explication), \emph{c'est ce que} (focalisation), \emph{lorsque} (circonstance), \emph{de surcroît} (renforcement), \emph{dès lors} (conséquence logique).
    \item \textbf{Énonciation :} \emph{il devient, c'est ce que met en lumière, on oppose, résume cette logique} (impersonnel, objectif). Pas de « je ».
    \item \textbf{Modalisation :} \emph{perçue comme, devient le symbole, garantirait} (verbes d'analyse et conditionnel de nuance, pas de jugements de valeur bruts).
    \item \textbf{Ponctuation :} point-virgule pour lier description et interprétation ; virgules pour les incises (\emph{conduites par Sadiku}).
    \item \textbf{Description :} \emph{miment la scène, bruit de la machine, parole vivante, rythme ancestral} (visuel, auditif, abstrait).
\end{itemize}
\textbf{Conclusion :} ce paragraphe respecte le format « argument $\rightarrow$ exemple précis (description/analyse) $\rightarrow$ interprétation $\rightarrow$ mini-conclusion partielle ».
\end{exoresolu}

\begin{methode}[Rédiger l'introduction et la conclusion]
    \textbf{Objectif :} maîtriser les deux « vitrines » de la dissertation, très surveillées à l'examen.
    \begin{enumerate}
        \item \textbf{L'introduction (entonnoir : large $\rightarrow$ précis) — 4 étapes obligatoires :}
            \begin{enumerate}
                \item \textbf{Accroche :} citation d'auteur (précise : auteur, œuvre), fait d'actualité vérifié, définition, paradoxe, référence au texte ouvroir (\emph{Le Lion et la Perle}). \textbf{Interdit :} « Depuis la nuit des temps... », « De tout temps les hommes... » (lieux communs vides).
                \item \textbf{Présentation du sujet / définitions :} reprise des termes du sujet, définitions (dénotation/connotation), mise en relation. Situer le débat.
                \item \textbf{Problématique :} \textbf{la question unique, problématisée, écrite en toutes lettres.} Souvent introduite par : « Nous nous demanderons dans quelle mesure... », « La question qui se pose est de savoir si... ».
                \item \textbf{Annonce du plan :} « Pour y répondre, nous analyserons d'abord... (I). Nous montrerons ensuite... (II). Enfin, nous verrons... (III). » \textbf{Utiliser les titres exacts des parties I, II, III.}
            \end{enumerate}
        \item \textbf{La conclusion (entonnoir inverse : précis $\rightarrow$ large) — 3 étapes obligatoires :}
            \begin{enumerate}
                \item \textbf{Bilan / synthèse :} reprise \emph{synthétique} (pas une liste) de la démonstration : « Au terme de cette analyse, il apparaît que... ». Réponse directe à la problématique.
                \item \textbf{Réponse définitive :} une phrase affirmative, nuancée, qui clôt le débat. « Ainsi, la tradition n'est pas l'ennemie du progrès, mais sa condition humaine. »
                \item \textbf{Ouverture :} prolongement pertinent (autre œuvre, autre discipline, actualité, question éthique). \textbf{Interdit :} « Pour conclure, on peut dire que... » (répétition), ouverture hors sujet, question rhétorique non traitée.
            \end{enumerate}
    \end{enumerate}
    \textbf{Réflexe Probatoire :} introduction et conclusion se rédigent \textbf{au propre directement} (ou au brouillon puis recopiées sans fautes). Elles doivent rester brèves (une dizaine de lignes chacune au maximum). Ne pas les négliger : elles pèsent lourd dans la notation.
\end{methode}

\begin{exoresolu}[Rédaction de l'introduction et de la conclusion — Sujet « Tradition et progrès »]
\textbf{Consigne :} rédiger l'introduction et la conclusion complètes pour le devoir, en appliquant la méthode (accroche — définitions — problématique — plan ; bilan — réponse — ouverture).

\tcblower
\textbf{Solution détaillée}

\paragraph{INTRODUCTION}
« La tradition n'est pas le culte des cendres, mais la conservation de la flamme. » Par cette formule lumineuse, attribuée à Jean Jaurès, se résume le paradoxe au cœur de notre sujet : la tradition, souvent perçue comme un poids mort (\emph{cendres}), porte en elle une énergie vitale (\emph{flamme}) capable d'éclairer l'avenir. Le sujet oppose deux forces majeures de l'histoire humaine : la \textbf{tradition}, ensemble des coutumes, valeurs et savoirs transmis (héritage, stabilité), et le \textbf{progrès}, dynamique d'amélioration technique, sociale et morale (rupture, avenir). L'une semble regarder en arrière, l'autre en avant ; l'une fige, l'autre transforme. Dès lors, la relation entre les deux apparaît conflictuelle : la tradition serait un obstacle que le progrès doit briser pour advenir. \textbf{Nous nous demanderons pourtant dans quelle mesure la tradition, loin de s'opposer systématiquement au progrès, peut en constituer le fondement nécessaire et le garde-fou éthique, à condition d'accepter de se réinventer.} Pour y répondre, nous analyserons d'abord les formes de conflictualité où la tradition apparaît comme un frein au changement technique et social \textbf{(I)}. Nous montrerons ensuite que la tradition fournit les repères identitaires et éthiques sans lesquels le progrès perd son sens humain \textbf{(II)}. Enfin, nous démontrerons que le véritable progrès naît d'une tradition vivante, capable d'intégrer la modernité par une réinterprétation critique \textbf{(III)}.

\paragraph{CONCLUSION}
Au terme de cette étude, il apparaît que la tradition n'est pas un bloc monolithique figé dans le passé. \textbf{Si} elle peut entraver le progrès lorsqu'elle se fige en traditionalisme — refus de l'école, pratiques néfastes, immobilisme coutumier (I) —, \textbf{elle} en constitue paradoxalement le socle indispensable : gardienne de l'identité face à l'uniformisation, réservoir d'humanisme face à la technique aveugle, source de savoirs endogènes durables (II). \textbf{Surtout}, la littérature (\emph{Le Lion et la Perle}) nous enseigne que le progrès réel, celui qui humanise, naît du métissage fécond : une tradition vivante qui trie, adapte et réinvente ses formes pour accueillir la modernité sans s'y dissoudre (III). \textbf{Ainsi, la tradition n'est pas l'obstacle au progrès, mais sa boussole éthique et culturelle.} Cette conclusion invite à repenser l'éducation au Cameroun : comment former des citoyens « enracinés et ouverts », capables de faire dialoguer la sagesse des ancêtres et les exigences du XXI\textsuperscript{e} siècle ?

\textbf{Analyse :}
\begin{itemize}
    \item \textbf{Introduction :} accroche citée, définitions liées, problématique \textbf{en gras} (repérable), annonce du plan avec numéros romains.
    \item \textbf{Conclusion :} bilan synthétique (si... si... surtout...), réponse définitive (\textbf{ainsi...}), ouverture ciblée (éducation, enracinement et ouverture).
\end{itemize}
\end{exoresolu}

\begin{methode}[Réaliser un compte rendu de lecture (Le Lion et la Perle) et s'auto-évaluer]
    \textbf{Objectif :} maîtriser l'exercice codifié du compte rendu (résumé + analyse + appréciation) et la grille de remédiation pour progresser.
    \begin{enumerate}
        \item \textbf{Structure type du compte rendu (4 parties distinctes) :}
            \begin{enumerate}
                \item \textbf{Identification / fiche technique :} titre, auteur, date, genre (comédie, théâtre), contexte (Nigéria, peuple yoruba, veille des indépendances).
                \item \textbf{Résumé de l'intrigue (15-20 lignes max) :} ordre chronologique, présent de narration. Focus sur le nœud (Sidi / Baroka / Lakunle), le climax (la ruse de Baroka), le dénouement (choix de Sidi, danse finale). \textbf{Pas de commentaire ici.}
                \item \textbf{Analyse / commentaire (cœur de l'exercice) :}
                    \begin{itemize}
                        \item \textbf{Thèmes :} tradition et modernité, pouvoir et séduction, condition féminine, langage et image.
                        \item \textbf{Personnages :} Baroka = le Lion, tradition rusée ; Lakunle = modernité caricaturale et stérile ; Sidi = la Perle, beauté devenue consciente d'elle-même.
                        \item \textbf{Procédés théâtraux :} mise en abyme (la danse du voyage), proverbes, chants et danses (oralité), satire et ironie.
                        \item \textbf{Portée :} universalité du conflit, réhabilitation de la tradition vivante.
                    \end{itemize}
                \item \textbf{Appréciation critique personnelle (5-10 lignes) :} « J'ai aimé / je n'ai pas aimé parce que... », argumentée (style, idées, personnages). Éviter « c'est bien / ennuyeux ».
            \end{enumerate}
        \item \textbf{Remédiation / auto-correction (grille type) :} après chaque devoir, remplir :
            \begin{center}
            \begin{tabularx}{\linewidth}{|l|X|X|}\hline
            \rowcolor{popblueL}\textbf{Compétence évaluée} & \textbf{Points forts (réussites)} & \textbf{Axes de progression (erreurs récurrentes)} \\ \hline
            Analyse du sujet / problématique & & \\ \hline
            Plan / structure & & \\ \hline
            Arguments / exemples & & \\ \hline
            Expression / langue (ponctuation, connecteurs) & & \\ \hline
            Introduction / conclusion & & \\ \hline
            \end{tabularx}
            \end{center}
        \item \textbf{Projet d'études :} fixer deux objectifs concrets pour le prochain devoir (ex. : « intégrer une référence précise par partie », « varier les connecteurs d'opposition », « soigner l'ouverture »).
    \end{enumerate}
    \textbf{Réflexe Probatoire :} le compte rendu du \emph{Lion et la Perle} est un classique de l'examen. Connaître \textbf{quelques répliques clés} et la structure en trois moments (Matin, Midi, Soir) est un atout majeur.
\end{methode}

\begin{exoresolu}[Compte rendu type — Le Lion et la Perle (Wole Soyinka)]
\textbf{Consigne :} rédiger un compte rendu de lecture structuré (fiche technique, résumé, analyse, appréciation) utilisable comme fiche de révision.

\tcblower
\textbf{Solution détaillée : fiche de lecture \emph{Le Lion et la Perle}}

\paragraph{1. Fiche technique}
\begin{itemize}
    \item \textbf{Titre :} \emph{The Lion and the Jewel} (publié en 1963) — traduit en français sous le titre \emph{Le Lion et la Perle}.
    \item \textbf{Auteur :} Wole Soyinka (Nigéria, prix Nobel de littérature 1986).
    \item \textbf{Genre :} comédie satirique en trois parties (Matin, Midi, Soir) — théâtre africain d'expression anglaise.
    \item \textbf{Contexte :} Nigéria à la veille de l'indépendance, village yoruba fictif d'Ilujinle. Tension entre tradition et modernité.
\end{itemize}

\paragraph{2. Résumé de l'intrigue}
\emph{Matin :} Lakunle, instituteur « moderne », veut épouser Sidi (la « Perle », beauté du village) sans payer la dot, coutume qu'il juge archaïque. Sidi refuse un mariage sans dot, qui la déshonorerait. Baroka (le « Lion », le Baalé, chef du village, soixante-deux ans) convoite lui aussi Sidi. Un photographe étranger a immortalisé Sidi : ses images paraissent dans un magazine et la rendent célèbre.
\emph{Midi :} les villageois miment la « danse du voyage de l'étranger » (l'arrivée de la voiture et du photographe). Sadiku, épouse aînée de Baroka, annonce à Sidi la demande en mariage du chef. Sidi refuse avec orgueil : sa photo lui a révélé sa valeur. Sadiku lui révèle alors le « secret » de Baroka : il serait devenu impuissant. Sidi, triomphante, décide d'aller narguer le vieux chef chez lui.
\emph{Soir :} chez Baroka, Sidi découvre que l'impuissance était une ruse. Baroka la séduit par son intelligence, son verbe et ses projets modernes (la machine à timbrer, le timbre-poste à l'effigie de Sidi). Sidi cède. De retour au village, elle annonce qu'elle épouse Baroka. Lakunle, soulagé de ne pas payer de dot, se tourne vers une autre jeune fille. La pièce s'achève par une danse collective.

\paragraph{3. Analyse}
\textbf{Thèmes majeurs :}
\begin{itemize}
    \item \cle{Tradition et modernité :} rapport non binaire. Baroka (tradition) utilise le moderne (timbre, machine, magazine). Lakunle (modernité) est ridicule, déraciné, stérile. Sidi synthétise les deux : beauté traditionnelle, conscience moderne de son image.
    \item \cle{Pouvoir et séduction :} Baroka = pouvoir politique, séduction masculine et ruse.
    \item \cle{Condition féminine :} Sidi passe du statut d'objet de convoitise (dot, photo) à celui de sujet qui choisit. Sadiku incarne le pouvoir féminin souterrain (messagère, complice).
    \item \cle{Langage et image :} la photo donne pouvoir à Sidi ; la parole (proverbes, ruses) est l'arme de Baroka.
\end{itemize}
\textbf{Personnages :}
\begin{itemize}
    \item \textbf{Baroka :} le « Lion » (force, ruse, longévité). Conservateur \emph{mais} adaptatif : il incarne la tradition vivante.
    \item \textbf{Lakunle :} l'instituteur caricatural. Modernité mimétique, vide, mal assimilée. Échec total.
    \item \textbf{Sidi :} la « Perle » (beauté, valeur). Conscience de sa valeur ; choix final à la fois pragmatique et identitaire.
\end{itemize}
\textbf{Procédés théâtraux :}
\begin{itemize}
    \item \cle{Mise en abyme :} la danse du voyage (Midi) résume en mime l'intrusion de la modernité dans la tradition.
    \item \cle{Oralité :} proverbes, chants, mimes, danses : un théâtre total.
    \item \cle{Ironie et satire :} Lakunle (cible principale), la bureaucratie (machine à timbrer), la polygamie (traitée avec humour).
\end{itemize}

\paragraph{4. Appréciation personnelle}
« J'ai particulièrement apprécié le refus du manichéisme : Soyinka ne sacralise pas la tradition (Baroka est manipulateur et polygame) ni ne diabolise la modernité (le timbre-poste est utile). La langue, vibrante d'oralité yoruba, fait entendre l'Afrique. La fin dansée, réconciliatrice, offre une vision optimiste d'une modernité assumée par la tradition. »

\paragraph{5. Repères à connaître pour l'examen}
\begin{enumerate}
    \item La structure en trois moments : Matin, Midi, Soir (le temps d'une seule journée).
    \item Les surnoms symboliques : Baroka = le Lion ; Sidi = la Perle.
    \item La ruse de Baroka : la fausse impuissance, ressort comique et dramatique majeur.
    \item Le refus de la dot par Lakunle : motif déclencheur du conflit.
    \item La danse finale : réconciliation de la tradition et de la modernité.
\end{enumerate}
\end{exoresolu}

\begin{attention}[Les 5 erreurs qui coûtent des points à l'examen (dissertation)]
    \begin{enumerate}
        \item \textbf{Problématique absente, floue ou hors sujet :} l'introduction se termine sans question claire, ou la question est une simple reformulation du sujet (« La tradition est-elle un obstacle ? »). \emph{Conséquence :} plan incohérent, note plafonnée. \textbf{Réflexe :} écrire la problématique \textbf{en toutes lettres}, la souligner, vérifier qu'elle contient une tension.
        \item \textbf{Plan « catalogue » (thématique au lieu de dialectique) :} I. La tradition. II. Le progrès. III. La tradition et le progrès. \emph{Conséquence :} pas d'argumentation, pas de progression logique, répétitions. \textbf{Réflexe :} les titres de parties doivent être des \textbf{arguments répondant à la problématique} (thèse / antithèse / synthèse).
        \item \textbf{Exemples vagues (« on dit que », « dans la vie quotidienne ») :} absence de références précises (auteur, œuvre, personnage, scène, lieu précis). \emph{Conséquence :} la compétence « exemples » n'est pas validée. \textbf{Réflexe :} se constituer une « banque d'exemples » par thème (littérature, histoire, actualité camerounaise, expérience vécue précise).
        \item \textbf{Mauvaise gestion de l'énonciation et langage familier :} « je pense que », « c'est clair que », « vraiment », « trop », phrases sans verbe, ponctuation aléatoire. \emph{Conséquence :} pénalité lourde sur l'expression. \textbf{Réflexe :} relire \textbf{uniquement pour la langue} (connecteurs, tournures impersonnelles, conditionnel, ponctuation).
        \item \textbf{Conclusion bâclée ou absente, ouverture « poubelle » :} arrêt brutal au milieu de la partie III ; ouverture sans rapport (« on peut aussi parler de l'écologie »). \emph{Conséquence :} perte de points faciles. \textbf{Réflexe :} gérer le temps, rédiger la conclusion \textbf{au propre} directement, choisir une ouverture en \textbf{prolongement logique} du sujet (éducation, politique culturelle, numérique).
    \end{enumerate}
\end{attention}

\newpage

\coursec{Exercices}

\courssub{Je vérifie mes connaissances}

\begin{exercice}[QCM : la dissertation]
Pour chaque question, une seule réponse est correcte. Entoure la lettre correspondante.
\begin{enumerate}
\item La problématique d'une dissertation est :
\begin{enumerate}
\item[a)] le titre du sujet recopié ;
\item[b)] la question centrale qui dirige toute la réflexion ;
\item[c)] la conclusion du devoir.
\end{enumerate}
\item Le plan dialectique comporte :
\begin{enumerate}
\item[a)] deux parties : thèse et antithèse ;
\item[b)] trois parties : thèse, antithèse, synthèse ;
\item[c)] une seule partie développée.
\end{enumerate}
\item Dans l'introduction, l'amorce (ou accroche) sert à :
\begin{enumerate}
\item[a)] donner la réponse finale ;
\item[b)] entrer dans le sujet en partant d'une idée générale ;
\item[c)] annoncer le plan.
\end{enumerate}
\item Une idée est « hors sujet » lorsqu'elle :
\begin{enumerate}
\item[a)] est mal rédigée ;
\item[b)] ne répond pas à la question posée par le sujet ;
\item[c)] est trop courte.
\end{enumerate}
\end{enumerate}
\end{exercice}

\begin{exercice}[Vrai ou faux]
Réponds par vrai ou faux, puis justifie ta réponse en une phrase.
\begin{enumerate}
\item L'annonce du plan se place à la fin de l'introduction.
\item Dans un plan détaillé, chaque sous-partie doit contenir au moins un exemple précis.
\item La dénotation d'un mot est le sens imagé ou affectif qu'il évoque.
\item On peut rédiger le développement avant d'avoir formulé la problématique.
\item Le point-virgule peut séparer deux propositions indépendantes de sens proche.
\end{enumerate}
\end{exercice}

\begin{exercice}[Définitions]
Définis en deux ou trois phrases chacun des termes suivants, puis donne un exemple pour les deux derniers :
\begin{enumerate}
\item la problématique ;
\item le plan détaillé ;
\item la dénotation ;
\item la connotation.
\end{enumerate}
\end{exercice}

\begin{exercice}[Paratexte et communication]
On donne la couverture du roman \emph{Le Lion et la perle} de Wole Soyinka : titre, nom de l'auteur, illustration représentant un jeune instituteur et une jeune fille en tenue traditionnelle, mention de l'éditeur.
\begin{enumerate}
\item Cite trois éléments du paratexte relevés sur cette couverture.
\item Pour chaque élément, indique ce qu'il permet de prévoir sur le contenu de l'œuvre.
\item Identifie dans la situation de lecture de ce roman : l'émetteur, le récepteur, le message, le canal et le code.
\end{enumerate}
\end{exercice}

\courssub{Je m'entraîne}

\begin{exercice}[Analyser un sujet]
Sujet : « L'école détruit-elle les valeurs traditionnelles africaines ? »
\begin{enumerate}
\item Dégage le ou les mots-clés du sujet et définis-les brièvement.
\item Reformule le sujet en une phrase déclarative.
\item Propose deux questions secondaires qui découlent du sujet.
\item Formule une problématique claire sous forme d'une seule question.
\end{enumerate}
\end{exercice}

\begin{exercice}[Corriger des problématiques]
Pour le sujet « Faut-il opposer tradition et modernité ? », trois élèves ont formulé les problématiques suivantes :
\begin{enumerate}
\item[a)] « La tradition et la modernité. »
\item[b)] « Pourquoi les anciens sont-ils meilleurs que les jeunes ? »
\item[c)] « Tradition et modernité sont-elles forcément incompatibles, ou peuvent-elles se concilier dans l'Afrique d'aujourd'hui ? »
\end{enumerate}
\begin{enumerate}
\item Explique pourquoi les problématiques a) et b) sont mal formulées.
\item Corrige chacune d'elles en une problématique acceptable.
\item Justifie tes corrections en t'appuyant sur les critères d'une bonne problématique (lien avec le sujet, forme interrogative, ouverture du débat).
\end{enumerate}
\end{exercice}

\begin{exercice}[Trier et regrouper des idées]
Sujet : « Le progrès technique est-il un danger pour la culture africaine ? » Voici un corpus d'idées désordonnées :
\begin{enumerate}
\item Les réseaux sociaux éloignent les jeunes des langues locales.
\item Baroka, dans \emph{Le Lion et la perle}, incarne la tradition.
\item Le téléphone portable facilite le commerce au village.
\item Les élèves doivent porter l'uniforme à l'école.
\item La photographie de Sidi publiée dans un magazine illustre la rencontre des deux mondes.
\item La télévision diffuse les danses traditionnelles à un large public.
\item Le football est le sport le plus pratiqué au Cameroun.
\item Lakunle veut moderniser Ilujinle sans en comprendre les valeurs.
\item Internet permet de préserver et de diffuser les contes oraux.
\item Les machines agricoles augmentent les récoltes.
\item Certains jeunes méprisent les rites d'initiation.
\item La ville de Douala compte de nombreux quartiers.
\end{enumerate}
\begin{enumerate}
\item Écarte les trois idées hors sujet en justifiant ton choix.
\item Regroupe les neuf idées restantes en trois ensembles cohérents.
\item Donne un titre à chaque ensemble et indique à quelle partie d'un plan dialectique il correspond.
\end{enumerate}
\end{exercice}

\begin{exercice}[Ponctuation et énonciation]
Ponctue correctement le paragraphe suivant (majuscules comprises), puis souligne trois marques explicites d'énonciation (pronoms personnels, verbes de parole, modalisateurs) :

\emph{le professeur demanda pourquoi lakunle refuse de payer la dot sidi répondit que la tradition devait être respectée n est-ce pas là le vrai problème le conflit entre tradition et modernité divise le village}

\begin{enumerate}
\item Relève dans le texte ponctué deux mots à valeur dénotative et deux mots à valeur connotative, en expliquant l'effet produit.
\item Indique la fonction de chaque signe de ponctuation que tu as ajouté.
\end{enumerate}
\end{exercice}

\courssub{Je me prépare au Bac}

\begin{exercice}[Sujet type Baccalauréat technique]
Sujet : « Le progrès technique est-il un danger pour la culture africaine ? »
\begin{enumerate}
\item Analyse le sujet : dégage les mots-clés, définis-les et repère l'enjeu du débat.
\item Formule la problématique sous forme d'une question unique et claire.
\item Recherche les idées : propose au moins trois arguments pour la thèse (le progrès menace la culture) et trois arguments pour l'antithèse (le progrès peut servir la culture), chacun accompagné d'un exemple tiré de \emph{Le Lion et la perle} ou de la réalité camerounaise.
\item Élabore le plan détaillé complet d'un plan dialectique : trois parties, chacune divisée en deux sous-parties avec arguments et exemples, sans oublier la synthèse.
\item Rédige l'annonce de plan telle qu'elle figurera à la fin de ton introduction.
\end{enumerate}
\end{exercice}

\begin{exercice}[Introduction complète]
Sujet : « Faut-il opposer tradition et modernité ? »
\begin{enumerate}
\item Rédige l'introduction complète de cette dissertation (15 à 20 lignes) en respectant les étapes : amorce, présentation du sujet, problématique, annonce du plan.
\item Ton introduction devra mobiliser au moins deux références à \emph{Le Lion et la perle} de Wole Soyinka (personnages de Baroka, Lakunle ou Sidi ; conflit des valeurs).
\item Soigne la ponctuation et utilise au moins deux connecteurs logiques que tu souligneras.
\end{enumerate}
\end{exercice}

\begin{exercice}[Évaluation croisée et compte rendu]
Voici le plan détaillé rédigé par un camarade pour le sujet : « La modernité doit-elle faire disparaître les traditions ? »

\textbf{I. Les traditions sont utiles} — 1. Elles transmettent les valeurs (exemple : les contes). — 2. Elles donnent une identité (exemple : les danses traditionnelles).

\textbf{II. La modernité est dangereuse} — 1. Elle détruit les langues locales (exemple : le français à l'école). — 2. Elle éloigne les jeunes du village (exemple : l'exode rural).

\textbf{III. Conclusion} : la modernité est mauvaise, il faut la refuser.

\begin{enumerate}
\item À l'aide de la grille critériée vue en classe (respect du type de plan, équilibre des parties, pertinence des arguments, présence des exemples, synthèse), relève au moins quatre défauts de ce plan.
\item Propose pour chaque défaut une correction précise.
\item Rédige un compte rendu d'évaluation de dix lignes adressé à ton camarade, comportant : les points forts, les points faibles et deux conseils de remédiation.
\end{enumerate}
\end{exercice}

\newpage

\coursec{Corrigés des exercices}

\coursec{Dissertation : produire le plan détaillé d'une dissertation}

\courssub{Entrée dans l'unité : paratextes et situations de communication}

\begin{aretenir}[Le paratexte : seuils de lecture]
Le \cle{paratexte} (G. Genette) désigne l'ensemble des éléments qui entourent le texte principal et médiatisent sa réception : \textbf{titre}, \textbf{nom de l'auteur}, \textbf{illustration de couverture}, \textbf{quatrième de couverture}, \textbf{préface}, \textbf{mentions éditoriales}. L'analyse du paratexte permet d'émettre des \textbf{hypothèses de lecture} (genre, ton, thèmes, époque, public visé) avant d'entrer dans l'œuvre.
\end{aretenir}

\begin{aretenir}[La situation de communication (schéma de Jakobson)]
Tout acte de communication repose sur six facteurs : \textbf{Émetteur} $\rightarrow$ \textbf{Message} $\rightarrow$ \textbf{Récepteur} ; \textbf{Contexte} (référent) ; \textbf{Code} (langue, genre littéraire) ; \textbf{Canal} (support : oral, écrit, numérique). \textbf{Fonctions} : référentielle (contexte), émotive (émetteur), conative (récepteur), poétique (message), phatique (canal), métalinguistique (code).
\end{aretenir}

\begin{corrige}[Exercice 1 — QCM : la dissertation]
\begin{enumerate}
\item Réponse \textbf{b)}. La problématique est la \cle{question centrale} qui naît de l'analyse du sujet et qui dirige toute la réflexion : tout le développement doit y répondre. L'option a) consiste à recopier le titre, ce qui ne constitue pas une problématique ; la conclusion (c) vient en fin de devoir.
\item Réponse \textbf{b)}. Le \cle{plan dialectique} comporte trois parties : la \textbf{thèse} (une première position), l'\textbf{antithèse} (la position opposée) et la \textbf{synthèse} (le dépassement du débat par une réconciliation ou une nuance).
\item Réponse \textbf{b)}. L'\cle{amorce} part d'une idée générale (un constat, une réalité sociale, une référence culturelle) pour conduire naturellement au sujet. Elle ne donne ni la réponse (a) ni le plan (c), qui sont des étapes distinctes de l'introduction.
\item Réponse \textbf{b)}. Une idée est \cle{hors sujet} lorsqu'elle ne répond pas à la question posée, même si elle est bien rédigée et développée : c'est la pertinence par rapport au sujet qui compte.
\end{enumerate}
\end{corrige}

\begin{corrige}[Exercice 2 — Vrai ou faux]
\begin{enumerate}
\item \textbf{Vrai.} L'annonce du plan est la dernière étape de l'introduction : elle vient après la problématique et présente les grandes parties du développement.
\item \textbf{Vrai.} Dans un \cle{plan détaillé}, chaque sous-partie associe un argument et au moins un exemple précis (œuvre étudiée, réalité sociale, fait d'actualité) : un plan sans exemple reste un plan sommaire.
\item \textbf{Faux.} La \cle{dénotation} est le sens premier, objectif et stable du mot (celui du dictionnaire) ; le sens imagé, affectif ou culturel relève de la \cle{connotation}.
\item \textbf{Faux.} La problématique doit être formulée \emph{avant} toute rédaction du développement : c'est elle qui oriente la recherche des idées et la construction du plan ; rédiger sans elle expose au hors sujet.
\item \textbf{Vrai.} Le \cle{point-virgule} sépare deux propositions indépendantes dont le sens est étroitement lié : \og Lakunle rêve de modernité ; Baroka défend la tradition \fg.
\end{enumerate}
\end{corrige}

\begin{corrige}[Exercice 3 — Définitions clés]
\begin{enumerate}
\item \textbf{La problématique} : question centrale, formulée à l'interrogatif, qui découle de l'analyse du sujet et résume le problème à discuter. Elle donne une direction unique à toute la dissertation : chaque partie du développement doit contribuer à y répondre.
\item \textbf{Le plan détaillé} : organisation complète du développement, rédigée au brouillon, qui présente les grandes parties (I, II, III), les sous-parties (1, 2), et pour chacune les arguments accompagnés d'exemples précis. Il sert de feuille de route pour la rédaction et garantit la cohérence et l'équilibre du devoir.
\item \textbf{La dénotation} : sens premier, objectif, stable et partagé d'un mot, tel que le donne le dictionnaire, indépendamment du contexte. \emph{Exemple} : \og lion \fg dénote un grand félin carnivore d'Afrique.
\item \textbf{La connotation} : ensemble des valeurs affectives, culturelles ou symboliques qu'un mot évoque en plus de son sens propre, selon le contexte. \emph{Exemple} : dans \emph{Le Lion et la perle}, \og lion \fg connote la force, la ruse et le pouvoir du chef Baroka ; \og perle \fg connote la beauté rare, précieuse et convoitée de Sidi.
\end{enumerate}
\end{corrige}

\begin{corrige}[Exercice 4 — Paratexte et communication (Le Lion et la perle)]
\begin{enumerate}
\item Trois éléments du \cle{paratexte} : le \textbf{titre} (\emph{Le Lion et la perle}), le \textbf{nom de l'auteur} (Wole Soyinka), l'\textbf{illustration de couverture} (souvent un jeune instituteur et une jeune fille en tenue traditionnelle). On peut ajouter la mention de l'éditeur ou la collection.
\item Ce que chaque élément permet de prévoir :
\begin{itemize}
\item Le \textbf{titre} annonce deux figures symboliques : une force (le lion) et une beauté (la perle) ; il laisse prévoir un affrontement, une conquête ou un face-à-face.
\item Le \textbf{nom de l'auteur}, dramaturge nigérian prix Nobel, fait attendre une œuvre engagée sur la société africaine, le choc des cultures et la langue théâtrale.
\item L'\textbf{illustration} annonce un cadre villageois africain (Ilujinle) et un conflit entre deux générations ou deux mondes : l'instituteur (modernité, école occidentale) face à la jeune fille en tenue traditionnelle (tradition, coutumes).
\end{itemize}
\item La situation de communication de la lecture :
\begin{center}
\begin{tabularx}{\linewidth}{|l|X|}\hline
\rowcolor{popblueL} \textbf{Élément (Jakobson)} & \textbf{Identification dans l'acte de lecture} \\ \hline
Émetteur & Wole Soyinka (l'auteur) \\ \hline
Récepteur & Le lecteur (l'élève de Première technique) \\ \hline
Message & Le récit théâtral du conflit entre tradition et modernité à Ilujinle \\ \hline
Contexte (Référent) & Le village yoruba fictif d'Ilujinle, Nigeria post-colonial \\ \hline
Code & La langue (anglais traduit en français), le code littéraire théâtral (didascalies, répliques) \\ \hline
Canal & Le livre imprimé (support écrit), support du cours \\ \hline
\end{tabularx}
\end{center}
\end{enumerate}
\end{corrige}

\courssub{Analyse du sujet et formulation de la problématique}

\begin{methode}[Analyser un sujet de dissertation — 4 étapes]
\begin{enumerate}
\item \textbf{Surligner les mots-clés} : termes forts, verbes d'action, notions abstraites. Les définir (dénotation) et questionner leur portée (connotation).
\item \textbf{Reformuler le sujet} : le transformer en phrase déclarative simple pour en vérifier la compréhension (paraphrase).
\item \textbf{Poser les questions secondaires} : explorer les oppositions, les nuances, les causes, les conséquences, les exemples possibles (thèse / antithèse).
\item \textbf{Formuler la problématique} : question unique, interrogative, qui reprend les mots-clés, ouvre un vrai débat (deux réponses légitimes) et engage la réflexion.
\end{enumerate}
\end{methode}

\begin{aretenir}[Critères d'une bonne problématique]
\begin{itemize}
\item \textbf{Lien au sujet} : reprend les termes ou synonymes du sujet.
\item \textbf{Forme interrogative} : se termine par un point d'interrogation.
\item \textbf{Ouverture du débat} : admet une thèse et une antithèse (pas de réponse évidente, pas de parti pris).
\item \textbf{Portée} : ni trop large (généralité), ni trop étroite (détail).
\end{itemize}
\end{aretenir}

\begin{corrige}[Exercice 5 — Analyser un sujet type]
\textit{Sujet : « L'école détruit-elle les valeurs traditionnelles africaines ? »}
\begin{enumerate}
\item \textbf{Mots-clés} : \og école \fg (institution d'enseignement moderne, souvent héritée de la colonisation) ; \og détruire \fg (anéantir, faire disparaître : verbe fort qui suppose un jugement) ; \og valeurs traditionnelles africaines \fg (coutumes, langues, rites, solidarités transmises par les anciens).
\item \textbf{Reformulation déclarative} : \og L'école, en imposant des savoirs et des langues étrangers, ferait disparaître les valeurs traditionnelles africaines. \fg
\item \textbf{Questions secondaires} : L'école éloigne-t-elle les jeunes des langues et des rites locaux ? L'école peut-elle au contraire transmettre et valoriser la culture traditionnelle ? L'opposition est-elle systématique ?
\item \textbf{Problématique} : \og L'école, en formant les jeunes selon des modèles étrangers, détruit-elle nécessairement les valeurs traditionnelles africaines, ou peut-elle contribuer à les préserver et à les renouveler ? \fg
\end{enumerate}
\textbf{Point de vigilance} : la problématique reprend les mots-clés du sujet, est formulée à l'interrogatif et ouvre un vrai débat (deux réponses possibles : destruction vs préservation/renouvellement).
\end{corrige}

\begin{corrige}[Exercice 6 — Corriger des problématiques défectueuses]
\textit{Sujet : « Tradition et modernité »}
\begin{enumerate}
\item \textbf{Pourquoi a) et b) sont fautives} :
\begin{itemize}
\item a) \og La tradition et la modernité \fg{} n'est qu'une \textbf{juxtaposition de mots} (nominalisation) : ce n'est pas une question, elle ne pose aucun problème et n'ouvre aucun débat.
\item b) \og Pourquoi les anciens sont-ils meilleurs que les jeunes ? \fg{} est un \textbf{hors sujet} : le sujet porte sur tradition et modernité (concepts), non sur la comparaison des générations (personnes) ; de plus, elle présuppose un jugement (\og meilleurs \fg) au lieu de l'interroger.
\end{itemize}
\item \textbf{Corrections proposées} :
\begin{itemize}
\item a) corrigée : \og Tradition et modernité sont-elles deux mondes irréconciliables ? \fg
\item b) corrigée : \og L'attachement des anciens aux traditions s'oppose-t-il nécessairement au désir de modernité des jeunes ? \fg
\end{itemize}
\item \textbf{Justification} : une bonne problématique respecte trois critères : le \textbf{lien avec le sujet} (les corrections reprennent les termes \og tradition \fg{} et \og modernité \fg), la \textbf{forme interrogative} (les deux corrections sont des questions), l'\textbf{ouverture du débat} (elles admettent une thèse et une antithèse : opposition ou conciliation). La formulation c) de l'énoncé original (\og Tradition et modernité : opposition ou conciliation ? \fg) les satisfait tous les trois, c'est pourquoi elle est la meilleure.
\end{enumerate}
\end{corrige}

\courssub{La recherche des idées}

\begin{methode}[Rechercher, trier et regrouper les idées — 3 temps]
\begin{enumerate}
\item \textbf{Remue-méninges (Brainstorming)} : noter \emph{toutes} les idées qui viennent (arguments, exemples, citations, souvenirs de cours) sans trier, pour chaque camp (thèse / antithèse).
\item \textbf{Triage (Élagage)} : éliminer les idées \textbf{hors sujet} (sans lien avec la problématique), les \textbf{doublons}, les \textbf{arguments trop faibles} ou non exemplifiables.
\item \textbf{Regroupement (Classification)} : rassembler les idées restantes en \textbf{3 ensembles cohérents} correspondant aux 3 parties du plan dialectique (Thèse / Antithèse / Synthèse). Chaque ensemble = une grande idée directrice (titre de partie). Vérifier l'équilibre (nombre d'arguments par partie).
\end{enumerate}
\end{methode}

\begin{corrige}[Exercice 7 — Trier et regrouper des idées]
\textit{Consigne : À partir d'une liste de 12 idées sur « Le progrès technique et la culture africaine », trier et structurer.}
\begin{enumerate}
\item \textbf{Idées hors sujet (3 éliminées)} :
\begin{itemize}
\item Idée 4 (\og les élèves doivent porter l'uniforme \fg) : règle scolaire sans rapport direct avec le progrès technique et la culture.
\item Idée 7 (\og le football est le sport le plus pratiqué au Cameroun \fg) : constat sportif, aucun lien avec le progrès technique.
\item Idée 12 (\og Douala compte de nombreux quartiers \fg) : simple information géographique, hors du débat.
\end{itemize}
\item \textbf{Regroupement des 9 idées restantes} :
\begin{itemize}
\item \textbf{Ensemble A} : idées 1, 8, 11 (réseaux sociaux et langues locales ; Lakunle qui méprise les valeurs d'Ilujinle ; jeunes qui méprisent les rites).
\item \textbf{Ensemble B} : idées 3, 6, 9, 10 (portable et commerce ; télévision et danses traditionnelles ; Internet et contes oraux ; machines agricoles).
\item \textbf{Ensemble C} : idées 2, 5 (Baroka incarne la tradition ; la photo de Sidi dans le magazine illustre la rencontre des deux mondes).
\end{itemize}
\item \textbf{Titres et place dans le plan dialectique} :
\begin{itemize}
\item Ensemble A : \og Le progrès technique menace les langues et les rites \fg{} $\rightarrow$ \textbf{Thèse} (Partie I).
\item Ensemble B : \og Le progrès technique peut servir et diffuser la culture \fg{} $\rightarrow$ \textbf{Antithèse} (Partie II).
\item Ensemble C : \og Tradition et modernité peuvent se rencontrer et composer \fg{} $\rightarrow$ \textbf{Synthèse} (Partie III).
\end{itemize}
\end{enumerate}
\textbf{Point de vigilance} : un regroupement est valide si chaque ensemble repose sur un critère unique (ici : menace / service / conciliation) et si chaque idée est à sa place ; vérifier toujours que le total des idées triées correspond au total de départ ($3 + 9 = 12$).
\end{corrige}

\courssub{L'élaboration du plan détaillé}

\begin{aretenir}[Structure type du plan dialectique (3 parties, 2 sous-parties chacune)]
\begin{center}
\begin{tabularx}{\linewidth}{|l|X|X|}\hline
\rowcolor{popblueL} \textbf{Partie} & \textbf{Rôle} & \textbf{Contenu type} \\ \hline
\textbf{I. Thèse} & Défendre une première réponse & 2 arguments + exemples (ex: le progrès menace) \\ \hline
\textbf{II. Antithèse} & Défendre la réponse opposée & 2 arguments + exemples (ex: le progrès sert la culture) \\ \hline
\textbf{III. Synthèse} & Dépasser l'opposition & 2 arguments + exemples (ex: conditions de conciliation, modernité enracinée) \\ \hline
\end{tabularx}
\end{center}
\textbf{Règle d'or} : Chaque sous-partie = \textbf{1 argument} + \textbf{1 exemple précis} (œuvre, fait social, chiffre, citation). La synthèse ne répète pas la thèse, elle \textbf{dépasse} le débat.
\end{aretenir}

\begin{methode}[Rédiger l'annonce de plan (dernière phrase de l'introduction)]
Utiliser des connecteurs logiques séquentiels : \og Dans un premier temps, nous montrerons que... \fg ; \og Dans un second temps, nous verrons que... \fg ; \og Enfin, nous dégagerons... \fg. L'annonce doit refléter \textbf{exactement} les titres des trois parties du plan détaillé.
\end{methode}

\begin{corrige}[Exercice 9 — Sujet type Baccalauréat technique : Plan détaillé complet]
\textit{Sujet : « Le progrès technique est-il un danger pour la culture africaine ? »}
\begin{enumerate}
\item \textbf{Analyse du sujet} : mots-clés : \og progrès technique \fg{} (inventions et outils modernes : téléphone, Internet, machines), \og danger \fg{} (ce qui menace, détruit), \og culture africaine \fg{} (langues, rites, arts, valeurs transmises). Enjeu : faut-il craindre que la technique fasse disparaître la culture, ou peut-elle la servir ?
\item \textbf{Problématique} : \og Le progrès technique menace-t-il nécessairement la culture africaine, ou peut-il devenir un instrument de sa sauvegarde et de son rayonnement ? \fg
\item \textbf{Recherche des idées} (sélectionnées et exemplifiées) :

\textbf{Thèse (le progrès menace la culture)} :
\begin{itemize}
\item Les réseaux sociaux détournent les jeunes des langues locales (exemple : des élèves camerounais qui n'écrivent plus en ewondo ou en duala).
\item Lakunle, l'instituteur, méprise la dot et les coutumes d'Ilujinle au nom d'une modernité mal comprise.
\item Certains jeunes citadins renient les rites d'initiation (exemple : l'abandon des cérémonies villageoises).
\end{itemize}

\textbf{Antithèse (le progrès peut servir la culture)} :
\begin{itemize}
\item Le téléphone portable facilite le commerce et maintient le lien avec le village (exemple : les transferts d'argent Mobile Money au Cameroun).
\item La télévision et Internet diffusent les danses et les contes traditionnels (exemple : les émissions culturelles de la CRTV).
\item Le magazine qui publie la photo de Sidi fait connaître la beauté d'Ilujinle au monde entier.
\end{itemize}
\item \textbf{Plan détaillé (dialectique)} :

\textbf{I. Le progrès technique apparaît comme un danger pour la culture africaine.}
\begin{itemize}
\item 1. Il éloigne les jeunes des langues et des rites (argument : uniformisation des pratiques ; exemple : réseaux sociaux, abandon des rites d'initiation).
\item 2. Il peut engendrer le mépris des valeurs traditionnelles (exemple : Lakunle qui refuse la dot sans comprendre Ilujinle).
\end{itemize}

\textbf{II. Pourtant, le progrès technique peut servir la culture.}
\begin{itemize}
\item 1. Il facilite la vie et renforce les solidarités (exemples : portable et commerce, Mobile Money, machines agricoles).
\item 2. Il diffuse et préserve le patrimoine (exemples : CRTV et danses traditionnelles, contes oraux numérisés, photo de Sidi dans le magazine).
\end{itemize}

\textbf{III. Tradition et modernité peuvent se concilier.}
\begin{itemize}
\item 1. La technique n'est qu'un outil : tout dépend de l'usage qu'on en fait (exemple : Baroka lui-même sait user du moderne pour défendre son village).
\item 2. Une modernité enracinée dans la culture est possible (exemple : les jeunes créateurs camerounais qui mêlent makossa et musiques urbaines).
\end{itemize}
\item \textbf{Annonce de plan} : \og Après avoir montré que le progrès technique peut menacer la culture africaine, nous verrons qu'il peut aussi la servir, avant de dégager les conditions d'une conciliation entre tradition et modernité. \fg
\end{enumerate}
\textbf{Barème indicatif (sur 20)} : analyse du sujet (3 pts), problématique (3 pts), recherche des idées (6 pts : 6 arguments exemplifiés), plan détaillé équilibré (6 pts), annonce de plan (2 pts).

\textbf{Points de vigilance} : problématique interrogative et liée au sujet ; trois parties équilibrées ; chaque sous-partie = argument + exemple précis ; la synthèse ne répète pas la thèse mais dépasse l'opposition.
\end{corrige}

\begin{corrige}[Exercice 10 — Introduction complète rédigée]
\textbf{Proposition d'introduction} (connecteurs logiques soulignés) :

Dans toute société, le passage des générations s'accompagne d'une tension entre ce qui est hérité et ce qui est inventé. L'Afrique contemporaine, marquée par l'école, la ville et les techniques nouvelles, vit cette tension avec une acuité particulière. C'est précisément le sujet que met en scène Wole Soyinka dans \emph{Le Lion et la perle} : à Ilujinle, l'instituteur Lakunle rêve de transformer le village et refuse de payer la dot pour épouser Sidi, tandis que le vieux chef Baroka défend avec ruse les coutumes de ses ancêtres. \underline{Dès lors}, une question s'impose : faut-il opposer tradition et modernité ? Autrement dit, ces deux forces sont-elles nécessairement ennemies, ou peuvent-elles se concilier pour construire l'Afrique de demain ? Pour répondre à cette problématique, nous montrerons \underline{dans un premier temps} que tradition et modernité semblent s'opposer radicalement ; nous verrons \underline{ensuite} que chacune possède ses richesses et ses limites ; \underline{enfin}, nous dégagerons les conditions d'un équilibre entre les deux.

\textbf{Points de vigilance et barème indicatif (sur 10)} : amorce générale (2 pts), présentation du sujet avec deux références précises à l'œuvre — Lakunle, la dot, Baroka, Sidi (3 pts), problématique interrogative (2 pts), annonce de plan en trois temps (2 pts), ponctuation soignée et connecteurs soulignés (1 pt). \textbf{Erreurs fréquentes} : commencer par une définition du dictionnaire, annoncer le plan avant la problématique, donner déjà son avis personnel (\og Je pense que... \fg).
\end{corrige}

\courssub{Outils linguistiques de la dissertation}

\begin{aretenir}[Ponctuation majeure pour la dissertation]
\begin{itemize}
\item \textbf{Deux-points (:)} : annoncent une explication, une énumération ou une citation (discours direct).
\item \textbf{Point-virgule (;)} : sépare deux propositions indépendantes liées par le sens (cause, conséquence, opposition, parallélisme).
\item \textbf{Guillemets (« »)} : encadrent les citations, les titres d'œuvres, les mots utilisés dans un sens particulier (ironie, métonymie).
\item \textbf{Parenthèses ( )} : insèrent une précision, une référence, une explication accessoire.
\item \textbf{Points de suspension (...)} : marquent l'hésitation, la réticence ou une énumération non close.
\end{itemize}
\end{aretenir}

\begin{aretenir}[Marques explicites de l'énonciation (repérage du locuteur)]
\begin{itemize}
\item \textbf{Personnes} : \og je \fg, \og nous \fg (énonciateur) ; \og tu \fg, \og vous \fg (allocutaire).
\item \textbf{Verbes de parole / pensée} : \og dire \fg, \og affirmer \fg, \og penser \fg, \og soutenir \fg (souvent au présent de narration ou passé simple).
\item \textbf{Adverbes / Modalisateurs} : \og certainement \fg, \og sans doute \fg, \og hélas \fg, \og heureusement \fg, \og évidemment \fg.
\item \textbf{Interrogations rhétoriques} : \og N'est-ce pas ? \fg, \og Comment admettre que... ? \fg (marquent la prise de position).
\item \textbf{Évaluation / Jugement} : adjectifs mélioratifs / péjoratifs (\og admirable \fg / \og détestable \fg).
\end{itemize}
\end{aretenir}

\begin{aretenir}[Dénotation vs Connotation : outil d'analyse et d'argumentation]
\begin{center}
\begin{tabularx}{\linewidth}{|l|X|X|}\hline
\rowcolor{popblueL} & \textbf{Dénotation} & \textbf{Connotation} \\ \hline
\textbf{Définition} & Sens propre, objectif, dictionnaire & Sens figuré, subjectif, culturel, affectif \\ \hline
\textbf{Stabilité} & Stable dans le temps & Variable selon contexte, époque, locuteur \\ \hline
\textbf{Exemple} & \og Rose \fg = fleur épineuse & \og Rose \fg = amour, beauté, éphémère, passion \\ \hline
\textbf{Usage dissertation} & Définir les termes du sujet (analyse) & Argumenter, persuader, nuancer (développement) \\ \hline
\end{tabularx}
\end{center}
\end{aretenir}

\begin{corrige}[Exercice 8 — Ponctuation, énonciation, dénotation/connotation]
\textbf{Texte à ponctuer} :
\og Le professeur demanda : \og Pourquoi Lakunle refuse-t-il de payer la dot ? \fg{} Sidi répondit que la tradition devait être respectée. N'est-ce pas là le vrai problème ? Le conflit entre tradition et modernité divise le village. \fg

\textbf{Trois marques explicites d'énonciation} (identifiées dans le texte) :
\begin{itemize}
\item Les verbes de parole : \og \underline{demanda} \fg{} et \og \underline{répondit} \fg (marquent l'énonciation rapportée).
\item Le modalisateur interrogatif : \og \underline{N'est-ce pas} \fg{} (marque la prise de position de l'énonciateur, invite l'adhésion du lecteur).
\end{itemize}

\begin{enumerate}
\item \textbf{Dénotation / Connotation} :
\begin{itemize}
\item \textbf{Mots dénotatifs} : \og professeur \fg{} (celui qui enseigne : sens objectif, fonction) ; \og dot \fg{} (biens versés par le prétendant à la famille de la fiancée : sens précis, culturellement défini, juridico-coutumier).
\item \textbf{Mots connotatifs} : \og tradition \fg{} (évoque le respect des anciens, la sagesse, la cohésion, mais aussi parfois l'immobilisme, le frein au changement) ; \og conflit \fg{} (évoque la violence, la division, la souffrance du village déchiré, l'affrontement irréductible). \textbf{Effet produit} : ces connotations engagent le lecteur affectivement et donnent au débat une dimension dramatique et existentielle.
\end{itemize}
\item \textbf{Fonction des signes de ponctuation ajoutés} :
\begin{itemize}
\item Les \textbf{deux-points} introduisent la parole rapportée au discours direct.
\item Les \textbf{guillemets} encadrent les propos cités (discours direct imbriqué).
\item Le \textbf{point d'interrogation} marque les phrases interrogatives (questions directes et rhétoriques qui font avancer la réflexion).
\item Le \textbf{point} marque la fin des phrases déclaratives (affirmations narratives).
\item Les \textbf{majuscules} signalent le début de chaque phrase et les noms propres (Lakunle, Sidi).
\end{itemize}
\end{enumerate}
\end{corrige}

\courssub{Compte rendu, remédiation et projet d'études}

\begin{methode}[Rédiger un compte rendu d'évaluation croisée (pair-à-pair)]
\begin{enumerate}
\item \textbf{Lire attentivement} la production du camarade (plan + introduction) sans juger a priori.
\item \textbf{Identifier les points forts} (structure respectée, exemples pertinents, clarté, langue soignée).
\item \textbf{Repérer les faiblesses} (hors sujet, déséquilibre, plan non dialectique, exemples manquants, formulation floue).
\item \textbf{Formuler des corrections précises} pour chaque faiblesse (pas seulement « c'est faux », mais « voici comment corriger »).
\item \textbf{Rédiger le compte rendu} : ton objectif, bienveillant, structuré (Points forts / Points faibles / Conseils concrets).
\end{enumerate}
\end{methode}

\begin{corrige}[Exercice 11 — Évaluation croisée et compte rendu]
\textit{Consigne : Évaluer le plan détaillé d'un camarade sur le sujet « Le progrès technique est-il un danger pour la culture africaine ? »}
\begin{enumerate}
\item \textbf{Au moins quatre défauts relevés dans la production de l'élève} :
\begin{itemize}
\item \textbf{Type de plan non respecté} : le sujet appelle un plan dialectique (thèse / antithèse / synthèse) ; or la partie III n'est pas une synthèse mais une conclusion personnelle (\og La modernité est mauvaise \fg).
\item \textbf{Déséquilibre et parti pris} : les deux premières parties défendent toutes les deux la tradition et accusent la modernité ; le point de vue opposé (les bienfaits de la modernité : santé, éducation, communication) n'est jamais examiné.
\item \textbf{Arguments discutables / réducteurs} : présenter le français à l'école comme une \og destruction \fg{} est un raccourci non nuancé (bilinguisme possible) ; l'exode rural n'est pas causé par la seule modernité.
\item \textbf{Pas de vraie synthèse} : la partie III conclut brutalement sans dépasser l'opposition ni répondre vraiment à la problématique (conditions de conciliation).
\item \textbf{Exemples insuffisamment exploités} : ils sont cités en un mot (ex: \og téléphone \fg), sans lien explicite avec l'argument ni référence à l'œuvre étudiée (\emph{Le Lion et la perle}).
\end{itemize}
\item \textbf{Corrections proposées (remédiation)} :
\begin{itemize}
\item Transformer la partie I en \textbf{Thèse} : \og Le progrès technique menace les repères culturels \fg (uniformisation, perte des langues, mépris des rites — ex: Lakunle).
\item Transformer la partie II en \textbf{Antithèse} : \og Le progrès technique offre des opportunités de préservation et de diffusion \fg (Mobile Money, CRTV, numérisation des contes, agriculture).
\item Rédiger une vraie \textbf{Synthèse (Partie III)} : \og Il ne faut ni rejeter ni subir la modernité, mais l'enraciner dans les valeurs traditionnelles \fg (ex: Baroka utilisant le magazine ; jeunes artistes mêlant tradition et modernité).
\item Nuancer les arguments : le français peut coexister avec les langues locales (bilinguisme officiel au Cameroun) ; l'exode rural a des causes économiques multiples.
\item Développer chaque exemple en 2-3 phrases : présenter l'exemple, l'analyser, le relier à l'argument.
\end{itemize}
\item \textbf{Compte rendu d'évaluation (modèle type)} :

\og Ton plan présente des points forts réels : il est lisible, chaque partie comporte deux sous-parties et chaque sous-partie contient un exemple, ce qui montre que tu as compris la structure attendue du plan détaillé. Cependant, il souffre de trois faiblesses majeures :
\begin{enumerate}
\item Il ne respecte pas le plan dialectique : ta troisième partie est une conclusion définitive (\og la modernité est mauvaise \fg) et non une synthèse qui dépasse le débat.
\item Il est déséquilibré : tu n'examines jamais les arguments contraires (les avantages de la modernité), ce qui donne l'impression d'un parti pris.
\item Ton jugement final est trop catégorique et ne répond pas à la problématique qui demande une nuance (\og nécessairement \fg / \og ou \fg).
\end{enumerate}
Pour progresser, je te conseille d'abord de toujours vérifier, avant de rédiger, que ton plan comporte une thèse, une antithèse et une synthèse qui propose une issue. Ensuite, entraîne-toi à formuler pour chaque sujet les arguments du point de vue opposé au tien : c'est la meilleure garantie d'un devoir équilibré et réussi. \fg
\end{enumerate}
\textbf{Points de vigilance} : le compte rendu doit être objectif, bienveillant et structuré (points forts, points faibles, conseils) ; chaque défaut relevé doit être accompagné d'une correction précise, pas d'une simple critique.
\end{corrige}

\newpage

\coursec{Fiche bilan}

\begin{popfigure}
\begin{tikzpicture}[
    mindmap,
    grow cyclic,
    every node/.style={concept, execute at begin node=\hskip0pt},
    concept color=popblue,
    level 1/.append style={level distance=4.5cm, sibling angle=60, font=\small\bfseries},
    level 2/.append style={level distance=3cm, sibling angle=45, font=\footnotesize},
    branch1/.style={concept color=poporange, parent anchor=children, child anchor=parent},
    branch2/.style={concept color=popgreen, parent anchor=children, child anchor=parent},
    branch3/.style={concept color=poppurple, parent anchor=children, child anchor=parent},
    branch4/.style={concept color=poppink, parent anchor=children, child anchor=parent},
    branch5/.style={concept color=popgold, parent anchor=children, child anchor=parent},
    branch6/.style={concept color=popblue, parent anchor=children, child anchor=parent}
]
\node[concept, text width=3.5cm, align=center, font=\normalsize\bfseries] {Dissertation\\Plan détaillé}
    child[branch1] { node[concept, text width=2.8cm, align=center] {Analyse du sujet}
        child { node[concept, text width=2.2cm, align=center] {Mots-clés \& Définitions} }
        child { node[concept, text width=2.2cm, align=center] {Formulation problèmatique} }
        child { node[concept, text width=2.2cm, align=center] {Type de sujet (ouvert/fermé)} }
    }
    child[branch2] { node[concept, text width=2.8cm, align=center] {Recherche d'idées}
        child { node[concept, text width=2.2cm, align=center] {Brainstorming / Mindmap} }
        child { node[concept, text width=2.2cm, align=center] {Arguments \& Exemples} }
        child { node[concept, text width=2.2cm, align=center] {Regroupement thématique} }
    }
    child[branch3] { node[concept, text width=2.8cm, align=center] {Structure du plan}
        child { node[concept, text width=2.2cm, align=center] {Introduction (AMI)} }
        child { node[concept, text width=2.2cm, align=center] {Développement (2-3 parties)} }
        child { node[concept, text width=2.2cm, align=center] {Conclusion (Bilan/Ouverture)} }
    }
    child[branch4] { node[concept, text width=2.8cm, align=center] {Types de plans}
        child { node[concept, text width=2.2cm, align=center] {Dialectique (Thèse/Antithèse/Synthèse)} }
        child { node[concept, text width=2.2cm, align=center] {Thématique (Aspects variés)} }
        child { node[concept, text width=2.2cm, align=center] {Analytique (Causes/Conséquences)} }
    }
    child[branch5] { node[concept, text width=2.8cm, align=center] {Outils linguistiques}
        child { node[concept, text width=2.2cm, align=center] {Ponctuation \& Connecteurs} }
        child { node[concept, text width=2.2cm, align=center] {Énonciation (Marqueurs)} }
        child { node[concept, text width=2.2cm, align=center] {Dénotation / Connotation} }
        child { node[concept, text width=2.2cm, align=center] {Texte descriptif} }
    }
    child[branch6] { node[concept, text width=2.8cm, align=center] {Écriture \& Révision}
        child { node[concept, text width=2.2cm, align=center] {Rédaction soignée} }
        child { node[concept, text width=2.2cm, align=center] {Relecture (Code/Com)} }
        child { node[concept, text width=2.2cm, align=center] {Remédiation / Projet} }
    };
\end{tikzpicture}
\legende{Carte mentale : Les 6 étapes clés pour produire un plan détaillé de dissertation en Première Technique.}
\end{popfigure}

\coursec{Bilan et Synthèse de la leçon}

\courssub{Définitions clés}

\begin{definition}[Sujet de dissertation]
Énoncé (citation, phrase affirmative, question) invitant à une réflexion argumentée. Il comporte des \cle{mots-clés} à définir et une \cle{tension} (problématique) à expliciter.
\end{definition}

\begin{definition}[Problématique]
Question centrale, ouverte et débattue, issue de l'analyse du sujet. Elle guide toute la démonstration. Elle ne se confond pas avec le sujet.
\end{definition}

\begin{definition}[Plan détaillé]
Squelette complet de la dissertation : Introduction (Accroche, Mise au point, Intérêt, Problématique, Annonce du plan), Développement (Parties I, II, III ; Sous-parties avec arguments précis, exemples, citations), Conclusion (Bilan synthétique, Réponse à la problématique, Ouverture).
\end{definition}

\begin{definition}[Types de plans]
\begin{itemize}
    \item \textbf{Dialectique :} Thèse / Antithèse / Synthèse (idéal pour sujet polémique).
    \item \textbf{Thématique :} Aspect 1 / Aspect 2 / Aspect 3 (idéal pour sujet ouvert « Quels sont les... ? »).
    \item \textbf{Analytique :} Causes / Conséquences / Solutions (idéal pour sujet « Pourquoi / Comment »).
\end{itemize}
\end{definition}

\courssub{Méthodes express (à mémoriser)}

\begin{methode}[Analyser un sujet (4 étapes)]
\begin{enumerate}
    \item \textbf{Cercler} les mots-clés ; les \textbf{définir} (sens littéral / figuré, dénotation / connotation).
    \item \textbf{Identifier} le type de sujet (question, affirmation, citation) et la discipline concernée.
    \item \textbf{Formuler} 2 à 3 problèmes possibles ; \textbf{choisir} la problématique la plus féconde (ni trop large, ni trop étroite).
    \item \textbf{Valider} : la problématique permet-elle de structurer un plan en 2 ou 3 parties équilibrées ?
\end{enumerate}
\end{methode}

\begin{methode}[Rechercher et organiser les idées (3 temps)]
\begin{enumerate}
    \item \textbf{Exhaustivité :} Noter \emph{toutes} idées, arguments, exemples, références (lecture ouvroir, culture personnelle) sans trier.
    \item \textbf{Regroupement :} Classer les idées en 2 ou 3 \cle{grands axes} (futures parties). Vérifier l'équilibre (même nombre d'arguments par partie).
    \item \textbf{Hiérarchisation :} Dans chaque partie, ordonner les sous-parties (logique : cause $\to$ conséquence, simple $\to$ complexe, général $\to$ particulier).
\end{enumerate}
\end{methode}

\begin{methode}[Rédiger le plan détaillé (Standard AMI + Développement)]
\begin{enumerate}
    \item \textbf{Introduction :} \cle{Accroche} (citation, fait d'actualité, définition) $\to$ \cle{Mise au point} (définitions, délimitation) $\to$ \cle{Intérêt} (enjeux) $\to$ \cle{Problématique} $\to$ \cle{Annonce du plan} (connecteurs : « Nous verrons d'abord... Ensuite... Enfin... »).
    \item \textbf{Développement :} Pour chaque partie (I, II, III) : \textbf{Titre de partie} (phrase nominale) $\to$ Sous-parties (A, B, C) avec \textbf{argument principal}, \textbf{exemple précis} (texte étudié, œuvre, fait social), \textbf{connecteur logique}.
    \item \textbf{Conclusion :} \cle{Bilan} (reprise synthétique des grandes parties) $\to$ \cle{Réponse directe} à la problématique $\to$ \cle{Ouverture} (autre œuvre, actualité, question nouvelle).
\end{enumerate}
\end{methode}

\courssub{Outils linguistiques indispensables (Tableau de révision)}

\begin{center}
\begin{tabularx}{\linewidth}{|l|X|X|}
\hline
\rowcolor{popblueL}
\textbf{Notion} & \textbf{Rappel essentiel} & \textbf{Effet en dissertation} \\ \hline
\cle{Ponctuation} & Point (fin assertion), deux-points (annonce/explication), point-virgule (lien fort), parenthèses (incise). & Structurer la pensée, marquer les relations logiques, aérer la lecture. \\ \hline
\cle{Connecteurs logiques} & \textbf{Addition :} de plus, en outre. \textbf{Opposition :} cependant, pourtant. \textbf{Conséquence :} par conséquent, dès lors. \textbf{Explication :} en effet, c'est-à-dire. & Assurer la \cle{cohérence} inter/intra-phrastique ; guider le correcteur. \\ \hline
\cle{Énonciation} & Marques \textbf{explicites} : je, nous, mon avis, selon moi. Marques \textbf{implicites} : adverbes (heureusement), modalisateurs (sans doute), passif. & Affirmer la \cle{subjectivité} maîtrisée (je argumentatif) ou l'\cle{objectivité} (impersonnel). \\ \hline
\cle{Dénotation / Connotation} & \textbf{Dénotation :} sens littéral, dictionnaire. \textbf{Connotation :} charge affective, culturelle, subjective (ex: « domicile » vs « taudis »). & Préciser les définitions (dénotation) ; enrichir l'argumentation par les valeurs (connotation). \\ \hline
\cle{Texte descriptif} & Verbes d'état, adjectifs qualitatifs, spatialisation (G$\to$D, H$\to$B), figures (comparaison, métaphore, personnification). & Illustrer un argument par l'exemple descriptif (portrait, lieu, atmosphère) pour convaincre par l'image. \\ \hline
\end{tabularx}
\end{center}

\courssub{Spécificités « Le Lion et la Perle » (Texte ouvroir)}
\begin{itemize}
    \item \textbf{Paratextes :} Titre (opposition/contraste), Auteur (Wole Soyinka), Genre (Comédie/Théâtre), Contexte (Nigéria post-colonial, tradition vs modernité).
    \item \textbf{Thèmes exploitables :} Conflit des générations, Place de la femme, Pouvoir du langage/baratin, Modernité vs Tradition, Identité africaine.
    \item \textbf{Références obligatoires :} Personnages (Baroka, Lakunlé, Sidi, Sadiku), Scènes clés (la danse, le timbre-poste, la séduction finale), Registres (comique, satirique, pathétique).
\end{itemize}

\begin{aretenir}[Checklist avant le Bac (Probatoire)]
\begin{itemize}
    \item $\square$ Je sais \textbf{distinguer} sujet, thèse et problématique.
    \item $\square$ Je sais \textbf{formuler} une problématique ouverte, débattue et opératoire.
    \item $\square$ Je maîtrise la \textbf{méthode AMI} pour l'introduction (Accroche, Mise au point, Intérêt, Problématique, Annonce).
    \item $\square$ Je sais \textbf{choisir le type de plan} adapté (Dialectique / Thématique / Analytique) et le justifier.
    \item $\square$ Je construis un \textbf{plan détaillé} : titres de parties (phrases nominales) + sous-parties (arguments + exemples précis).
    \item $\square$ J'utilise des \textbf{connecteurs logiques} variés et adaptés à la relation d'idées (opposition, concession, conséquence...).
    \item $\square$ Je mobilise le \textbf{texte ouvroir} (\emph{Le Lion et la Perle}) comme référence principale (citations, scènes, personnages).
    \item $\square$ Je soigne la \textbf{ponctuation} et la syntaxe pour marquer l'articulation de ma pensée.
    \item $\square$ Je gère l'\textbf{énonciation} : « je » argumentatif assumé ou impersonnalisation selon la convention attendue.
    \item $\square$ Je rédige une \textbf{conclusion en 3 temps} : Bilan synthétique $\to$ Réponse nette $\to$ Ouverture pertinente.
    \item $\square$ Je relis pour \textbf{corriger} (orthographe, grammaire, cohérence) et \textbf{vérifier} l'adéquation plan/rédaction.
    \item $\square$ Je respecte le \textbf{format d'examen} : brouillon structuré, copie finale lisible, gestion du temps (3h-4h).
\end{itemize}
\end{aretenir}

\findecours

\end{document}
